% !TeX root = ../Book3.tex
% 纲要标题：# 第四编　交换代数进阶与几何初步
% 纲要位置：计划纲要.md，第 2184 行。

\BookPart{交换代数进阶与几何初步}{b3:part:04}
\BookPartGuide
本编把前三编的局部方法用于三个问题。局部上同调计算由给定理想确定的支集上的模信息，Matlis 对偶\footnote{马特利斯（Eben Matlis，1923-2015），美国数学家，研究交换环上的内射模与对偶理论。}与局部对偶再把这些计算同深度和自由消解联系起来；在完备局部环中，Hensel 引理追问模极大理想的根或因子能否提升；把局部化中的函数和模截面按开集组织并黏合起来，则得到结构层、概形和拟凝聚层。

\ProseParagraphBreak
\chapref{b3:ch:22}从指定理想的挠子模出发，以 Čech 复形\footnote{切赫（Eduard Čech，1893-1960），捷克数学家，研究拓扑与微分几何；以开覆盖构造的 Čech 上同调以他命名。}计算局部上同调，再证明深度的首次非消失刻画。在完备 Noether 局部环上，内射包络与 Matlis 对偶把有限生成模和 Artin 模互相转换，双对偶返回原模。给定完备 Noether 正则局部环的满射表示 \(S\twoheadrightarrow R\)，本章通过有限自由消解建立局部对偶，并证明一般 Bass 判别\footnote{巴斯（Hyman Bass，1932-），美国数学家，研究代数 K 理论、环和模。}；一般完备 Noether 局部环的相应满射表示由附录D的 Cohen 结构定理提供。这里需要第二卷的 Baer 判据、足够多内射模，以及第13章的完备化和第15—19章的 Koszul、深度、消解与 Cohen–Macaulay 理论。只希望先认识局部上同调的读者可以先读\secref{b3:sec:2201}，把内射包络和对偶留作进一步学习。

\ProseParagraphBreak
\chapref{b3:ch:23}接续第14章的微分和平方零提升。在完备局部环中，单根与互素因式的 Hensel 提升在本章证明；Hensel 化的存在则以证明概要采用永田的构造与有限单根邻域描述，再展开局部性、忠实平坦性及 Noether 情形的有限阶商比较。对于每个平方零提升问题，形式光滑要求提升存在，形式非分歧要求提升至多唯一，形式平展要求存在且唯一。局部标准形与 Jacobian 计算把这些条件落实到具体方程；Noether 局部形式光滑蕴含平坦的结论采用文献中的局部判据，并给出核对拓扑假设的概要，在此基础上证明任意底环上光滑同态的平坦性。读这一章主要需要第12—14章；比较正则性与光滑性时还会用到第18章。尖点及 Frobenius 的例子\footnote{弗罗贝尼乌斯（Ferdinand Georg Frobenius，1849-1917），德国数学家，研究群表示、矩阵和有限域。}用来区分绝对的环性质与相对于基底的性质。

\ProseParagraphBreak
\chapref{b3:ch:24}先把素谱上的分式逐点黏合成结构层，再建立仿射概形和拟凝聚层的代数描述，由分次环构造 Proj。前三节的主要基础是第1—3章的素谱、模局部化与有限性、第11章的分次环和第12章的局部自由判据。第24.4节通过纤维积、对角和齐次环证明射影空间及其闭子概形的固有性。第24.5节引入层上同调：Noether 仿射概形上的拟凝聚层消失给出证明概要，所需的一般截面延拓仍采用层论结果；在这些输入下，对所有非空有限交仍仿射的有限仿射开覆盖，展开 Čech 比较的双复形证明，并计算射影直线扭曲层的上同调。这一节还需要内射消解与右导出函子的基本知识。

\ProseParagraphBreak
对偶与 Gorenstein 环可参考松村英之《Commutative Ring Theory》\cite[\S 18]{Matsumura1989CommutativeRingTheory}，光滑性与平展性可参考同书\cite[\S 28]{Matsumura1989CommutativeRingTheory}。概形部分对应 Hartshorne《Algebraic Geometry》\cite[Chapter II, \S\S 1--5; Chapter III, \S\S 2--4]{Hartshorne1977}；各章导读与具体结果处再给出所用的定理、页码及适用范围。

\ProseParagraphBreak
读过本编，应能把环与模组织成概形和拟凝聚层，辨认局部化所得的茎、基变换所得的纤维以及黏合所需的相容条件。对具体方程，可以检查根与同态的提升条件，用局部上同调和对偶描述闭支集附近的信息，并在满足交集条件的覆盖上计算层上同调，判断局部截面为什么能够或不能成为一个全局截面。

% !TeX root = ../../Book3.tex
% 纲要标题：## 第22章*　局部上同调、典范模与 Gorenstein 环
% 纲要位置：计划纲要.md，第 2667 行。

\OptionalChapter{局部上同调、典范模与 Gorenstein 环}{b3:ch:22}
\BookChapterNumber{b3:ch:22}

\begin{chapterabstract}
本章由 Čech 复形计算局部上同调，并用首次非零次数刻画深度。完备局部环上的 Matlis 对偶，将 Artin 模与有限生成模相互联系。把局部环写成完备正则局部环的商后，环境环上的 Ext 给出典范模；Cohen–Macaulay 条件使对偶信息集中在一个次数，Gorenstein 自对偶\footnote{戈伦斯坦（Daniel Gorenstein，1923-1992），美国数学家，研究有限单群分类；其早年关于平面曲线的工作引出了以他命名的环。}则进一步表现为典范模与环自身同构。非 Cohen–Macaulay 情形须保留多个对偶次数。
\end{chapterabstract}

\begin{learninggoals}
\begin{itemize}
\item 用有限 Čech 复形计算局部上同调，并解释深度和维数给出的两端消失区间。
\item 检验内射包络的本质扩张条件，写出有限生成模与 Artin 模的 Matlis 求值同构，并指出完备性的作用。
\item 用来自完备正则局部环的满射表示构造典范模，解释 CM 条件如何使局部对偶集中在一个次数。
\item 由 Koszul 对偶计算完全交的典范模，用基座或具体乘法配对检验零维 Gorenstein 性。
\end{itemize}
\end{learninggoals}

设 \(k\) 为域，令 \(R=k[x,y]\)，把 \(R_x,R_y\) 放在同一个 Laurent 多项式环 \(R_{xy}\) 中。\(1/(xy)\) 不能写成这两个子环中元素之差：\(R_x\) 的单项式有非负的 \(y\) 指数，\(R_y\) 的单项式有非负的 \(x\) 指数，故两者之和中 \(x^{-1}y^{-1}\) 的系数总为零。这就给出 \(R_{xy}/(R_x+R_y)\) 的一个非零类，显示交叠处的分式未必来自两边局部化数据的差。

\ProseParagraphBreak

Čech 复形把这种“局部有解、全局不能拼合”的差异延伸到多个开集，并产生局部上同调。它的首次非零次数与前面用正则序列定义的深度相遇；而在完备局部情形，与内射模中的对偶计算相配，又可把上同调转成 Ext 的信息。这些 Ext 模是有限生成的，因而原本含有无界分母的局部上同调，可以借助有限生成模的消解来研究。

\ProseParagraphBreak

本章需要\chapref{b3:ch:13}的 Artin–Rees 与完备化、\chapref{b3:ch:15}的 Koszul 复形、\chapref{b3:ch:16}的深度，以及\chapref{b3:ch:18}和\chapref{b3:ch:19}的正则局部环与 CM 环。第二卷第7.2节的 Baer 判据、足够多内射模定理也将直接使用。内射消解、比较映射及单行双复形计算在使用处给出构造；一般导出函子和导出范畴的系统论述安排在第四卷。

\ProseParagraphBreak
设 \(R\) 是来自完备正则局部环 \(S\) 的满射像。典范模可以从 \(\operatorname{Ext}_S(R,S)\) 的一个次数取出；若 \(R\) 为 Cohen–Macaulay 环，其余次数消失，便得到只涉及这个模的局部对偶。完全交的 Koszul 消解具有自对偶性，因此典范模同构于 \(R\)；不满足 CM 条件的商环则可能保留多个非零次数。完备局部环具有这类满射表示，是 Cohen 结构定理的结论，见\appref{b3:app:D}。

\ProseParagraphBreak
局部上同调的导出函子定义、Koszul 计算及深度判据，可参阅 Weibel《An Introduction to Homological Algebra》\cite[\S 4.6, pp. 115--119]{Weibel1994HomologicalAlgebra}。一般 Gorenstein 环的判别见松村英之《Commutative Ring Theory》\cite[\S 18, Theorem 18.1, pp. 141--144]{Matsumura1989CommutativeRingTheory}；内射包络和 Matlis 对偶见同书\cite[\S 18, Theorems 18.4--18.6, pp. 145--149]{Matsumura1989CommutativeRingTheory}。典范模与局部对偶的进一步论述，见 Bruns 与 Herzog《Cohen--Macaulay Rings》修订版\cite[\S\S 3.3, 3.5]{BrunsHerzog1998CohenMacaulay}。

% !TeX root = ../../Book3.tex
% 纲要标题：### 22.1 局部上同调
% 纲要位置：计划纲要.md，第 2669 行。

\section{局部上同调}
\label{b3:sec:2201}

模中的元素可能在局部化后消失，几份局部化中的元素也可能无法同时由原模中的元素给出。支集挠子函子的导出把这些现象按次数区分，有限 Čech 复形则给出可直接计算的表示。


% 纲要内容（仅作写作计划，尚非教材正文）：
%
% * 支集挠子函子（即本章的 $I$-挠子函子 $\Gamma_I$，英文 torsion functor；不是 torsor）
% * 内射消解负责导出定义，有限 Čech 复形负责计算；双复形两种计算、二元算例与平坦基变换
% * 逐元素幂挠与统一幂零化的区别；正则元商的深度归纳与参数个数给出的高次消失
%
% ---
%

\subsection{支集挠子函子}
\label{b3:subsec:2201:01}

局部化 \(M_f\) 保留在 \(f\ne0\) 的地方能够看见的元素；映射 \(M\to M_f\) 的核则记录离开 \(V(f)\) 后消失的部分。对于一个由多个元素生成的理想，除了这种核，还会出现局部数据不能拼合的障碍。局部上同调把这些障碍按次数组织起来。

\begin{definition}\label{b3:def:support-torsion-functor}
设 \(R\) 为交换 Noether 环，\(I\subseteq R\) 为理想，\(M\) 为 \(R\) 模。令
\[
\Gamma_I(M)\coloneqq\{m\in M:\text{存在}\;n\ge1\;\text{使}\;I^nm=0\}.
\]
把 \(M\) 送到这个子模、把模同态送到其限制，得到的函子称为\term[zhijinaozihanzi]{支集挠子函子}[torsion functor]，也称 \(I\)-挠子函子。
\end{definition}
\symidx{\(\Gamma_I(M)\)：支集挠子模}

这里的“支集”有直接含义：\(m\in\Gamma_I(M)\) 等价于循环模 \(Rm\) 的支集包含于 \(V(I)\)。事实上，循环模的支集为 \(V(\operatorname{Ann}_R(m))\)，包含关系等价于 \(I\subseteq\sqrt{\operatorname{Ann}_R(m)}\)；有限组生成元各有一个幂落入零化子，取足够大的总次数就得到 \(I^nm=0\)。定义中的量词是“对每个元素存在一个指数”，这个指数可以随元素改变；即使 \(\Gamma_I(M)=M\)，也未必存在一个统一的 \(N\) 使 \(I^NM=0\)。

\ProseParagraphBreak
若 \(0\to M'\to M\to M''\) 正合，则 \(0\to\Gamma_I(M')\to\Gamma_I(M)\to\Gamma_I(M'')\) 正合：前两处只是在原模中检验相同的零化条件。因此它是左正合函子；但右端未必满射。例如对域 \(k\)，在 \(R=k[x]\)、\(I=(x)\) 下，满射 \(R\to R/(x)\) 的两端挠子模分别为零与 \(k\)。为记录这一失败，需要把原函子延伸到正次数。

\ProseParagraphBreak
内射消解使这一延伸成为 \(\Gamma_I\) 的右导出函子，并由比较映射和短正合列的消解得到函子性与长正合列。随后要证明的 Čech 比较定理承担另一项任务：把这个由内射消解定义的对象写成有限的局部化复形，供具体计算使用。

\begin{definition}\label{b3:def:local-cohomology}
设 \(R\) 为交换 Noether 环，\(I\subseteq R\) 为理想，\(M\) 为 \(R\) 模。选一个内射消解
\[
0\longrightarrow M\longrightarrow E^0\xrightarrow{\mathrm{d}^0}E^1
\xrightarrow{\mathrm{d}^1}E^2\longrightarrow\cdots.
\]
定义 \(M\) 关于 \(I\) 的第 \(i\) 个\term[jubushangtongdiao]{局部上同调模}[local cohomology module]为
\[
H_I^i(M)\coloneqq H^i(\Gamma_I(E^\bullet))\quad(i\ge0),
\]
并令负次数的局部上同调为零。亦记 \(H_I^i=R^i\Gamma_I\)，即 \(\Gamma_I\) 的第 \(i\) 个右导出函子。
\end{definition}
\symidx{\(H_I^i(M)\)：局部上同调模}

第二卷第7.2节的足够多内射模定理已保证这种消解存在：先嵌入 \(E^0\)，再把余核嵌入 \(E^1\)，逐次进行。消解之间的比较也由目标项的内射性构造。给定 \(M\to N\)，先将它延拓到次数零；在下一次数，先在前一微分的像上规定映射，再延拓到整个项。链映射关系保证这个规定与核相容，因而可逐次完成。

\ProseParagraphBreak
若两条比较链映射提升同一个 \(M\to N\)，它们的差在次数零的 \(M\) 上为零，故先经过 \(E^0/M\cong\operatorname{im}(E^0\to E^1)\) 分解，再利用目标次数零项的内射性将该映射延拓到 \(E^1\)，得到第一项同伦。逐次减去已构造的同伦项，在下一次数重复，得到差为 \(\mathrm{d}h+h\mathrm{d}\)。加性函子 \(\Gamma_I\) 将这个同伦限制到挠子子模，所以两条链映射诱导相同的上同调映射。将此用于恒等映射的两次延拓，便得到不同消解之间互逆的同构；用于一般模同态，则得到所述函子性。

\begin{proposition}\label{b3:prop:local-cohomology-basic}
设 \(R\) 为交换 Noether 环，\(I,J\subseteq R\) 为理想，\(M\) 为 \(R\) 模。则 \(H_I^0(M)=\Gamma_I(M)\)。短正合的 \(R\) 模列给出各阶 \(H_I^i\) 的长正合列；若 \(\sqrt I=\sqrt J\)，则 \(H_I^i(M)\cong H_J^i(M)\) 自然成立。
\end{proposition}
\begin{proof}
零次结论就是左正合性。对于 \(0\to M'\to M\to M''\to0\)，给两端选内射嵌入 \(M'\to E'\)、\(M''\to E''\)。把第一条映射延拓到 \(M\)，再合上 \(M\to M''\to E''\)，得到单射 \(M\to E'\oplus E''\)。其余核与两端的余核成短正合列，故可重复此过程，得到逐项分裂的三份内射消解。加性函子 \(\Gamma_I\) 保持这些逐项分裂列，圈的提升、取微分和取商遂给出长正合列，连接同态的良定义及正合性与\cref{b3:lem:ext-basic-properties}中的复形论证相同。

若两理想根相等，有限生成性给出 \(I^a\subseteq J\)、\(J^b\subseteq I\)。因此两种挠子条件相同，\(\Gamma_I=\Gamma_J\)；对同一内射消解取上同调即可。
\end{proof}

\subsection{Čech 复形与导出函子的比较}
\label{b3:subsec:2201:02}

内射消解适合定义，但往往不便计算。对于 \(I=(f_1,\ldots,f_r)\)，仅用局部化就可以写出一个有限复形。

\begin{definition}\label{b3:def:cech-complex}
设 \(R\) 为交换环，\(f_1,\ldots,f_r\in R\)，\(M\) 为 \(R\) 模。对任意 \(f\in R\)，取两项上链复形 \(C^\bullet(f;R)=[R\to R_f]\)，次数为零和一。令
\[
C^\bullet(f_1,\ldots,f_r;M)
=M\otimes_R C^\bullet(f_1;R)\otimes_R\cdots\otimes_R C^\bullet(f_r;R),
\]
张量微分在越过次数 \(q\) 的因子时带符号 \((-1)^q\)。这个复形称为所给元素组关于 \(M\) 的\term[Cechfuxing]{Čech 复形}[Čech complex]。
\end{definition}
\symidx{\(C^\bullet(f_1,\ldots,f_r;M)\)：Čech 复形}

具体地，第 \(p\) 项为所有 \(M_{f_{i_1}\cdots f_{i_p}}\) 的有限直和；次数零为 \(M\)。微分是带交错符号的局部化映射。两元情形为
\[
0\longrightarrow M\xrightarrow{m\mapsto(m,m)}M_f\oplus M_g
\xrightarrow{(a,b)\mapsto b-a}M_{fg}\longrightarrow0.
\]
这里零次上同调是那些在两份局部化中都消失的原元素。一次上同调由在 \(M_{fg}\) 中相等的元素对组成，再商去来自 \(M\) 的对角像，因而测量相容的局部数据能否由同一个原元素给出。二次上同调则是 \(M_{fg}/(\operatorname{im}M_f+\operatorname{im}M_g)\)，测量交叠处的元素能否写成两份局部数据之差。若以空元素组生成零理想，则复形只有次数零的 \(M\)。有限复形并不意味着其上同调有限生成：\(k[x,x^{-1}]/k[x]\) 已经是一个反例。

\ProseParagraphBreak
同一组局部化还会在\secref{b3:sec:2405}的开覆盖计算中出现，但起始次数不同。对 \(I=(f,g)\)，两种排列为
\[
\begin{array}{c|ccccc}
\text{次数}&0&&1&&2\\\hline
\text{支集复形}&M&\longrightarrow&M_f\oplus M_g&\longrightarrow&M_{fg}\\
\text{开覆盖复形}&M_f\oplus M_g&\longrightarrow&M_{fg}&\longrightarrow&0
\end{array}
\]
各行的箭头都是上面的局部化映射。第二行删去第一行的 \(M\)，再将其余项左移一位，因此正次数的上同调与第一行相差一次。这里的两行都只是局部化模构成的复形；第二行作为 \(D(f)\cup D(g)\) 开覆盖模型的解释，以及它的零次项与一般开集截面的关系，将在\secref{b3:sec:2405}定义相关几何对象后写成正合列。

\begin{lemma}\label{b3:lem:injective-cech-acyclic}
设 \(R\) 为交换 Noether 环，\(I\subseteq R\) 为理想，\(E\) 为内射 \(R\) 模。则 \(\Gamma_I(E)\) 仍内射。对任意 \(f\in R\)，映射 \(E\to E_f\) 满射。因此，任意生成组 \(I=(f_1,\ldots,f_r)\) 满足
\[
H^0(C^\bullet(f_1,\ldots,f_r;E))=\Gamma_I(E),\qquad
H^i(C^\bullet(f_1,\ldots,f_r;E))=0\quad(i>0).
\]
\end{lemma}
\begin{proof}
先证明挠子部分仍然内射，再证明局部化映射满射。这使一元 Čech 复形在内射模上只留下零次上同调；由于挠子部分仍内射，多元情形便能逐个加入其余平坦的局部化因子。

用第二卷第7.2节已经证明的 Baer 判据检验第一项。给理想 \(\mathfrak a\subseteq R\) 和同态 \(u:\mathfrak a\to\Gamma_I(E)\)。因 \(\mathfrak a\) 有限生成，某个 \(I^n\) 零化全部像。对有限生成的子模 \(\mathfrak a\subseteq R\) 应用\cref{b3:thm:artin-rees}，给出足够大的 \(N\)，使
\(\mathfrak a\cap I^N\subseteq I^n\mathfrak a\)。Artin–Rees 在这里把“像被 \(I^n\) 零化”转成一个统一的交理想条件，使延拓能够在 \(I^N\) 上取零。于是 \(u\) 在该交上为零，可延拓为
\(\mathfrak a+I^N\to E\)、\(a+b\mapsto u(a)\)。再用 \(E\) 的内射性延拓到 \(R\)。设一的像为 \(e\)，则 \(I^Ne=0\)，所以这个延拓实际取值于 \(\Gamma_I(E)\)。

为了证明满射性，给 \(e/f^n\in E_f\)。环的升链条件使 \(0:_Rf^k=0:_Rf^{k+n}\) 对充分大的 \(k\) 成立。因而
\[
f^{n+k}R\longrightarrow E,\qquad f^{n+k}a\longmapsto f^kae
\]
良定义。延拓到 \(R\to E\)，设一的像为 \(z\)，则 \(f^{n+k}z=f^ke\)，即 \(z/1=e/f^n\)。当 \(f\) 幂零时局部化为零，也包含在此证明中。

现在对 \(r\) 归纳。已证的一元情形给出短正合列
\[
0\longrightarrow\Gamma_{(f_1)}(E)\longrightarrow E\longrightarrow E_{f_1}\longrightarrow0.
\]
因而次数零的 \(\Gamma_{(f_1)}(E)\) 嵌入两项复形 \([E\to E_{f_1}]\)，商复形正合。将这一复形短正合列与余下 Čech 复形的每个平坦项张量，得到逐项短正合列，且商双复形的各行仍正合，故总复形的上同调等于
\(C^\bullet(f_2,\ldots,f_r;\Gamma_{(f_1)}(E))\) 的上同调。这里利用这些正合行逐次消去圈的分量，有限项数保证终止；所用微分符号正是张量符号。第一项已经证明中间模内射，归纳假设可用。最后，先被 \(f_1\) 的某幂零化、再被其余生成元各自的某幂零化，等价于被 \(I\) 的某幂零化，所以零次为 \(\Gamma_I(E)\)。
\end{proof}

\begin{theorem}\label{b3:thm:cech-local-cohomology}
设 \(R\) 为交换 Noether 环，\(I=(f_1,\ldots,f_r)\)，\(M\) 为任意 \(R\) 模。则有自然同构
\[
H_I^i(M)\cong H^i(C^\bullet(f_1,\ldots,f_r;M))\qquad(i\ge0).
\]
特别地，右侧不依赖所选生成组，每个 \(H_I^i(M)\) 都是 \(I\)-挠子模。
若 \(R\to B\) 为交换 Noether 环之间的平坦同态，则对所有 \(i\ge0\) 有自然同构
\[
H_I^i(M)\otimes_RB\cong H_{IB}^i(M\otimes_RB).
\]
具体地，对域 \(k\)、\(R=k[x,y]\) 和 \(I=(x,y)\)，有
\[
H_I^0(R)=H_I^1(R)=0,\qquad
H_I^2(R)=\bigoplus_{a,b\ge1}k x^{-a}y^{-b},
\]
且其余次数为零。这里单项式表示它在 \(R_{xy}/(R_x+R_y)\) 中的商类，所列直和给出 \(k\) 向量空间的基，\(R\) 作用由这个余核继承。
\end{theorem}
\begin{proof}
双复形要把两种计算放在同一个总复形中：先计算内射消解方向，得到 \(M\) 的 Čech 复形；先计算 Čech 方向，则只剩定义中的挠子复形。

选内射消解 \(M\to E^\bullet\)，令
\[
\mathcal B^{p,q}=C^p(f_1,\ldots,f_r;R)\otimes_R E^q,
\qquad 0\le p\le r,\quad q\ge0.
\]
总次数为 \(p+q\)，在 \((p,q)\) 项上总微分为 \(\mathrm{d}_C+(-1)^p\mathrm{d}_E\)。先沿 \(q\) 方向取上同调：每个 Čech 项平坦，所以只剩 \(q=0\) 行的 \(C^\bullet(f_1,\ldots,f_r;M)\)。先沿 \(p\) 方向取上同调：上一个引理使每行只剩 \(p=0\) 列的 \(\Gamma_I(E^q)\)，从而得到定义中的复形。

两条边缘增广分别给出复形单射
\[
C^\bullet(f;M)\longrightarrow\operatorname{Tot}(\mathcal B),\qquad
\Gamma_I(E^\bullet)\longrightarrow\operatorname{Tot}(\mathcal B).
\]
对第一条取商后，每列底项换成 \(\mathcal B^{p,0}/(C^p(f;R)\otimes_RM)\)，其余项不变，各列正合。对第二条取商后，每行左端换成 \(\mathcal B^{0,q}/\Gamma_I(E^q)\)，其余项不变，各行正合。两个商总复形都正合，这可以直接作有限消元：列正合时，对次数 \(n\) 的圈选最小的非零横指标 \(p\)，则该分量 \(z_{p,q}\) 满足 \(\mathrm d_Ez_{p,q}=0\)。若 \(q=0\)，列底端正合已使它为零；否则选 \(w\) 使 \((-1)^p\mathrm d_Ew=z_{p,q}\)，减去总边界 \(\mathrm dw\) 便消去该分量，新增修正只在 \((p+1,q-1)\)。横指标最多到 \(n\)，所以有限步结束。行正合时交换两个指标作相同消元。由这两个正合商复形，边缘增广都诱导上同调同构，以上两种计算因而给出同一个总上同调。所有映射来自局部化及消解比较，所以同构自然。不同生成组的 Čech 复形可以有不同项数和微分；这里的生成组无关性指它们的上同调自然同构。

把 Čech 复形局部化于任一 \(f_j\)，其中的两项因子变成恒等映射 \(R_{f_j}\to R_{f_j}\)，可缩，故所有上同调局部化后为零。对于上同调中的一个元素，每个 \(f_j\) 的某个幂都将它零化。有限组生成元的任意足够高次乘积必含这样的幂，故 \(I\) 的某幂零化该元素。

对基变换公式，任意环同态 \(R\to B\) 都给出与局部化微分相容的逐项同构
\[
M_{f_{j_1}\cdots f_{j_p}}\otimes_RB
\cong(M\otimes_RB)_{f_{j_1}\cdots f_{j_p}}.
\]
因此 \(C^\bullet(f;M)\otimes_RB\cong C^\bullet(f;M\otimes_RB)\)，右边在 \(B\) 上计算。平坦性在下一步使用：张量保持微分的核、像及其商，所以
\(H^i(C^\bullet)\otimes_RB\cong H^i(C^\bullet\otimes_RB)\)。在两环上应用已证的 Čech 比较，就得到所述公式。

最后计算二元例子。将 \(R_x,R_y\) 看作 \(R_{xy}=k[x,x^{-1},y,y^{-1}]\) 的子环。零次核为零；在 \(R_x\) 中 \(y\) 的指数非负，在 \(R_y\) 中 \(x\) 的指数非负，故 \(R_x\cap R_y=R\)。一次复形的核因而恰为 \(R\) 的对角像，\(H_I^1(R)=0\)。二次余核 \(R_{xy}/(R_x+R_y)\) 商去了至少有一个指数非负的全部 Laurent 单项式；剩下的 \(x^{-a}y^{-b}\)、\(a,b\ge1\)，线性独立并张成该余核，所以构成所列基。复形长度为二，其余次数消失。
\end{proof}

例如 \(R=k[x]\)、\(I=(x)\) 时，\(H_I^0(R)=0\)，而
\[
H_I^1(R)=k[x,x^{-1}]/k[x]
=\bigoplus_{j\ge1}k x^{-j}.
\]
乘 \(x\) 把 \(x^{-j}\) 送到 \(x^{-j+1}\)，其中 \(x^{-1}\) 被送到零。每个元素被某个 \(x\) 的幂零化，但整个模不被统一的幂零化，也不有限生成。这是局部上同调与有限生成模上的 Ext 之间的一处显著差别。

\ProseParagraphBreak
这个模还具体记录了节首所见的满射失败。对短正合列
\[
0\longrightarrow R\xrightarrow{x}R\longrightarrow k\longrightarrow0,
\qquad R=k[x],
\]
关于 \((x)\) 的局部上同调长正合列含有
\[
0\longrightarrow k\xrightarrow{\delta}H^1_{(x)}(R)
\xrightarrow{x}H^1_{(x)}(R)\longrightarrow0.
\]
连接同态把 \(1\in k\) 送到 \(x^{-1}+R\)：先在中间的 \(R\) 中提升为 \(1\)，它的 Čech 微分是 \(1\in R_x\)，再沿前一项的乘 \(x\) 映射提升为 \(x^{-1}\)。因而 \(\ker(x:H^1_{(x)}(R)\to H^1_{(x)}(R))=k x^{-1}\)，恰好补上原来 \(\Gamma_{(x)}(R)=0\) 无法满射到 \(\Gamma_{(x)}(k)=k\) 的部分。高次局部上同调在这里把左正合函子的失效变成了可计算的正合列。

\subsection{局部上同调的消失与深度}
\label{b3:subsec:2201:03}

深度要求连续选取非零因子，局部上同调则测量支集在闭点的障碍。前几个上同调消失，恰好允许逐次选取相应长度的正则序列。

\begin{theorem}\label{b3:thm:depth-local-cohomology}
设 \(R\) 为交换 Noether 局部环，极大理想为 \(\mathfrak m\)，\(M\ne0\) 为有限生成的 \(R\) 模。则
\[
\operatorname{depth}_R M
=\min\{i\ge0:H_{\mathfrak m}^i(M)\ne0\}.
\]
\end{theorem}
\begin{proof}
商去一个正则元素使深度降一。长正合列将利用低次局部上同调上的乘法单射，结合逐元素的幂零化，使低次消失；再把商模的首次非零项送入原模高一次的上同调。

对 \(t=\operatorname{depth}_R M\) 归纳。若 \(t=0\)，\cref{b3:cor:depth-zero-socle}给出非零的 \(\operatorname{Soc}(M)\)，其每个元素被 \(\mathfrak m\) 零化，故 \(H_{\mathfrak m}^0(M)\ne0\)。

若 \(t>0\)，取 \(M\) 正则元素 \(x\in\mathfrak m\)。商 \(\bar M=M/xM\) 非零，且由\cref{b3:cor:regular-quotient-depth}有深度 \(t-1\)。原短正合列的局部上同调长正合列中含有
\[
H_{\mathfrak m}^{i-1}(\bar M)\longrightarrow
H_{\mathfrak m}^i(M)\xrightarrow{x}H_{\mathfrak m}^i(M).
\]
当 \(i<t\) 时左端为零，所以乘 \(x\) 在中间模上单射。该模为 \(\mathfrak m\)-挠子模；一个既被某个 \(x\) 的幂零化、又处在 \(x\) 单射的模中的元素只能为零。因此这些次数全部消失。取 \(i=t\)，长正合列给出非零模
\(H_{\mathfrak m}^{t-1}(\bar M)\) 到 \(H_{\mathfrak m}^t(M)\) 的单射，因为前一项 \(H_{\mathfrak m}^{t-1}(M)\) 已为零。于是第 \(t\) 次非零。
\end{proof}

\begin{theorem}\label{b3:thm:local-cohomology-dimension-vanishing}
设 \(R\) 为交换 Noether 局部环，极大理想为 \(\mathfrak m\)，\(M\ne0\) 为有限生成模，\(d=\dim_RM\)。则
\[
H_{\mathfrak m}^i(M)=0\quad(i>d).
\]
因此，\(M\) 为 Cohen--Macaulay 模当且仅当其局部上同调仅在次数 \(d\) 非零。
\end{theorem}
\begin{proof}
维数 \(d\) 允许选取 \(d\) 个参数，把支集切到闭点；这些参数给出的长度 \(d\) 的 Čech 复形就足以计算极大理想支集上的上同调，因此更高次数没有复形项。

取 \(M\) 的参数系 \(x_1,\ldots,x_d\)，其存在性由\cref{b3:thm:system-of-parameters}应用于 \(A=R/\operatorname{Ann}_RM\) 保证。在 \(A\) 中，所生成理想的根就是极大理想。Čech 复形的项只涉及 \(M\) 的局部化，所以不论把 \(M\) 看作 \(R\) 模还是 \(A\) 模，其复形相同。先在 \(A\) 上应用理想根无关性，再使用 Čech 比较定理，便可用 \(x_1,\ldots,x_d\) 的长度 \(d\) 复形计算 \(H_{\mathfrak m}^i(M)\)。更高次数没有项，因而消失。

若 \(M\) 为 CM 模，深度为 \(d\)，前一定理使低于 \(d\) 的次数消失并使第 \(d\) 次非零；反向同样由首次非零次数得到深度 \(d\)。
\end{proof}

对域 \(k\)，以 \(R=k[\![x,y]\!]/(x^2,xy)\) 为例，\((y)\) 的根为极大理想 \(\mathfrak m=(x,y)\)。每个元素唯一写成 \(a(y)+bx\)，其中 \(a(y)\in k[\![y]\!]\)、\(b\in k\)。由 \(yx=0\)，映射 \(R\to R_y=k(\!(y)\!)\) 的核是 \(kx\)，像是 \(k[\![y]\!]\)。因此
\[
H_{\mathfrak m}^0(R)=kx,\qquad
H_{\mathfrak m}^1(R)=k(\!(y)\!)/k[\![y]\!].
\]
环维数为一，深度却为零。被嵌入的零维部分 \(kx\) 在零次局部上同调中保留下来，一维主分量则在一次留下负幂的类，两部分信息出现在不同次数。CM 条件所要求的是全部局部上同调集中在模维数这一次数；下一节的对偶将说明，这种集中性为什么能用一个有限生成模表达。

% !TeX root = ../../Book3.tex
% 纲要标题：### 22.2 Matlis 对偶、典范模与局部对偶
% 纲要位置：计划纲要.md，第 2675 行。

\section{Matlis 对偶、典范模与局部对偶}
\label{b3:sec:2202}

局部上同调通常不是有限生成模；以环本身为值的 Hom 对偶甚至可能检测不到非零元素。内射包络与 Matlis 对偶提供合适的取值对象。若一个完备局部环来自完备正则局部环的满射表示，就能把它的局部上同调转为环境环上的 Ext；Cohen–Macaulay 性使对偶信息集中在一个次数，进而用典范模写出环自身上的对偶公式。

\ProseParagraphBreak
这些计算从满射 \(S\twoheadrightarrow R\) 出发，其中 \(S\) 是交换完备 Noether 正则局部环。幂级数商环天然带有这样的环境环；一般交换完备 Noether 局部环的相应满射来自\appref{b3:app:D}的 Cohen 结构定理。

% 纲要内容（仅作写作计划，尚非教材正文）：
%
% * 从离散赋值环的环值对偶失效，说明内射包络作为取值对象的作用
% * 内射包络、有限长度对偶与完备 Noether 局部环上的 Matlis 对偶；非完备双对偶反例及 Artin 不等于有限长度
% * 极小内射消解的基座计算；由 Koszul 对偶与内射模的极大理想幂挠部分计算正则局部环的最高局部上同调
% * 给定满射 $S\twoheadrightarrow R$，$S$ 为 $n$ 维完备 Noether 正则局部环、$d=\dim R$，构造典范模 $\omega_R=\operatorname{Ext}_S^{n-d}(R,S)$；此存在性不要求 $R$ 为 CM 环
% * 若 $R$ 为 CM 环，环境环上的 Ext 集中于 $n-d$ 次；换环引理给出单个典范模的局部对偶公式
% * 对有限生成的 $R$ 模比较局部上同调与 $\operatorname{Ext}_R^{d-i}(-,\omega_R)$ 的消失区间；非 CM 环以多个非零 Ext 次数说明单模公式的限制
% * 一般完备局部环上这类满射表示的存在性由附录 D 的 Cohen 结构定理说明
%
% ---
%

\subsection{环值对偶的不足}
\label{b3:subsec:2202:01}

设 \(k\) 为域，先看离散赋值环 \(A=k[\![t]\!]\)。由两项 Čech 复形得到
\[
H^1_{(t)}(A)=k(\!(t)\!)/k[\![t]\!].
\]
右端含有任意高的负幂，不能由有限个元素生成。普通的环值对偶在此为零。事实上，令 \(K=k(\!(t)\!)\)，则
\[
 \operatorname{Hom}_A(K/A,A)=0.
\]
这是因为每个 \(K/A\) 中的元素都被 \(t\) 的某个幂零化，而 \(A\) 没有非零的 \(t\)-挠元。因此以 \(A\) 为值的对偶无法检测这个非零模，需要改变取值对象。

\ProseParagraphBreak
这个对偶的取值对象将是剩余域的内射包络。它既能接受子模上同态的延拓，又能检测任意非零元素；在完备 Noether 局部环上，它把有限生成模与 Artin 模互相转换。因此，原来的局部上同调模虽然通常不有限生成，它的对偶却有可能由有限组生成元及关系描述。

\ProseParagraphBreak
给定满射 \(S\twoheadrightarrow R\)，其中 \(S\) 是 \(n\) 维交换完备 Noether 正则局部环，\(R\) 是 \(d\) 维商环，与最高次局部上同调对应的有限生成模是
\[
\operatorname{Ext}_S^{n-d}(R,S).
\]
它能够用 \(R\) 的有限长度、各项有限秩自由的 \(S\) 消解计算；下面的\cref{b3:def:canonical-module-regular-quotient}将从环自身的最高次局部上同调定义典范模，\cref{b3:prop:canonical-module-quotient}再识别它与这个环境 Ext 模。余维 \(n-d\) 出现在这里，是因为正则局部环的最高局部上同调位于次数 \(n\)，而 \(R\) 的维数为 \(d\)。若 \(R\) 是 Cohen–Macaulay 环，其余局部上同调消失，整个对偶公式便集中到这一个模上。

\subsection{离散赋值环的内射对偶}
\label{b3:subsec:2202:02}

先把上述离散赋值环的例子算得更具体。令 \(A=k[\![t]\!]\)、\(K=k(\!(t)\!)\)、\(E=K/A\)。第二卷第7.2节的 Baer 判据说明 \(E\) 是内射 \(A\) 模：任意非零理想为 \(t^rA\)，若 \(f(t^r)=a+A\)，则取 \(g(1)=t^{-r}a+A\)，就得到延拓 \(g:A\rightarrow E\)。零理想的映射用零映射延拓。

\ProseParagraphBreak
\(E\) 中每个非零元素都能乘以某个 \(t\) 的幂，得到非零的 \(t^{-1}\) 项。因此子模 \(kt^{-1}\cong k\) 与 \(E\) 的任意非零子模相交非零。这说明内射性没有加入与剩余域毫无关系的直和因子；这正是下文“内射包络”中包络条件的含义。

\ProseParagraphBreak
记 \(D_A(N)=\operatorname{Hom}_A(N,E)\)。这里有
\[
D_A(A)=E,
\qquad
D_A(A/t^rA)=(0:_Et^r)=t^{-r}A/A\cong A/t^rA.
\]
最后一个同构把 \(a+t^rA\) 送到 \(at^{-r}+A\)。这使有限长度模的对偶仍是有限长度模，而自由模的对偶是非有限生成的 \(E\)。离散赋值环的计算对应于一般局部对偶公式
\[
H^i_{\mathfrak m}(M)
\cong D_R\bigl(\operatorname{Ext}_R^{d-i}(M,\omega_R)\bigr),
\]
其中 \(R=S/I\) 是来自完备 Noether 正则局部环 \(S\) 的 \(d\) 维 Cohen–Macaulay 商环，\(\mathfrak m\) 是其极大理想，\(M\) 是有限生成的 \(R\) 模；\(\omega_R\) 是将在\cref{b3:def:canonical-module-regular-quotient}定义的典范模。完整陈述和证明见\cref{b3:thm:canonical-module-local-duality}。

\subsection{局部对偶的适用条件}
\label{b3:subsec:2202:03}

局部对偶中的三种条件各有作用。Noether 性保证有限生成模有有限展示，并使内射包络的元素能用极大理想的幂控制。完备性使取两次对偶后的环仍是原环；若环不完备，首先出现的是它的完备化。Cohen–Macaulay 性则把环自身的局部上同调集中在唯一的次数 \(d\)，使这些对偶公式都可以用一个模 \(\omega_R\) 表示。

\ProseParagraphBreak
这些条件不能由“局部环”三个字替代。一般非 Cohen–Macaulay 环可能同时有几个非零局部上同调模，不能把上述公式中的 \(\omega_R\) 原样用于所有次数。即使某个环具有典范模，也不表示它的全部对偶信息已经集中在该模上。\subsecref{b3:subsec:2202:06}将给出一个含有嵌入点的幂级数商环，直接检验这种失败。

\ProseParagraphBreak
若只给出一个交换完备 Noether 局部环 \(R\)，Cohen 结构定理提供来自完备正则局部环的满射表示；相应的系数域与系数环构造见\appref{b3:app:D}。

\subsection{内射包络与有限生成模、Artin 模的 Matlis 对偶}
\label{b3:subsec:2202:04}

先说明如何在所有可能的内射扩张中选择一个尽量小的对象。这里使用第二卷第7.2节已证明的两个结果：每个模都可嵌入内射模；Baer 判据用理想上的同态延拓检验内射性。内射包络及以下对偶定理在本节证明。松村英之《Commutative Ring Theory》的对应论述见\cite[\S 18, Theorems 18.4--18.7, pp. 145--150]{Matsumura1989CommutativeRingTheory}。

\begin{definition}\label{b3:def:injective-hull}
设 \(R\) 为交换环，\(N\subseteq M\) 为 \(R\) 模。若 \(M\) 的每个非零子模都与 \(N\) 相交非零，则此包含为\term[benzhi-kuozhang]{本质扩张}[essential extension]。若 \(M\) 还为内射模，就称 \(N\subseteq M\) 为 \(N\) 的\term[neishe-baoluo]{内射包络}[injective hull]，记一个这样的模为 \(E_R(N)\)。
\end{definition}
\symidx{\(E_R(N)\)：模 \(N\) 的内射包络}

本质性有一个常用的检验形式：若映射 \(f:M\rightarrow L\) 在 \(N\) 上单射，则它在整个 \(M\) 上单射，因为非零的 \(\ker f\) 必与 \(N\) 相交。

\ProseParagraphBreak
本质扩张要求每个非零子模都遇到原来的 \(N\)，所以不能额外加入一个与 \(N\) 交零的非零直和因子。前面的 \(kt^{-1}\subset k(\!(t)\!)/k[\![t]\!]\) 就体现了这一点：任意非零子模都有元素经乘法落到这份剩余域中。下面的存在性证明先在一个内射模内取极大本质扩张，再证明它本身内射；唯一性同样由本质性排除补因子。

\begin{proposition}\label{b3:prop:injective-hull-existence}
设 \(R\) 为交换环，\(N\) 为 \(R\) 模。\(N\) 有内射包络，且任意两个内射包络之间存在限制在 \(N\) 上为恒等的同构。这里不声称这个同构唯一。
\end{proposition}
\begin{proof}
把 \(N\) 嵌入内射模 \(J\)。在 \(J\) 中所有本质包含 \(N\) 的子模中按包含排序；一条链的并仍本质包含 \(N\)，因为并中的非零元素已落在链的某个成员中。Zorn 引理给出极大者 \(E\)。

\(E\) 没有真本质扩张。事实上，若 \(E\subseteq X\) 本质，则本质性的传递性使 \(N\subseteq X\) 本质；内射性把 \(E\hookrightarrow J\) 延拓为 \(X\rightarrow J\)。此映射的核与 \(E\) 交零，故为零。其像在 \(J\) 中仍本质包含 \(N\)，极大性使像等于 \(E\)。映射又在 \(E\) 上为恒等，故 \(X=E\)。

再把 \(E\) 嵌入内射模 \(J'\)，选取与 \(E\) 交零的极大子模 \(L\subseteq J'\)。商模 \(J'/L\) 本质包含 \(E\)：若一个非零子模与 \(E\) 的像交零，其原像就是真包含 \(L\) 且与 \(E\) 交零的子模，矛盾。由上一段，\(J'/L=E\) 的像，故 \(J'=E\oplus L\)。内射模的直和因子内射，得存在性。

若 \(E,E'\) 都是包络，由 \(E'\) 的内射性延拓 \(N\) 上的恒等为 \(u:E\rightarrow E'\)。本质性使 \(u\) 单射，而内射子模 \(u(E)\) 在 \(E'\) 中分裂。补因子与 \(N\) 交零，本质性迫使补因子为零，所以 \(u\) 是同构。
\end{proof}

以下固定交换 Noether 局部环 \((R,\mathfrak m,k)\) 及包络 \(k\subseteq E=E_R(k)\)。

\ProseParagraphBreak
其基座只有指定的一份剩余域，每个元素又被极大理想的某个幂零化。这两点控制 \(E\) 中的有限层；第三项元素检测性质则保证，以 \(E\) 为值的对偶不会把非零模全部送到零。

\begin{proposition}\label{b3:prop:injective-hull-local-properties}
设 \((R,\mathfrak m,k)\) 为交换 Noether 局部环，\(E=E_R(k)\)。则
\[
\operatorname{Soc}_R E=(0:_E\mathfrak m)=k,
\qquad
E=\bigcup_{r\ge1}(0:_E\mathfrak m^r).
\]
此外，对任意 \(R\) 模 \(M\) 及 \(0\ne x\in M\)，存在 \(f:M\rightarrow E\) 使 \(f(x)\ne0\)。
\end{proposition}
\begin{proof}
\((0:_E\mathfrak m)\) 是 \(k\) 向量空间；它的每条一维子空间都须与指定的 \(k\) 相交，所以它恰为该 \(k\)。

固定 \(e\in E\)。若 \(e=0\)，它已被 \(\mathfrak m\) 零化；以下设 \(e\ne0\)。若 \(\mathfrak p\in\operatorname{Ass}_R(Re)\)，则 \(Re\) 有一个同构于 \(R/\mathfrak p\) 的子模。该子模与 \(k\) 相交非零；而 \(R/\mathfrak p\) 的每个非零元素的零化子都是 \(\mathfrak p\)，相交元素的零化子又是 \(\mathfrak m\)，故 \(\mathfrak p=\mathfrak m\)。对由 \(e\) 生成的循环模 \(Re\) 应用\cref{b3:prop:minimal-primes-associated}，得到 \(\sqrt{\operatorname{Ann}_R(e)}=\mathfrak m\)。\(\mathfrak m\) 有限生成，因而某个 \(\mathfrak m^r\) 包含于该零化子，得并集公式。

对非零 \(x\)，\(\operatorname{Ann}_R(x)\) 是真理想。复合
\(Rx\cong R/\operatorname{Ann}_R(x)\twoheadrightarrow k\hookrightarrow E\)
把 \(x\) 送到非零元素，再用内射性将它延拓到 \(M\)。
\end{proof}

\begin{definition}\label{b3:def:matlis-dual}
设 \((R,\mathfrak m,k)\) 为交换 Noether 局部环，固定 \(E=E_R(k)\)。对 \(R\) 模 \(M\)，模
\[
D_R(M)\coloneqq\operatorname{Hom}_R(M,E)
\]
为它的\term[Matlis-duiou]{Matlis 对偶}[Matlis dual]。同态 \(u:M\rightarrow N\) 对应预复合映射 \(D_R(u):D_R(N)\rightarrow D_R(M)\)。
\end{definition}
\symidx{\(D_R(M)\)：模 \(M\) 的 Matlis 对偶}

\(D_R\) 是反变正合函子，因为 \(E\) 内射。上一命题还说明 \(D_R(M)=0\) 蕴含 \(M=0\)。求值映射
\[
\delta_M:M\longrightarrow D_RD_R(M),
\qquad \delta_M(x)(f)=f(x)
\]
对每个模都单射。不过，满射性需要限制模的类别。

\begin{lemma}\label{b3:lem:matlis-finite-length}
设 \((R,\mathfrak m,k)\) 为交换 Noether 局部环，\(D_R(-)=\operatorname{Hom}_R(-,E_R(k))\)。若 \(M\) 为有限长度 \(R\) 模，\(\delta_M:m\mapsto(\varphi\mapsto\varphi(m))\) 为求值映射，则
\[
\ell_R D_R(M)=\ell_R M,
\qquad \delta_M:M\xrightarrow{\sim}D_RD_R(M).
\]
并且乘法给出环同构
\[
\widehat R\xrightarrow{\sim}\operatorname{End}_R\bigl(E_R(k)\bigr),
\qquad \widehat R=\varprojlim_r R/\mathfrak m^r.
\]
具体地，若 \(k\) 为域、\(R_0=k[t]_{(t)}\)、\(E_0=k(t)/R_0\)、\(D_{R_0}=\operatorname{Hom}_{R_0}(-,E_0)\)，则 \(E_0\) 为 \(k\) 的内射包络，且
\[
D_{R_0}D_{R_0}(R_0)\cong k[\![t]\!].
\]
自然求值映射对应完备化映射 \(R_0\to k[\![t]\!]\)，它不是满射。
\end{lemma}
\begin{proof}
先由 \(D_R(k)=\operatorname{Soc}E\cong k\) 及正合性，对组成列长度归纳得到长度等式。求值映射已知单射，且两端长度相等，故为同构。

令 \(E_r=(0:_E\mathfrak m^r)=D_R(R/\mathfrak m^r)\)。前述有限长度对偶给出
\(D_R(E_r)\cong R/\mathfrak m^r\)，且元素 \(a+\mathfrak m^r\) 对应 \(e\mapsto ae\)。由 \(E=\bigcup_rE_r\)，一个 \(E\) 上的同态等价于各 \(E_r\) 上相容的同态族。因此
\[
\operatorname{End}_R(E)
=\varprojlim_r\operatorname{Hom}_R(E_r,E)
\cong\varprojlim_rR/\mathfrak m^r.
\]
这些同构把复合变为乘法，故是环同构。

对具体例子，\(R_0=k[t]_{(t)}\) 是 DVR，每个非零理想主生成，而 \(E_0=k(t)/R_0\) 可除。对主理想上同态，除以其生成元就能延拓到 \(R_0\)，Baer 判据因而给出内射性。其基座是 \(kt^{-1}\)，任意非零元素取适当的 \(t\) 幂后都成为这个基座中的非零元素，所以 \(k\subseteq E_0\) 本质，\(E_0\) 是内射包络。\(D_{R_0}(R_0)=E_0\)，故前面已证的自同态计算给出
\[
D_{R_0}D_{R_0}(R_0)=\operatorname{End}_{R_0}(E_0)
\cong\widehat R_0=k[\![t]\!].
\]
求值作用在 \(E_0\) 上就是乘法，因而自然映射正是完备化映射。它不满：级数 \(f(t)=\sum_{j\ge1}t^{j!}\) 不在 \(R_0\) 中。若 \(f=p/q\)，其中 \(p,q\in k[t]\)、\(q(0)\ne0\)，取 \(j\) 足够大，使 \(j!>\deg p\) 且 \(j!-(j-1)!>\deg q\)。在 \(qf=p\) 中比较 \(t^{j!}\) 的系数，左边只有 \(q(0)\ne0\)，右边为零，矛盾。
\end{proof}

完备性在这里有了具体作用：只有 \(R=\widehat R\) 时，\(D_RD_R(R)=D_R(E)\) 才由同一个环 \(R\) 表示。以下的\term[Matlis-duiou-dingli]{Matlis 对偶定理}[Matlis duality theorem]将这一等式推广到两个完整的模类别。

\begin{theorem}[Matlis 对偶]\label{b3:thm:matlis-duality}
设 \((R,\mathfrak m,k)\) 为交换完备 Noether 局部环，\(E=E_R(k)\)，\(D_R(-)=\operatorname{Hom}_R(-,E)\)。则 \(E\) 是 Artin 模；\(D_R\) 把有限生成的 \(R\) 模变成 \(R\) 上的 Artin 模，把 \(R\) 上的 Artin 模变成有限生成的 \(R\) 模。对这两类中的任意模 \(M\)，求值映射 \(\delta_M:M\to D_RD_R(M)\)、\(m\mapsto(\varphi\mapsto\varphi(m))\) 是同构。因此 \(D_R\) 在有限生成模的范畴与 Artin 模的范畴之间给出反变等价。
\end{theorem}
\begin{proof}
证明先将 \(E\) 的子模降链转成 \(R\) 的理想升链，得到 \(E\) 的 Artin 性。有限生成模的对偶随后嵌入有限份 \(E\)，Artin 模本身也将嵌入有限份 \(E\)，但后一项须借助降链条件和元素检测。有限展示给出有限生成一侧的双对偶同构，再用求值映射的相容关系返回 Artin 一侧。

先证 \(E\) 满足降链条件。对子模 \(L\subseteq E\) 记
\[
L^\perp=\{\,a\in R\mid aL=0\,\},
\qquad
J^\perp=(0:_EJ)\quad(J\subseteq R\;\text{为理想}).
\]
这里 \(L^\perp\) 是 \(R\) 中的理想，而 \(J^\perp\) 是 \(E\) 中的子模，两次取垂直才回到原来的对象类型。
显然 \(L\subseteq L^{\perp\perp}\)。若 \(e\notin L\)，元素检测性质在 \(E/L\) 上给出同态 \(f:E/L\rightarrow E\)，使 \(f(e+L)\ne0\)。其与商映射的复合是 \(E\) 的自同态，前一引理和完备性使它为乘以某个 \(a\in R\)。此时 \(aL=0\) 而 \(ae\ne0\)，所以 \(e\notin L^{\perp\perp}\)。故 \(L=L^{\perp\perp}\)。同样，对 \(a\notin J\)，在 \(R/J\) 上检测 \(a+J\) 的同态给出 \(e\in J^\perp\) 且 \(ae\ne0\)，故 \(J=J^{\perp\perp}\)。在当前完备 Noether 局部环的假设下，这两个取垂直的操作互为逆，且反转包含关系：
\[
\begin{tikzcd}[column sep=5em]
\{L\subseteq E:L\;\text{为子模}\}
\arrow[r, shift left=1.2ex, "L\mapsto L^\perp"] &
\{J\subseteq R:J\;\text{为理想}\}
\arrow[l, shift left=1.2ex, "J\mapsto J^\perp"]
\end{tikzcd}
\]
\[
L^{\perp\perp}=L,\qquad J^{\perp\perp}=J,\qquad
L_1\subseteq L_2\ \Longleftrightarrow\ L_2^\perp\subseteq L_1^\perp.
\]
因此 \(E\) 的子模降链变成 \(R\) 的理想升链；后者稳定，取垂直后前者也稳定，故 \(E\) 为 Artin 模。

若 \(M\) 有限生成，由满射 \(R^r\twoheadrightarrow M\) 得单射 \(D_R(M)\hookrightarrow E^r\)，故 \(D_R(M)\) 为 Artin 模。\(R\) 的 Noether 性又给出有限展示 \(R^s\rightarrow R^r\rightarrow M\rightarrow0\)。\(D_RD_R\) 是正合协变函子，\(\delta_R\) 由前一引理是同构，所以将该展示取两次对偶并比较余核，得到 \(\delta_M\) 同构。

若 \(A\) 为 Artin 模，考察全部映射 \(A\rightarrow E^r\) 的核，其中 \(r\) 可为任意非负整数。降链条件保证这些核中有极小者 \(K\)。若 \(0\ne x\in K\)，可再加一个分量 \(A\rightarrow E\) 检测 \(x\)，把核严格缩小，矛盾。因此某个映射给出单射 \(A\hookrightarrow E^r\)。取对偶得到满射 \(R^r\twoheadrightarrow D_R(A)\)，故 \(D_R(A)\) 有限生成。

最后，直接计算求值映射可得
\[
D_R(\delta_A)\circ\delta_{D_R(A)}
=1_{D_R(A)}.
\]
因为 \(D_R(A)\) 有限生成，\(\delta_{D_R(A)}\) 已是同构，故 \(D_R(\delta_A)\) 同构。\(D_R\) 正合且检测非零模，遂使 \(\delta_A\) 的核和余核均为零。求值映射的自然性说明这些同构与一切模同态相容，给出所述反变等价。
\end{proof}

Artin 性与有限长度仍须区分。设 \(k\) 为域、\(A=k[\![t]\!]\)，则 \(E=k(\!(t)\!)/A\) 由上述定理为 Artin 模，满足子模降链条件；但它包含严格上升的子模链 \(t^{-1}A/A\subsetneq t^{-2}A/A\subsetneq\cdots\)，所以不是 Noether 模，更不是有限长度模。有限长度要求升链、降链同时稳定，Matlis 对偶的 Artin 一侧允许这里无界的分母。

\subsection{典范模及 Cohen–Macaulay 情形}
\label{b3:subsec:2202:05}

正则局部环的最高次局部上同调给出典范模构造中的对偶对象。计算时，内射消解中被极大理想的幂零化的部分由基座决定；极小内射消解又把各项的基座直接转为剩余域上的 Ext。

\begin{lemma}\label{b3:lem:injective-maximal-torsion}
设 \((A,\mathfrak n,k)\) 为交换 Noether 局部环，\(E=E_A(k)\)，\(J\) 为内射 \(A\) 模。若 \(\{u_\lambda\}_{\lambda\in\Lambda}\) 为 \(\operatorname{Soc}_A J\) 的一个 \(k\) 基，则
\[
\Gamma_{\mathfrak n}(J)\cong E^{(\Lambda)}.
\]
其中 \(E^{(\Lambda)}\) 为直和，指标数等于该基座的维数。
\end{lemma}
\begin{proof}
先注意，在 Noether 环上内射模的任意直和仍内射。由 Baer 判据，只需把理想 \(L\) 上的映射延拓到 \(A\)。因为 \(L\) 有限生成，其像落在有限个直和项中；逐分量延拓后再合成即可。

把每个映射 \(k\rightarrow J\)、\(1\mapsto u_\lambda\) 延拓到 \(E\rightarrow J\)，并合成为 \(v:E^{(\Lambda)}\rightarrow J\)。直和的每个元素由有限个分量组成，且由\cref{b3:prop:injective-hull-local-properties}被某个 \(\mathfrak n\) 的幂零化。因此它的任何非零子模都含有被 \(\mathfrak n\) 零化的非零元素：取零化幂数最小的元素，再乘 \(\mathfrak n\) 的适当幂。故 \(k^{(\Lambda)}\) 在该直和中本质。\(v\) 在这一个基座上单射，遂整体单射。

其像内射，故 \(J=v(E^{(\Lambda)})\oplus J'\)。\(J'\) 的基座为零，因为 \(u_\lambda\) 已张成 \(J\) 的全部基座。若 \(\Gamma_{\mathfrak n}(J')\ne0\)，同样取最小零化幂数就会找到非零基座元素，矛盾。因此极大理想幂挠部分恰为第一直和项。
\end{proof}

\begin{lemma}\label{b3:lem:minimal-injective-resolution-socle}
设 \((A,\mathfrak n,k)\) 为交换 Noether 局部环，\(M\) 为 \(A\) 模。令 \(C^0=M\)，对 \(q\ge0\) 递归取内射包络 \(C^q\subseteq J^q=E_A(C^q)\)、商模 \(C^{q+1}=J^q/C^q\)。以商映射与下一次包络的包含复合为微分，所得 \(M\to J^\bullet\) 是 \(M\) 的极小内射消解，且
\[
\operatorname{Hom}_A(k,J^\bullet)\;\text{的微分均为零},
\qquad
\operatorname{Soc}_A J^q\cong\operatorname{Ext}_A^q(k,M).
\]
若某个整数 \(n\ge0\) 满足 \(\operatorname{Ext}_A^{n+1}(A/L,M)=0\) 对所有理想 \(L\subseteq A\) 成立，则 \(C^n\) 内射，可在第 \(n\) 项结束此消解。
\end{lemma}
\begin{proof}
第十七章的极小自由消解张量剩余域后微分为零。极小内射消解相应地给出零映射
\[
\operatorname{Hom}_A(k,J^q)\xrightarrow{0}
\operatorname{Hom}_A(k,J^{q+1})\qquad(q\ge0).
\]
前者逐层选择极小生成组，后者逐层取本质的内射包络，使每个简单子模都留在当前微分的核中。

由构造，\(J^q\to J^{q+1}\) 的核是 \(C^q\)，在 \(q\ge1\) 时又等于前一微分的像，所以增广列正合。若 \(u\in\operatorname{Soc}_A J^q\) 非零，则 \(Au\cong k\) 是简单子模。本质性使 \(Au\cap C^q\ne0\)，故 \(Au\subseteq C^q\)。因此全部基座元素都在微分的核中，取 \(\operatorname{Hom}_A(k,-)\) 后微分为零。

内射消解计算 Ext 与自由消解计算 Ext 给出同一个模。具体地，对自由消解 \(P_\bullet\to k\)，取双复形 \(\operatorname{Hom}_A(P_p,J^q)\)，其中 \(p,q\ge0\)、总次数为 \(p+q\)。沿 \(q\) 方向，\(P_p\) 的投射性使 Hom 保持消解正合，只留下 \(q=0\) 行的 \(\operatorname{Hom}_A(P_\bullet,M)\)；沿 \(p\) 方向，\(J^q\) 的内射性使正次数消失，只留下 \(p=0\) 列的 \(\operatorname{Hom}_A(k,J^\bullet)\)。两个消解虽可无限延伸，第一象限中的每条固定总次数对角线仍只有有限项，逐次提升圈并消去正合行、列中的分量便得到两种计算之间的同构。后一个复形微分为零，故其第 \(q\) 次上同调就是 \(\operatorname{Hom}_A(k,J^q)=\operatorname{Soc}_A J^q\)。

最后，短正合列 \(0\to C^q\to J^q\to C^{q+1}\to0\) 逐次作维数平移，给出
\[
\operatorname{Ext}_A^1(A/L,C^n)
\cong\operatorname{Ext}_A^{n+1}(A/L,M)=0.
\]
Baer 判据使 \(C^n\) 内射。它既内射又在 \(J^n\) 中本质，便只能是 \(J^n\) 自身，于是 \(C^{n+1}=0\)。当 \(n=0\) 时，维数平移等式不需经过中间项，结论仍成立。
\end{proof}

\begin{proposition}\label{b3:prop:regular-local-top-cohomology}
设 \((S,\mathfrak n,k)\) 为 \(n\) 维交换 Noether 正则局部环。则
\[
H^q_{\mathfrak n}(S)=0\quad(q\ne n),
\qquad
H^n_{\mathfrak n}(S)\cong E_S(k).
\]
此处不要求 \(S\) 完备，最后的同构依赖于一个剩余域基座的识别。
\end{proposition}
\begin{proof}
先用正则参数的 Koszul 对偶计算 \(\operatorname{Ext}_S^q(k,S)\)，再由极小内射消解把这些 Ext 识别为各项基座。内射模的极大理想幂挠部分由基座决定，因而最终只在最高次数留下剩余域的内射包络。

取 \(\mathfrak n\) 的正则生成系 \(x_1,\ldots,x_n\)，其存在由\cref{b3:thm:regular-local-characterizations}保证。\cref{b3:thm:koszul-regular-resolution}给出 \(k\) 的 Koszul 自由消解。外积配对
\(\bigwedge^qS^n\times\bigwedge^{n-q}S^n\rightarrow\bigwedge^nS^n\cong S\)
把其对偶复形识别为反向的 Koszul 复形。具体地，令 \(\Phi_q(v)(u)\) 为 \(u\wedge v\) 的最高外积系数。对 \(u\in K_{q+1}\)、\(v\in K_{n-q}\)，等式 \(\mathrm{d}(u\wedge v)=0\) 给
\(\delta\Phi_q(v)=(-1)^q\Phi_{q+1}(\mathrm{d}v)\)。故以 \((-1)^{q(q-1)/2}\Phi_q\) 作第 \(q\) 项的识别，所有微分均相容。因此
\[
\operatorname{Ext}_S^q(k,S)=
\begin{cases}k,&q=n,\\0,&q\ne n.\end{cases}
\]
等价地，也可将 \(n\) 个两项复形 \(S\xrightarrow{x_i}S\) 逐个对偶，每对偶一个因子便把它的非零余核提高一个次数。

取\cref{b3:lem:minimal-injective-resolution-socle}构造的 \(S\) 的极小内射消解 \(J^\bullet\)。由该引理，\(\operatorname{Soc}_S J^q\) 的维数在 \(q=n\) 时为一，在其余次数为零。又由\cref{b3:thm:regular-local-global-dimension}，对所有理想 \(L\subseteq S\) 都有 \(\operatorname{Ext}_S^{n+1}(S/L,S)=0\)，所以同一引理保证消解在第 \(n\) 项结束。

\cref{b3:lem:injective-maximal-torsion}于是给出
\[
\Gamma_{\mathfrak n}(J^q)\cong
\begin{cases}E_S(k),&q=n,\\0,&q\ne n.\end{cases}
\]
这个复形只有一个非零项。按支集挠子函子的导出函子计算局部上同调，即得结论；这也包括 \(n=0\)、\(S=k\) 的情形。
\end{proof}

\begin{theorem}\label{b3:thm:regular-local-duality}
设 \((S,\mathfrak n,k)\) 为 \(n\) 维交换完备 Noether 正则局部环，\(E=E_S(k)\)，\(D_S(-)=\operatorname{Hom}_S(-,E)\)。选定同构 \(H^n_{\mathfrak n}(S)\cong E\) 后，对每个有限生成的 \(S\) 模 \(N\) 和每个整数 \(i\)，有关于 \(N\) 自然的同构
\[
H^i_{\mathfrak n}(N)
\cong D_S\bigl(\operatorname{Ext}_S^{n-i}(N,S)\bigr),
\qquad
D_S\bigl(H^i_{\mathfrak n}(N)\bigr)
\cong\operatorname{Ext}_S^{n-i}(N,S).
\]
负次局部上同调和负次 Ext 均约定为零。各 \(H^i_{\mathfrak n}(N)\) 均为 Artin 模。
\end{theorem}
\begin{proof}
要比较局部上同调与环境 Ext，可以用同一个双复形的两种计算。自由消解方向先计算，会返回 \(N\) 的 Čech 上同调；Čech 方向先计算，则由前一命题只留下最高次数的内射包络，再用有限秩自由性把所得复形识别为 Ext 计算的 Matlis 对偶。

令 \(C^\bullet\) 为正则生成系 \(x_1,\ldots,x_n\) 的 Čech 复形。\cref{b3:thm:regular-local-homological-criterion}保证 \(N\) 有长度至多 \(n\)、各项有限秩自由的消解 \(P_\bullet\rightarrow N\)。将 \(P_j\) 放在上链次数 \(-j\)，取双复形
\[
C^q\otimes_SP_j,\qquad 0\le q\le n,\quad 0\le j\le n,
\]
其总次数为 \(q-j\)。总微分采用
\(\mathrm{d}(c\otimes p)=\mathrm{d}_Cc\otimes p+(-1)^{\deg c}c\otimes \mathrm{d}_Pp\)，故平方为零。

每个 \(C^q\) 都平坦，所以沿 \(P\) 方向取同调只在 \(j=0\) 留下 \(C^q\otimes_SN\)。因此总上同调为 \(H^i_{\mathfrak n}(N)\)，这里使用\cref{b3:thm:cech-local-cohomology}。换一个方向，前一命题给 \(H^q(C^\bullet)=0\) 对 \(q\ne n\)，而 \(H^n(C^\bullet)=E\)。各 \(P_j\) 自由，故这次只留下 \(q=n\) 行的 \(E\otimes_SP_j\)。该行的微分为 \((-1)^n\mathrm{d}_P\)，在第 \(j\) 项乘以 \((-1)^{nj}\) 就将其识别为通常的链复形 \(E\otimes_SP_\bullet\)。两次比较用有限滤过完成：按一方向次数过滤总复形，把圈从最外侧的分量开始处理，利用正合性减去边界，有限次后只剩非零同调所在的行；边界的比较同样进行。两个方向均有界，故过程终止，并与复形映射相容。

最高行位于 Čech 次数 \(n\)，自由消解的链次数 \(j\) 将总次数降到 \(n-j\)，所以在总次数 \(i\) 必有 \(j=n-i\)。这给出
\[
H^i_{\mathfrak n}(N)
\cong H_{n-i}(E\otimes_SP_\bullet).
\]
有限自由性给出逐项同构
\[
E\otimes_SP_j
\cong\operatorname{Hom}_S\bigl(\operatorname{Hom}_S(P_j,S),E\bigr),
\quad e\otimes p\longmapsto[\lambda\mapsto\lambda(p)\cdot e].
\]
此同构与微分相容。\(D_S\) 正合，故右侧复形的同调为 \(D_S(\operatorname{Ext}_S^{n-i}(N,S))\)，得到第一个同构。Ext 模有限生成，Matlis 对偶定理使其对偶为 Artin 模；再取一次对偶并使用有限生成模的求值同构，得到第二个同构。
\end{proof}

现在可以准确地命名所需的有限生成模。

\begin{definition}\label{b3:def:canonical-module-regular-quotient}
设 \((R,\mathfrak m,k)\) 为 \(d\) 维交换完备 Noether 局部环，\(D_R=\operatorname{Hom}_R(-,E_R(k))\)。若有限生成的 \(R\) 模 \(\omega_R\) 满足
\[
D_R(\omega_R)\cong H^d_{\mathfrak m}(R),
\]
则称 \(\omega_R\) 为 \(R\) 的\term[dianfanmo]{典范模}[canonical module]。
\end{definition}
\symidx{\(\omega_R\)：局部环 \(R\) 的典范模}

Matlis 对偶定理说明，一旦典范模存在，对定义中的同构再取对偶便得
\[
\omega_R\cong D_R\bigl(H^d_{\mathfrak m}(R)\bigr).
\]
这从 \(R\) 自身的最高次局部上同调确定典范模的同构类；下面的环境 Ext 则给出它的一个具体实现。这里不指定某一个同构，所以“典范”不表示所有实现之间有唯一的同构，与内射包络的唯一性一样，它确定的是同构类。

\begin{proposition}\label{b3:prop:canonical-module-quotient}
设 \((S,\mathfrak n,k)\) 为 \(n\) 维交换完备 Noether 正则局部环，\(I\subsetneq S\) 为理想，\(R=S/I\)、\(\mathfrak m=\mathfrak n/I\)、\(d=\dim R\)。则
\[
\omega_R=\operatorname{Ext}_S^{n-d}(R,S)
\]
是 \(R\) 的典范模。若 \(R\) 还是 Cohen–Macaulay 环，令 \(c=n-d\)，则
\[
\operatorname{Ext}_S^q(R,S)=0\quad(q\ne c),
\qquad
\operatorname{depth}_R\omega_R=d.
\]
从不同的完备正则局部环到 \(R\) 的满射表示产生同构的典范模。
\end{proposition}
\begin{proof}
先比较两个环上的对偶对象。模
\[
E_R=\operatorname{Hom}_S(R,E_S(k))=(0:_{E_S(k)}I)
\]
作为 \(R\) 模内射，因为对每个 \(R\) 模 \(M\)，在 \(1\) 处求值给出自然同构
\[
\operatorname{Hom}_R(M,E_R)\cong\operatorname{Hom}_S(M,E_S(k)).
\]
右端是正合函子。\(E_R\) 中的 \(k\) 本质，因为它已经在 \(E_S(k)\) 中本质。因此 \(E_R\) 就是 \(E_R(k)\) 的一个实现，而 \(D_R(M)=D_S(M)\)。由同一组生成元的 Čech 复形，\(H^i_{\mathfrak m}(M)=H^i_{\mathfrak n}(M)\)。正则局部对偶应用于 \(M=R\)，便给出典范模定义要求的同构；\cref{b3:prop:finite-ext-localization,b3:lem:ext-basic-properties}分别说明 Ext 有限生成且被第一变量的零化子 \(I\) 零化，故是有限生成的 \(R\) 模。

现假设 \(R\) 为 CM 环。\cref{b3:thm:depth-local-cohomology}与\cref{b3:thm:local-cohomology-dimension-vanishing}使 \(H^i_{\mathfrak m}(R)\) 仅在 \(i=d\) 非零。正则局部对偶立即给出 Ext 的集中次数。

进一步，\(R\) 作为 \(S\) 模的深度也是 \(d\)，因为 \(\mathfrak n\) 中元素在 \(R\) 上的作用只由其在 \(\mathfrak m\) 中的像决定。由\cref{b3:thm:auslander-buchsbaum}，\(\operatorname{pd}_SR=c\)。取长度 \(c\)、各项有限秩自由的消解 \(F_\bullet\rightarrow R\)。对偶复形 \(\operatorname{Hom}_S(F_\bullet,S)\) 仅在最后一项有上同调 \(\omega_R\)；反向排列后，它就是 \(\omega_R\) 的长度 \(c\) 的自由消解。再对偶一次，逐项双对偶回到 \(F_\bullet\)，所以
\[
\operatorname{Ext}_S^q(\omega_R,S)=
\begin{cases}R,&q=c,\\0,&q\ne c.\end{cases}
\]
故 \(\operatorname{pd}_S\omega_R=c\)，再次应用 Auslander–Buchsbaum 公式得深度 \(n-c=d\)。\(S\) 与 \(R\) 上的深度相同，得到所述结论。不同满射表示所构造的典范模具有相同的同构类型，因为它们都同构于 \(D_R(H^d_{\mathfrak m}(R))\)。
\end{proof}

存在性一项本身没有要求 CM 性：最高次局部上同调对应的环境 Ext 总能给出典范模。CM 性所增加的是其他 Ext 次数全部消失，以及典范模的深度达到 \(d\)。环境对偶信息因此集中在一个次数，下一小节才能把全部局部对偶写成只含 \(\omega_R\) 的公式。

\subsection{局部上同调与 Ext 的 Matlis 对偶}
\label{b3:subsec:2202:06}

为了从环境环 \(S\) 转到商环 \(R\)，需要把两个 Ext 计算联系起来。若 \(R\) 相对于 \(S\) 的对偶信息集中在次数 \(c\)，先通过 \(R\) 计算任意 \(R\) 模的环境 Ext，就只产生整体的次数平移 \(c\)。下面的引理将把这项平移写成两个 Ext 的自然同构。

\begin{lemma}\label{b3:lem:canonical-change-of-rings}
设 \(S\) 为交换 Noether 环，\(I\subseteq S\) 为理想，\(R=S/I\)。假设 \(S\) 作为自身上的模有有界内射消解，且存在整数 \(c\ge0\) 使
\(\operatorname{Ext}_S^q(R,S)=0\) 对 \(q\ne c\) 成立。记 \(W=\operatorname{Ext}_S^c(R,S)\)。则对每个有限生成的 \(R\) 模 \(M\) 及每个 \(p\ge0\)，有自然同构
\[
\operatorname{Ext}_R^p(M,W)
\cong\operatorname{Ext}_S^{p+c}(M,S),
\]
且 \(\operatorname{Ext}_S^q(M,S)=0\) 对 \(q<c\) 成立。
\end{lemma}
\begin{proof}
取内射 \(S\) 消解 \(J^\bullet\)，选整数 \(h\ge0\) 使 \(J^q=0\) 对 \(q>h\) 成立，再取每项有限秩自由的 \(R\) 消解 \(P_\bullet\rightarrow M\)。后一个消解可以无限延伸。令
\[
T^q=\operatorname{Hom}_S(R,J^q),
\qquad B^{p,q}=\operatorname{Hom}_R(P_p,T^q).
\]
每个 \(T^q\) 都是内射 \(R\) 模，理由是
\(\operatorname{Hom}_R(-,T^q)\cong\operatorname{Hom}_S(-,J^q)\)
正合。\(B\) 的指标范围为 \(p\ge0\)、\(0\le q\le h\)，总次数为 \(p+q\)。在 \((p,q)\) 项上，总微分取来自 \(P\) 的 Hom 微分加 \((-1)^p\) 倍来自 \(T\) 的微分，使两方向反交换。第一象限与 \(q\) 方向有界保证每条总次数对角线有限。

两种计算的目标分别是原来的环境 Ext 和经过 \(R\) 的 Ext。先沿 \(p\) 方向取上同调，\(T^q\) 的内射性使 \(p>0\) 时消失，\(p=0\) 为 \(\operatorname{Hom}_R(M,T^q)\cong\operatorname{Hom}_S(M,J^q)\)，所以总上同调是 \(\operatorname{Ext}_S^*(M,S)\)。再沿 \(q\) 方向取上同调，\(P_p\) 自由使 Hom 保持核、像及其商，故
\[
H^q(B^{p,\bullet})
=\operatorname{Hom}_R(P_p,H^q(T^\bullet)),
\]
它只在 \(q=c\) 留下 \(\operatorname{Hom}_R(P_p,W)\)。因此总次数 \(p+c\) 的上同调就是 \(\operatorname{Ext}_R^p(M,W)\)；低于 \(c\) 的总次数没有这一行的项，上同调为零。这里的单行比较使用次数滤过：逐行提升圈，利用 \(q\ne c\) 时的正合性消去相应分量，剩余行上的边界恰为总边界的像。每个固定总次数及相邻次数都只有有限项，因此圈与边界的比较只经过有限步。此构造与同态诱导的消解映射相容，给出自然性。
\end{proof}

\begin{theorem}[局部对偶]\label{b3:thm:canonical-module-local-duality}
设 \((S,\mathfrak n,k)\) 为 \(n\) 维交换完备 Noether 正则局部环，\(I\subsetneq S\) 为理想，且 \((R,\mathfrak m,k)=(S/I,\mathfrak n/I,k)\) 为 \(d\) 维 Cohen–Macaulay 环。取
\[
\omega_R=\operatorname{Ext}_S^{n-d}(R,S),
\qquad D_R=\operatorname{Hom}_R(-,E_R(k)).
\]
则对每个有限生成的 \(R\) 模 \(M\) 和每个整数 \(i\)，有关于 \(M\) 自然的同构
\begin{equation}\label{b3:eq:canonical-local-duality}
\begin{aligned}
H^i_{\mathfrak m}(M)
&\cong D_R\bigl(\operatorname{Ext}_R^{d-i}(M,\omega_R)\bigr),\\
D_R\bigl(H^i_{\mathfrak m}(M)\bigr)
&\cong\operatorname{Ext}_R^{d-i}(M,\omega_R).
\end{aligned}
\end{equation}
负次局部上同调和负次 Ext 均取零。所有 \(H^i_{\mathfrak m}(M)\) 都是 Artin 模，并且
\[
\operatorname{Ext}_R^p(M,\omega_R)=0\qquad(p>d).
\]
\end{theorem}
\begin{proof}
令 \(c=n-d\)。\cref{b3:prop:canonical-module-quotient}给出 \(\operatorname{Ext}_S^q(R,S)\) 只在 \(q=c\) 非零。\(S\) 的有界内射消解已在\cref{b3:prop:regular-local-top-cohomology}的证明中构造。故前一引理给出
\[
\operatorname{Ext}_S^{n-i}(M,S)
\cong\operatorname{Ext}_R^{d-i}(M,\omega_R)
\]
在 \(d-i\ge0\) 时成立；若 \(d-i<0\)，左端也由引理的低次消失为零。再用\cref{b3:thm:regular-local-duality}，以及已验证的 \(D_S(M)=D_R(M)\)、\(H^i_{\mathfrak n}(M)=H^i_{\mathfrak m}(M)\)，得到两行公式。Artin 性来自正则局部对偶；当 \(i<0\) 时，局部上同调为零，第二行给出最后的 Ext 消失。
\end{proof}

上式称为\term[jubu-duiou]{局部对偶}[local duality]。次数变换 \(q=d-i\) 中的 \(d\) 是底环 \(R\) 的维数；模自身的维数则限制哪些局部上同调次数可能出现，两者不能互换。公式允许用有限生成模的消解研究通常不有限生成的局部上同调，但没有把后者断言为有限生成模。例如在 \(k[\![t]\!]\) 上，\(\omega_R\cong R\)，公式的最高次仍然是 \(H^1_{(t)}(R)\cong E_R(k)\)，与节首计算一致。

\ProseParagraphBreak
设 \(R=S/I\) 满足上述局部对偶定理的假设，即 \(S\) 是完备 Noether 正则局部环，\(R\) 是 \(d\) 维 Cohen–Macaulay 商环，\(\mathfrak m\) 是 \(R\) 的极大理想，\(\omega_R\) 是所得典范模。对非零有限生成 \(R\) 模 \(M\)，令 \(\alpha=\operatorname{depth}_R M\)、\(t=\dim_R M\)。Matlis 对偶检测非零模，因此次数变换 \(q=d-i\) 同时保留消失与非消失：
\[
H^i_{\mathfrak m}(M)\ne0
\quad\Longleftrightarrow\quad
\operatorname{Ext}_R^{d-i}(M,\omega_R)\ne0.
\]
\cref{b3:tab:local-duality-dictionary}把深度、维数与 CM 条件的表述集中在一起。其中 \(M\) 为 CM 模并不要求 \(t=d\)。

\begin{table}[htbp]
\centering
\caption[局部对偶中的消失区间]{局部对偶中的消失区间。\(R\) 为上述完备 CM 商环，\(M\ne0\) 有限生成，\(\alpha=\operatorname{depth}_R M\)、\(t=\dim_R M\)。}
\label{b3:tab:local-duality-dictionary}
\begin{tabularx}{\linewidth}{@{}p{0.15\linewidth}XX@{}}
\toprule
数学信息 & 局部上同调 \(H^i_{\mathfrak m}(M)\) & \(\operatorname{Ext}_R^q(M,\omega_R)\) \\
\midrule
深度 \(\alpha\) & \(i<\alpha\) 时为零，\(i=\alpha\) 时非零 & \(q>d-\alpha\) 时为零，\(q=d-\alpha\) 时非零 \\
维数 \(t\) & \(i>t\) 时为零 & \(q<d-t\) 时为零 \\
\(M\) 为 CM 模 & 只在 \(i=t\) 非零 & 只在 \(q=d-t\) 非零 \\
\bottomrule
\end{tabularx}
\end{table}

可能非零的局部上同调次数满足 \(\alpha\le i\le t\)，次数变换将其翻转为 \(d-t\le q\le d-\alpha\)。深度给出 Ext 的最高非零次数 \(d-\alpha\)；模维数给出最低可能次数 \(d-t\)，即其下方全部消失。在 CM 情形 \(\alpha=t\)，整个区间缩成一个确实非零的次数。

\ProseParagraphBreak
先看模维数与底环维数确实不同的情形。取域 \(k\)，令 \(S=R=k[\![x,y]\!]\)、\(\mathfrak m=(x,y)\)、\(M=R/(x)\)。此时环境环维数 \(n\) 与底环维数 \(d\) 均为二，而 \(t=\dim_R M=1\)、\(\alpha=\operatorname{depth}_R M=1\)；\(y\) 在 \(M\cong k[\![y]\!]\) 上是非零因子。\(R\) 的典范模为 \(R\)，自由消解
\[
0\longrightarrow R\xrightarrow{x}R\longrightarrow M\longrightarrow0
\]
对偶后给出 \(\operatorname{Hom}_R(M,R)=0\)、\(\operatorname{Ext}^1_R(M,R)=R/(x)\)，且更高 Ext 为零。因此局部对偶给出
\[
D_R\bigl(H^1_{\mathfrak m}(M)\bigr)
\cong\operatorname{Ext}^{2-1}_R(M,R)
\cong R/(x).
\]
翻转次数中的二来自底环维数 \(d\)，不是模维数 \(t\)；模维数一则决定这里的局部上同调只可能在一次非零。

\ProseParagraphBreak
最后用\subsecref{b3:subsec:2201:03}的例子检验 CM 假设的作用。取域 \(k\) 及
\[
S=k[\![x,y]\!],\qquad R=S/(x^2,xy),\qquad\mathfrak m=(\bar x,\bar y).
\]
此前已算得
\[
\dim R=1,
\qquad H^0_{\mathfrak m}(R)=k\bar x\ne0.
\]
这个环还满足 \(H^1_{\mathfrak m}(R)=k(\!(y)\!)/k[\![y]\!]\ne0\)。环境环上的正则局部对偶分别给出
\[
\begin{aligned}
D_R(H^1_{\mathfrak m}(R))&\cong\operatorname{Ext}_S^1(R,S),\\
D_R(H^0_{\mathfrak m}(R))&\cong\operatorname{Ext}_S^2(R,S).
\end{aligned}
\]
两项均非零。\cref{b3:prop:canonical-module-quotient}仍给出典范模 \(\omega_R=\operatorname{Ext}_S^1(R,S)\)，但它只保留与最高次局部上同调对应的一项。若把局部对偶的单模公式用于 \(M=R\)、\(i=0\)，就会得到
\[
0\ne D_R\bigl(H^0_{\mathfrak m}(R)\bigr)
\cong\operatorname{Ext}_R^1(R,\omega_R)=0,
\]
矛盾。真正保留这部分信息的是另一个非零模 \(\operatorname{Ext}_S^2(R,S)\cong D_R(H^0_{\mathfrak m}(R))\)。因此一般情形须保留具有多个上同调的对偶复形，不能只取最高次局部上同调对应的那一个模。

\subsection{Cohen 结构定理与正则局部环的满射表示}
\label{b3:subsec:2202:07}

本节所有存在性及局部对偶证明都从给定满射 \(S\twoheadrightarrow R\) 开始。对显式幂级数商环，这份数据已经在题目中；例如 \(k[\![x_1,\ldots,x_n]\!]/I\) 的环境环就是所需的完备正则局部环。混合特征时也可以从给定完备离散赋值环 \(V\) 上的幂级数环 \(V[\![x_1,\ldots,x_r]\!]\) 开始，而不把系数环误写成域。

\ProseParagraphBreak
对一个未给出满射表示的交换完备 Noether 局部环，附录\cref{b3:thm:coefficient-field-existence,b3:thm:equicharacteristic-power-series-cover,b3:thm:mixed-characteristic-coefficient-ring}给出所需的环境环与满射。等特征时先取系数域，再按极大理想阶数展开，得到幂级数环满射；混合特征时以 Cohen 环替代系数域，同样逐阶选择剩余系数并取完备极限，得到其上幂级数环的满射。两种环境环均完备正则局部。相应结构定理还可对照松村英之\cite[\S 29, Theorems 29.3--29.4, pp. 225--226]{Matsumura1989CommutativeRingTheory}。

\ProseParagraphBreak
因此每个交换完备 Noether 局部环都有本节意义下的典范模；若该环还为 CM 环，就可使用\eqref{b3:eq:canonical-local-duality}。对于非完备局部环，这一构造首先给出的是完备化上的典范模；能否下降为原环上的有限生成模，是另一个需要条件的问题。

% !TeX root = ../../Book3.tex
% 纲要标题：### 22.3 Gorenstein 环
% 纲要位置：计划纲要.md，第 2204 行。

\section{Gorenstein 环}
\label{b3:sec:2203}

环自身的有限内射消解，比所有模具有有限自由消解的要求宽松。一般 Bass 判别把这种有限性转为 Cohen–Macaulay 性与顶次 Ext 的一维条件，完全交则可由 Koszul 对偶满足这一条件。


% 纲要内容：
%
% * Noether 局部环的有限自内射维数
% * 与 Cohen–Macaulay 性的联系
% * 完全交蕴含 Gorenstein 的适当版本
%

\subsection{Noether 局部环的有限自内射维数}
\label{b3:subsec:2203:01}

正则局部环用有限自由消解控制所有模；Gorenstein 性则问环本身能否有有限的内射消解。两种要求的方向不同。特别是，剩余域有无限自由消解，并不排除环作为自身的模内射。

\begin{definition}\label{b3:def:injective-dimension-gorenstein}
设 \(R\) 为交换环，\(N\ne0\) 为 \(R\) 模。若存在正合列
\[
0\longrightarrow N\longrightarrow E^0\longrightarrow\cdots
\longrightarrow E^s\longrightarrow0
\]
且每个 \(E^i\) 内射，则满足此条件的最小 \(s\) 称为 \(N\) 的\term[neisheweishu]{内射维数}[injective dimension]，记为 \(\operatorname{id}_RN\)；若不存在有限长度的此种消解，则记为 \(\infty\)。对交换 Noether 局部环 \(R\)，条件 \(\operatorname{id}_RR<\infty\) 成立时，称 \(R\) 为 \term[Gorensteinhuan]{Gorenstein 环}[Gorenstein ring]。交换 Noether 环在每个素理想处的局部环均有此性质时，也称为 Gorenstein 环。
\end{definition}
\symidx{\(\operatorname{id}_RN\)：内射维数}

内射消解中的模通常不是有限模。譬如离散赋值环 \(R\) 的分式域 \(K\) 给出
\[
0\longrightarrow R\longrightarrow K\longrightarrow K/R\longrightarrow0.
\]
两个右端模都可除，由主理想整环上的 Baer 判据都内射；而 \(R\) 不是内射，因为均匀参数的倒数不在 \(R\) 中。因此 \(\operatorname{id}_RR=1\)。这里仍是用一个有限长度的消解描述一个一维环，却不能要求消解各项有限生成。

\begin{lemma}\label{b3:lem:injective-dimension-test}
设 \(R\) 为交换环，\(N\ne0\) 为 \(R\) 模，\(s\ge0\)。下列两条件等价：
\[
\operatorname{id}_RN\le s;\qquad
\operatorname{Ext}_R^{s+1}(R/J,N)=0
\text{ 对每个理想 }J\subseteq R\text{ 成立}.
\]
\end{lemma}
\begin{proof}
第一条件推出更高次的 Ext 全部为零。反过来，在一个内射消解中截取
\(0\to N\to E^0\to\cdots\to E^{s-1}\to C\to0\)；\(s=0\) 时令 \(C=N\)。逐次使用第二变量长正合列，得到
\(\operatorname{Ext}_R^1(R/J,C)\cong\operatorname{Ext}_R^{s+1}(R/J,N)=0\)。
由 \(0\to J\to R\to R/J\to0\) 的 Hom 长正合列，每个 \(J\to C\) 均可延拓到 \(R\)。第二卷第7.2节的 Baer 判据使 \(C\) 内射，故上述截断就是长度至多 \(s\) 的内射消解。

这里使用的内射计算与此前用自由消解定义的 Ext 一致：把自由消解 \(P_\bullet\to R/J\) 和内射消解 \(N\to E^\bullet\) 组成 \(\operatorname{Hom}_R(P_p,E^q)\)，两个方向分别由自由性和内射性保持正合，因而分别只剩 \(\operatorname{Hom}_R(P_\bullet,N)\) 与 \(\operatorname{Hom}_R(R/J,E^\bullet)\)。给定总次数中只有有限个项，双复形的圈消去论证给出相同上同调。
\end{proof}

\subsection{Gorenstein 性与 Cohen–Macaulay 性}
\label{b3:subsec:2203:02}

有限自内射维数不仅给出一项同调界，还强迫环达到 Cohen--Macaulay 深度。下面的判别不要求环完备，也不要求事先给定正则局部覆盖。证明采用素点间的 Ext 比较与正则元素换环，可对照松村英之\cite[\S 18, Lemmas 1--4 and Theorem 18.1, pp. 139--144]{Matsumura1989CommutativeRingTheory}。

\begin{theorem}[Gorenstein 环的判别]\label{b3:thm:gorenstein-bass-criterion}
设 \((R,\mathfrak m,k)\) 为交换 Noether 局部环，\(d=\dim R\)。下列条件等价：
\begin{equivalences}
\item \(R\) 为 Gorenstein 环；
\item \(R\) 为 Cohen--Macaulay 环，且
\(\dim_k\operatorname{Ext}_R^d(k,R)=1\)；
\item \(\operatorname{Ext}_R^i(k,R)=0\) 对 \(i\ne d\) 成立，且
\(\operatorname{Ext}_R^d(k,R)\cong k\)。
\end{equivalences}
这些条件成立时，\(\operatorname{id}_RR=d\)。
\end{theorem}
\begin{proof}
先说明三个计算事实。第一，\cref{b3:lem:injective-dimension-test}中的理想可只取素理想：每个有限模有有限素循环滤过，长正合列把素循环商上的 Ext 消失传到任意有限模，特别是各个 \(R/J\)。

第二，若 \(x\in\mathfrak m\) 是 \(R\) 非零因子，令 \(\overline R=R/xR\)。对每个 \(\overline R\) 模 \(N\)，有
\[
\operatorname{Hom}_R(N,R)=0,\qquad
\operatorname{Ext}_R^{j+1}(N,R)
\cong\operatorname{Ext}_{\overline R}^{j}(N,\overline R)
\quad(j\ge0).
\]
确实，二项自由消解 \(0\to R\xrightarrow{x}R\to\overline R\to0\) 给出
\(\operatorname{Ext}_R^q(\overline R,R)\) 只在 \(q=1\) 为 \(\overline R\)。把 \(N\) 的自由 \(\overline R\) 消解与 \(R\) 的内射消解组成 Hom 双复形，按\cref{b3:lem:canonical-change-of-rings}证明中的单行消去，便得上述换环同构；每个总次数只有有限项，不涉及无限乘积。

第三，对有限 \(R\) 模 \(M\)、素理想 \(\mathfrak p\) 及 \(t=\dim R/\mathfrak p\)，有
\[
\operatorname{Ext}_R^{j+t}(k,M)=0
\ \Longrightarrow\
\operatorname{Ext}_{R_{\mathfrak p}}^j
  (\kappa(\mathfrak p),M_{\mathfrak p})=0.
\]
先设 \(t=1\)，取 \(a\in\mathfrak m\setminus\mathfrak p\)。模 \(R/(\mathfrak p,a)\) 有有限长度，故给定的 \((j+1)\) 次消失也对它成立。由
\(0\to R/\mathfrak p\xrightarrow{a}R/\mathfrak p\to R/(\mathfrak p,a)\to0\)，乘以 \(a\) 在有限模 \(\operatorname{Ext}_R^j(R/\mathfrak p,M)\) 上满射；Nakayama 引理使该模为零，局部化即得断言。一般 \(t\) 时，取长度为 \(t\) 的最长素理想链，从 \(\mathfrak m\) 到 \(\mathfrak p\) 逐个作这一步；每个相邻商的局部维数为一，Ext 与有限模的局部化相容。

现在设 \(r=\operatorname{id}_RR<\infty\)。取满足 \(\dim R/\mathfrak p=d\) 的极小素理想，则零维局部环 \(R_{\mathfrak p}\) 有非零基座。第三个事实的逆否命题给出 \(\operatorname{Ext}_R^d(k,R)\ne0\)，所以 \(r\ge d\)。若 \(r>0\)，则 \(\operatorname{Ext}_R^r(k,R)\ne0\)：第一事实给出某个 \(\mathfrak p\) 使 \(\operatorname{Ext}_R^r(R/\mathfrak p,R)\ne0\)；若 \(\mathfrak p\ne\mathfrak m\)，取 \(a\in\mathfrak m\setminus\mathfrak p\)，同一短正合列及 \((r+1)\) 次 Ext 消失会使乘以 \(a\) 满射，再由 Nakayama 引理得到矛盾。

此时 \(\mathfrak m\notin\operatorname{Ass}R\)。否则有单射 \(k\to R\)，其余核的 \((r+1)\) 次 Ext 为零，迫使
\(0=\operatorname{Ext}_R^r(R,R)\to\operatorname{Ext}_R^r(k,R)\) 满射，矛盾。因此可选 \(x\in\mathfrak m\) 为非零因子。第二事实和第一事实使 \(\operatorname{id}_{\overline R}\overline R=r-1\)：上界由全部素循环商的换环同构得到，下界由 \(\operatorname{Ext}_R^r(k,R)\ne0\) 得到。对 \(r\) 归纳，\(r=0\) 时已有 \(d=0\)，得到 \(R\) 为 CM 环且 \(r=d\)。

归纳还给出条件3。零维且自内射时，任取单射 \(k\to R\)，内射性使 \(\operatorname{Hom}_R(R,R)\to\operatorname{Hom}_R(k,R)\) 满射。后者由该单射生成且被 \(\mathfrak m\) 零化，故同构于 \(k\)，正次 Ext 全为零。正维时，\(\operatorname{Hom}_R(k,R)=0\)，其余各次由第二事实下降到 \(\overline R\)。这证明条件1推出条件3；条件3结合深度的 Ext 刻画显然推出条件2。

最后设条件2成立。取参数正则序列 \(x_1,\ldots,x_d\)，令 \(T=R/(x_1,\ldots,x_d)\)。反复换环得到
\(\operatorname{Soc}(T)\cong\operatorname{Ext}_R^d(k,R)\cong k\)。Artin 局部环中每个非零理想都与基座非零相交：对该理想反复乘以极大理想，取最后一个非零幂次即可。将基座同构 \(\operatorname{Soc}(T)\to k\subseteq E_T(k)\) 延拓到 \(T\to E_T(k)\)，本质性使其为单射；由\cref{b3:lem:matlis-finite-length}，
\(\ell_T E_T(k)=\ell_T D_T(T)=\ell_T T\)，故它是同构，\(T\) 自内射。

沿正则序列逐项返回。若 \(\operatorname{id}_{\overline R}\overline R\le s\)，则对含 \(x\) 的素理想，换环给出 \((s+2)\) 次 Ext 消失；对不含 \(x\) 的素理想，模 \(R/(\mathfrak p,x)\) 被 \(x\) 零化，短正合列使乘以 \(x\) 在 \(\operatorname{Ext}_R^{s+2}(R/\mathfrak p,R)\) 上满射，Nakayama 引理再次使其为零。第一事实于是给出 \(\operatorname{id}_RR\le s+1\)。从 \(T\) 返回 \(R\) 得 \(\operatorname{id}_RR\le d\)，而顶次 Ext 非零给出等号，完成证明。
\end{proof}

对一个 \(d\) 维 CM 局部环 \(R\)，整数
\(\dim_k\operatorname{Ext}_R^d(k,R)\) 称为它的\term[CohenMacaulayleixing]{Cohen–Macaulay 类型}[Cohen--Macaulay type]。深度的 Ext 刻画保证这个整数为正。判别定理说明 Gorenstein 环恰是类型为一的 CM 环；只有“CM”仍不足够。零维环的类型就是其基座的维数，所以 \(k[x,y]/(x,y)^2\) 是 CM 环但不是 Gorenstein 环。

\begin{proposition}\label{b3:prop:gorenstein-completion-localization}
设 \((R,\mathfrak m,k)\) 为交换 Noether 局部环。则 \(R\) 为 Gorenstein 环当且仅当其 \(\mathfrak m\)-进完备化 \(\widehat R\) 为 Gorenstein 环。若 \(R\) 为 Gorenstein 环，则每个 \(R_{\mathfrak p}\)、\(\mathfrak p\in\operatorname{Spec}R\)，也为 Gorenstein 环。
\end{proposition}
\begin{proof}
完备化保维数并忠实平坦，且由\cref{b3:thm:cm-completion}保持及反映 CM 性。取 \(k\) 的有限秩自由 \(R\) 消解，逐项张量 \(\widehat R\) 后仍是 \(k\) 的自由消解。有限秩保证 Hom 与此张量相容，平坦性保证上同调与张量相容，故
\[
\operatorname{Ext}_R^i(k,R)\otimes_R\widehat R
\cong\operatorname{Ext}_{\widehat R}^i(k,\widehat R).
\]
左侧原本被 \(\mathfrak m\) 零化，张量后不变。应用上一判别定理，便得完备化等价性。

为处理局部化，先说明内射模 \(E\) 的局部化 \(E_{\mathfrak p}\) 对 \(R_{\mathfrak p}\) 内射。该局部环的每个理想都是某个有限生成理想 \(J\subseteq R\) 的局部化。Noether 性使 \(J\) 有限展示，从而
\[
\operatorname{Hom}_{R_{\mathfrak p}}(J_{\mathfrak p},E_{\mathfrak p})
\cong\operatorname{Hom}_R(J,E)_{\mathfrak p}.
\]
原内射性使 \(\operatorname{Hom}_R(R,E)\to\operatorname{Hom}_R(J,E)\) 满射，局部化后仍满射，再用 Baer 判据即可。把 \(R\) 的有限内射消解逐项局部化，就得到 \(R_{\mathfrak p}\) 的有限内射消解。
\end{proof}

在给定正则覆盖的情形，典范模提供另一种直接的充分条件，并说明自内射维数为何能由局部对偶控制。

\begin{proposition}\label{b3:prop:canonical-free-gorenstein}
设 \((S,\mathfrak n)\) 为 \(n\) 维交换 Noether 完备正则局部环，\(I\subsetneq S\) 为理想，\(R=S/I\) 为 \(d\) 维 Cohen--Macaulay 环，典范模为
\(\omega_R=\operatorname{Ext}_S^{n-d}(R,S)\)。若 \(\omega_R\cong R\) 为 \(R\) 模，则 \(R\) 为 Gorenstein 环。
\end{proposition}
\begin{proof}
对每个理想 \(J\subseteq R\)，模 \(R/J\) 有限。\cref{b3:thm:canonical-module-local-duality}给出
\[
\operatorname{Ext}_R^{d+1}(R/J,\omega_R)
\cong D_R\bigl(H_{\mathfrak m}^{-1}(R/J)\bigr)=0.
\]
用给定的 \(\omega_R\cong R\) 替换第二变量，再由\cref{b3:lem:injective-dimension-test}得 \(\operatorname{id}_RR\le d\)。此外 CM 性及\cref{b3:thm:depth-ext}使 \(\operatorname{Ext}_R^d(k,R)\ne0\)，所以内射维数恰为 \(d\)。
\end{proof}

此处典范模同构是一个 \(R\) 模同构，通常没有指定的自然选择。把典范模选成 \(R\) 也不等于对每个有限模使用普通线性对偶 \(\operatorname{Hom}_R(-,R)\)：局部对偶中的次数 \(d-i\) 和 Matlis 对偶仍然不可省略。

\subsection{完全交的 Gorenstein 性}
\label{b3:subsec:2203:03}

完全交商的 Koszul 消解有一种对称性：取 \(S\) 线性对偶后，复形各项的顺序倒过来，微分仍是原正则序列给出的矩阵。这使典范模的计算十分具体。

\begin{proposition}\label{b3:prop:ci-canonical-module}
设 \((S,\mathfrak n)\) 为 \(n\) 维交换 Noether 完备正则局部环，\(f_1,\ldots,f_c\in\mathfrak n\) 为 \(S\) 正则序列，\(R=S/(f_1,\ldots,f_c)\)。则 \(d=\dim R=n-c\)，而且
\[
\operatorname{Ext}_S^j(R,S)=
\begin{cases}
R,&j=c,\\
0,&j\ne c.
\end{cases}
\]
特别地，\(R\) 的典范模 \(\omega_R\) 同构于 \(R\)。
\end{proposition}
\begin{proof}
\cref{b3:thm:complete-intersection-cm}及\cref{b3:prop:given-complete-intersection}给出 CM 性和维数。用\cref{b3:prop:ci-koszul-resolution}的 Koszul 消解 \(K_\bullet(f;S)\) 计算 Ext。外积给出完美配对
\[
\bigwedge^jS^c\times\bigwedge^{c-j}S^c
\longrightarrow\bigwedge^cS^c\cong S,
\]
最后一个同构取有序基 \(e_1\wedge\cdots\wedge e_c\mapsto1\)。若 \(u\) 的次数为 \(j\)，Koszul 微分满足
\(\mathrm{d}(u\wedge v)=\mathrm{d}u\wedge v+(-1)^ju\wedge \mathrm{d}v\)。因此配对使对偶微分与倒序 Koszul 微分相差一个只依赖次数的符号；逐次把某些次数上的配对乘以 \(-1\)，便得到复形同构
\[
\operatorname{Hom}_S(K_\bullet(f;S),S)^j
\cong K_{c-j}(f;S).
\]
正则序列使右侧只有链次数零的同调 \(R\)，所以对偶上同调只在 \(j=c\) 为 \(R\)。\(c=0\) 时消解只有 \(S\)，结论同样成立。由于 \(n-d=c\)，这正是典范模定义要求的次数。
\end{proof}

\begin{theorem}\label{b3:thm:complete-intersection-gorenstein}
每个交换 Noether 完全交局部环都是 Gorenstein 环。因此，对交换 Noether 局部环有
\[
\text{正则}\ \Longrightarrow\ \text{完全交}\
\Longrightarrow\ \text{Gorenstein}\
\Longrightarrow\ \text{Cohen--Macaulay}.
\]
\end{theorem}
\begin{proof}
按照\cref{b3:def:local-complete-intersection}，完全交局部环 \(R\) 的完备化具有
\(\widehat R=S/(f_1,\ldots,f_c)\) 的表示，其中 \(S\) 完备正则局部，\(f_1,\ldots,f_c\) 为正则序列。前一命题使 \(\omega_{\widehat R}\cong\widehat R\)，再由\cref{b3:prop:canonical-free-gorenstein}得 \(\widehat R\) 为 Gorenstein 环。\cref{b3:prop:gorenstein-completion-localization}把此性质下降到 \(R\)。链中的第一项由\cref{b3:thm:complete-intersection-cm}证明中的空正则序列论证成立，最后一项来自\cref{b3:thm:gorenstein-bass-criterion}。
\end{proof}

完全交的定义已给出 \(\widehat R\) 的正则局部覆盖及其正则序列关系，所以此处可直接计算典范模。对一般交换完备 Noether 局部环，覆盖的存在性由\appref{b3:app:D}的 Cohen 结构定理保证，但其关系未必由正则序列生成。对于给定的正则局部环 \(S\) 及非零非单位 \(f\)，商 \(S/(f)\) 是超曲面，因而也是 Gorenstein 环；它是否正则仍须另外判断。

% !TeX root = ../../Book3.tex
% 纲要标题：### 22.4 Gorenstein 环的对偶例子
% 纲要位置：计划纲要.md，第 2692 行。

\section{Gorenstein 环的对偶例子}
\label{b3:sec:2204}

零维 Gorenstein 性可以通过基座及乘法配对检验，因而适合具体计算。正维及非 Cohen–Macaulay 情形还要保留对偶的次数；一个典范模未必足以记录全部数据。


% 纲要内容（仅作写作计划，尚非教材正文）：
% * Artin Gorenstein 环的基座与乘法配对
% * 超曲面及完全交的环境环 Ext 计算
% * 环境环上的对偶复形：有限长度和各项有限秩的自由消解、对偶微分及移位；区分各项的 $S$ 模结构与上同调的 $R$ 模结构
% * 非 CM 例子的两个 Ext 次数与局部上同调的 Matlis 对应
% ---

\subsection{Artin Gorenstein 环}
\label{b3:subsec:2204:01}

零维时，Gorenstein 性可以通过一个很小的线性空间辨认。基座包含那些被极大理想全部零化的元素；它们是乘法再也不能向下推进的部分。

\begin{proposition}\label{b3:prop:artinian-gorenstein-socle}
设 \((R,\mathfrak m,k)\) 为非零交换 Artin 局部环，并记 \(\operatorname{Soc}(R)=(0:_R\mathfrak m)\)。则下列条件等价：\(R\) 为 Gorenstein 环；\(R\) 为内射 \(R\) 模；\(\dim_k\operatorname{Soc}(R)=1\)。
\end{proposition}
\begin{proof}
零维时，\cref{b3:thm:gorenstein-bass-criterion}给出内射维数为零，并使
\(\operatorname{Hom}_R(k,R)=\operatorname{Soc}(R)\) 为一维。下面也直接说明基座条件如何产生自内射性。Artin 局部环完备，其有限长度模 \(R\) 的 Matlis 对偶 \(D_R(R)=E_R(k)\) 与 \(R\) 有相同长度；这是\cref{b3:thm:matlis-duality}在组成列上的应用，因为 \(D_R(k)\cong k\) 且对偶正合。

若基座一维，把 \(\operatorname{Soc}(R)\cong k\) 嵌入 \(E_R(k)\)，利用内射性将它延拓为 \(u:R\to E_R(k)\)。任一非零理想 \(J\) 都与基座非零相交：因 \(\mathfrak m\) 幂零，选最大 \(r\) 使 \(\mathfrak m^rJ\ne0\)，则该模被 \(\mathfrak m\) 零化。因此 \(\ker u\) 若非零，必与基座相交，与 \(u\) 在基座上单射矛盾。故 \(u\) 单射，两端长度相同使它为同构，\(R\) 遂内射。
\end{proof}

若 \(R=A\) 还是剩余域为 \(k\) 的有限维交换局部 \(k\) 代数，上述结论可直接写成乘法配对。选 \(k\) 线性泛函 \(\lambda:A\to k\)，使它在一维基座上非零。则
\[
A\longrightarrow\operatorname{Hom}_k(A,k),\qquad
a\longmapsto\bigl(b\longmapsto\lambda(ab)\bigr)
\]
单射：若 \(a\ne0\)，理想 \(Aa\) 包含非零基座元素，所以存在 \(b\) 使 \(\lambda(ab)\ne0\)。两端维数相同，故这是同构。这里的目标模取作用 \((c\phi)(b)=\phi(cb)\)；它内射，因为
\[
\operatorname{Hom}_A(N,\operatorname{Hom}_k(A,k))\cong\operatorname{Hom}_k(N,k).
\]
同构把 \(f\) 送到 \(n\mapsto f(n)(1)\)，逆映射把 \(u\) 送到 \(n\mapsto[a\mapsto u(an)]\)。向量空间对偶正合，从而此 Hom 函子正合。

\ProseParagraphBreak
例如 \(A=k[x,y]/(x^a,y^b)\)、\(a,b\ge1\)，基座由 \(x^{a-1}y^{b-1}\) 生成。取 \(\lambda\) 为此单项式的系数，基单项式 \(x^iy^j\) 与 \(x^{a-1-i}y^{b-1-j}\) 配对为一，其余达到不同次数的乘积对该系数贡献为零。这个有限维配对把“类型一”的条件变成能够逐项检验的矩阵。同样长度很短的 \(k[x,y]/(x,y)^2\) 却有二维基座 \(kx+ky\)，所以不存在这样的非退化乘法配对。

\ProseParagraphBreak
对后一环 \(A=k[x,y]/(x,y)^2\)，典范模仍然存在，但不再同构于环本身。令 \(W=\operatorname{Hom}_k(A,k)\)，以 \(\phi_0,\phi_x,\phi_y\) 表示基 \(1,x,y\) 的对偶基。按照 \((a\phi)(b)=\phi(ab)\) 的作用，直接计算得
\[
\begin{aligned}
x\phi_x&=\phi_0,& y\phi_y&=\phi_0,\\
x\phi_y&=0,& y\phi_x&=0,\\
x\phi_0&=0,& y\phi_0&=0.
\end{aligned}
\]
于是 \(\operatorname{Soc}_A(W)=k\phi_0\)。任意非零子模中，若一个元素的 \(\phi_x\) 或 \(\phi_y\) 系数非零，乘 \(x\) 或 \(y\) 就得到非零的 \(\phi_0\)；否则它本身已在 \(k\phi_0\) 中。因此 \(k\phi_0\subseteq W\) 本质。上面的 Hom 计算又说明 \(W\) 内射，故 \(W\cong E_A(k)=D_A(A)\)。零维时 \(H^0_{\mathfrak m}(A)=A\)，Matlis 双对偶遂使 \(W\) 成为典范模。可是 \(\mathfrak mW=k\phi_0\)，所以 \(W/\mathfrak mW\) 为二维，\(W\) 至少需要两个生成元；\(A\) 作为自身模只需一个。这个作用表具体说明了 \(\omega_A\not\cong A\)，也区分了典范模的存在与 Gorenstein 性。

\subsection{超曲面与完全交的对偶例子}
\label{b3:subsec:2204:02}

设 \(k\) 为域，取 \(S=k[\![x,y]\!]\)、\(R=S/(xy)\)。环境环中的自由消解是
\[
0\longrightarrow S\xrightarrow{xy}S\longrightarrow R\longrightarrow0.
\]
对偶后仍为乘 \(xy\)，所以 \(\omega_R=\operatorname{Ext}_S^1(R,S)\cong R\)。这是一维 Gorenstein 环，但在原点有两条分支，因而不是整环，也不是正则局部环。Gorenstein 条件描述对偶性质，不能替代不可约性或光滑性。

\ProseParagraphBreak
该例的局部上同调也可以算出。因为在 \(R\) 中 \(xy=0\)，有 \(R_{xy}=0\)，所以两元 Čech 复形给出
\[
H_{\mathfrak m}^1(R)=\bigl(R_x\oplus R_y\bigr)/R
=\bigl(k((x))\oplus k((y))\bigr)/R.
\]
此处 \(R\) 在两支中的像是所有常数项相等的幂级数对：任一元素唯一写成
\(c+x\cdot a(x)+y\cdot b(y)\)。故上式中的 \(R\to R_x\oplus R_y\) 为单射，\(H_{\mathfrak m}^0(R)=0\)。局部对偶进一步给出
\(D_R(H_{\mathfrak m}^1(R))\cong R\)，并由 Matlis 双对偶得到
\(H_{\mathfrak m}^1(R)\cong E_R(k)\)。最后这个同构反映的是内射包络结构，不应被误写成 \(H_{\mathfrak m}^1(R)\cong R\)。

\ProseParagraphBreak
再设 \(k\) 为域，取 \(S=k[\![x,y,z]\!]\)、\(R=S/(x^2,y^3)\)。以 \((x^2,y^3)\) 为序，消解为
\[
0\longrightarrow S\xrightarrow{\binom{-y^3}{x^2}}S^2
\xrightarrow{(x^2\ \ y^3)}S\longrightarrow R\longrightarrow0.
\]
取对偶后，零次核为零；若 \(-y^3a+x^2b=0\)，则互素性给出
\((a,b)=(x^2c,y^3c)\)，故一次上同调为零；二次余核为 \(R\)。所以 \(\omega_R\cong R\)，且 \(\dim R=1\)。参数 \(z\) 为非零因子，商去 \(z\) 后得到前面有非退化配对的 Artin 环 \(k[x,y]/(x^2,y^3)\)。

\ProseParagraphBreak
这些对偶计算全部在环境环 \(S\) 上进行。商环 \(R\) 上的剩余域自由消解通常无限长；例如 \(k[\varepsilon]/(\varepsilon^2)\) 上由乘 \(\varepsilon\) 构成的无限消解已在\secref{b3:sec:1903}计算，而环本身却内射。必须在记号中保留 \(\operatorname{Ext}_S\)、\(\operatorname{Ext}_R\) 的下标，才能区分这两项事实。

\subsection{环境环上的对偶复形}
\label{b3:subsec:2204:03}

对完备正则局部环 \(S\) 的一个 CM 商 \(R=S/I\)，自由消解的对偶只有一个非零上同调模，因而典范模能够集中表达对偶信息。若 \(R\) 不是 CM 环，不同次数上的对偶信息可能同时存在；这时保留整个对偶复形，比只取最高次局部上同调的对偶更合适。

\begin{definition}\label{b3:def:regular-cover-dual-complex}
设 \((S,\mathfrak n,k)\) 为 \(n\) 维交换完备 Noether 正则局部环，\(I\subsetneq S\) 为理想，\(R=S/I\)。取有限长度、各项均为有限秩自由 \(S\) 模的消解 \(P_\bullet\rightarrow R\)，微分记为 \(\mathrm d_j:P_j\rightarrow P_{j-1}\)。令
\(C^j=\operatorname{Hom}_S(P_j,S)\)，微分为
\(\delta^j(\varphi)=\varphi\circ\mathrm d_{j+1}\)。把
\[
\operatorname{Hom}_S(P_\bullet,S)[n]
\]
看作上链复形，称它为给定满射表示 \(S\twoheadrightarrow R\) 下的一个归一化对偶复形模型。移位约定为
\(C[n]^j=C^{j+n}\)、\(\mathrm{d}_{C[n]}=(-1)^n\delta\)，所以对每个整数 \(i\)，该模型的上同调为
\[
H^{-i}(C[n])=\operatorname{Ext}_S^{n-i}(R,S).
\]
未出现的 \(P_j\) 及负次 Ext 均记为零。
\end{definition}

上述有限消解可以取得：\cref{b3:thm:regular-local-global-dimension}给出正则局部环 \(S\) 的整体维数为 \(n\)，而有限生成模的极小自由消解在每一项只有有限个生成元。若 \(P_\bullet\) 与 \(Q_\bullet\) 都是这样的消解，恒等映射 \(R\rightarrow R\) 可提升为两个方向的比较映射；其复合与各自的恒等映射链同伦。取 \(S\) 对偶并移位后，仍得到互为逆的上链同伦等价。因此这里的消解选择不改变模型的链同伦类型，也不改变各次上同调。

\ProseParagraphBreak
该模型的各项是 \(S\) 模，未必是 \(R\) 模，因为 \(I\) 未必零化这些自由项。不过，对任意 \(a\in I\)，消解上的乘 \(a\) 映射诱导 \(R\) 上的零映射，比较映射的同伦唯一性使它链同伦于零。对偶后乘 \(a\) 同样同伦于零，故 \(I\) 零化每个上同调模。这才给出上述 Ext 模的 \(R\) 模结构。

\ProseParagraphBreak
若 \(R\) 为 \(d\) 维 CM 环，\cref{b3:thm:regular-local-duality}及局部上同调的消失说明只有
\(\operatorname{Ext}_S^{n-d}(R,S)=\omega_R\) 非零。因此归一化后的模型在次数 \(-d\) 的上同调为 \(\omega_R\)，其余次数均为零；这个次数与局部对偶公式中的 \(d-i\) 相对应。

\ProseParagraphBreak
最后设 \(k\) 为域，取 \(S=k[\![x,y]\!]\)、\(R=S/(x^2,xy)\)，并记 \(\mathfrak m=(x,y)R\)。若 \(x^2a+xyb=0\)，在整环 \(S\) 中消去 \(x\) 得 \(xa+yb=0\)；由于 \(x,y\) 互素，必有 \((a,b)=c(-y,x)\)。故有自由消解
\[
0\longrightarrow S\xrightarrow{\binom{-y}{x}}S^2
\xrightarrow{(x^2\ \ xy)}S\longrightarrow R\longrightarrow0.
\]
取 \(S\) 对偶后，两个微分分别是
\[
S\xrightarrow{\binom{x^2}{xy}}S^2
\xrightarrow{(-y\ \ x)}S.
\]
若 \(-ya+xb=0\)，同样的整除论证给出 \((a,b)=c(x,y)\)。所以一次核由 \((x,y)\) 生成，而前一像由 \(x(x,y)\) 生成；二次余核则是 \(S/(x,y)\)。因此
\[
\operatorname{Ext}_S^1(R,S)\cong S/(x),\qquad
\operatorname{Ext}_S^2(R,S)\cong S/(x,y)=k.
\]
这里 \(S\) 的维数为二，归一化对偶模型在次数 \(-1\) 和零都有非零上同调。\subsecref{b3:subsec:2201:03}给出的两个局部上同调模正好与之对应。具体地，\(R\) 为完备 Noether 局部环，取
\(D_R(-)=\operatorname{Hom}_R(-,E_R(k))\)。由\cref{b3:thm:regular-local-duality}及商环上的对偶识别，得到\cref{b3:tab:non-cm-dual-degrees}中的两项同构；Ext 的下标仍是环境环 \(S\)。

\begin{table}[htbp]
\centering
\caption{非 Cohen--Macaulay 环的两个对偶次数}
\label{b3:tab:non-cm-dual-degrees}
\begin{tabularx}{\linewidth}{@{}cXX@{}}
\toprule
次数 \(i\) & 局部上同调 \(H^i_{\mathfrak m}(R)\) & Matlis 对偶 \(D_R(H^i_{\mathfrak m}(R))\) \\
\midrule
\(1\) & \(k((y))/k[\![y]\!]\) & \(\operatorname{Ext}_S^1(R,S)\cong S/(x)\) \\
\(0\) & \(kx\cong k\) & \(\operatorname{Ext}_S^2(R,S)\cong k\) \\
\bottomrule
\end{tabularx}
\end{table}

典范模 \(\omega_R=\operatorname{Ext}_S^1(R,S)\) 记录了一次局部上同调的对偶，却没有记录由 \(kx\) 产生的另一个非零次数。这个例子说明，非 CM 环的局部对偶信息通常需要多个模及其次数，不能仅由一个典范模替代。


\begin{chaptersummary}
局部上同调来自一个左正合函子的导出，但在 Noether 环上可由有限的局部化复形计算。局部环上有限生成模的深度给出首次非零次数，维数给出最高可能次数；CM 性恰好把全部信息集中到一个次数。Matlis 对偶使这些通常不有限生成的 Artin 模，能够与有限生成模相互转换，完备性保证双对偶返回原模。

\ProseParagraphBreak
给定从完备正则局部环到所研究局部环的满射后，环境环上的 Ext 构造典范模；CM 假设使其余次数消失，从而得到含单个典范模的局部对偶。完全交的 Koszul 自对偶给出 \(\omega_R\cong R\)，继而推出有限自内射维数。零维时，这一条件又成为基座一维和非退化乘法配对。非 CM 例子保留多个非零 Ext，说明一般对偶理论为何必须研究复形。
\end{chaptersummary}

\BookExercises

\begin{optionalexercise}\label{b3:ex:22-unbounded-torsion}
设 \(k\) 为域。在 \(R=k[x]\) 上令 \(M=\bigoplus_{n\ge1}R/(x^n)\)。证明 \(\Gamma_{(x)}(M)=M\)，但不存在统一的 \(N\) 使 \(x^NM=0\)。再计算全部 \(H_{(x)}^i(M)\)。
\end{optionalexercise}
\begin{solution}
每个元素只有有限个非零分量，取这些分量指标的最大值即可零化它。对于任意 \(N\)，第 \(N+1\) 个分量中的一不被 \(x^N\) 零化，所以没有统一指数。局部化 \(M_x=0\)，两项 Čech 复形为 \(M\to0\)，故零次为 \(M\)，所有正次数为零。
\end{solution}

\begin{optionalexercise}\label{b3:ex:22-principal-nonreduced}
设 \(k\) 为域。令 \(R=k[x,y]/(xy)\)，\(I=(x)\)。求 \(H_I^0(R)\)、\(H_I^1(R)\)，并说明它们分别记录了哪一部分元素。
\end{optionalexercise}
\begin{solution}
局部化后 \(x\) 可逆迫使 \(y=0\)，所以 \(R_x=k[x,x^{-1}]\)。映射 \(R\to R_x\) 的核为 \((y)\)，像为 \(k[x]\)。故
\[
H_I^0(R)=(y),\qquad
H_I^1(R)=k[x,x^{-1}]/k[x].
\]
第一项记录 \(y\) 轴这一不可约分量上被 \(x\) 零化的元素；第二项由负次幂的类组成，记录局部化后新增、却不能由原环提升的元素。两项性质不同，不能把局部上同调全都看成原模的子模。
\end{solution}

\begin{optionalexercise}\label{b3:ex:22-two-variable-cech}
设 \(k\) 为域，\(R=k[x,y]\)、\(\mathfrak m=(x,y)\)、\(M=R/(xy)\)。用二元 Čech 复形计算全部 \(H_{\mathfrak m}^i(M)\)，并给出每个非零模的 \(k\) 基。说明为什么一次上同调除两个分支的负次幂之外，还含有一份常数差。
\end{optionalexercise}
\begin{solution}
\(M_x\cong k[x,x^{-1}]\)、\(M_y\cong k[y,y^{-1}]\)，而同时逆转 \(x,y\) 会使 \(xy=0\) 变成单位为零，故 \(M_{xy}=0\)。每个 \(M\) 元素可唯一写成 \(c+xf(x)+yg(y)\)，它在两个局部化中的像为 \((c+xf(x),c+yg(y))\)。这个映射单射，且其像恰是常数项相等的多项式对。因此零次上同调为零，所有二次及更高上同调也为零；一次上同调是这份直和模掉上述像。

令 \(\delta\) 为 \((1,0)\) 的类，则 \((0,1)\) 的类等于 \(-\delta\)。全部 \(k\) 基为
\[
\delta,\quad (x^{-a},0)\ (a\ge1),\quad(0,y^{-b})\ (b\ge1).
\]
正次幂可从像中消去，负次幂不可能出现在像中，常数对则只商去对角线，故这些类线性无关且穷尽商。常数差 \(\delta\) 记录了两分支在原点处的兼容条件；它被 \(x,y\) 同时零化，但仍是一次上同调中的类，并非 \(M\) 自身的幂挠元素。
\end{solution}

\begin{optionalexercise}\label{b3:ex:22-flat-base-change}
设 \(k\) 为域，\(R=k[x]\)、\(I=(x)\)、\(B=R/(x)\)。计算 \(H_I^i(R)\otimes_RB\) 与 \(H_{IB}^i(B)\)，指出不平坦的基变换在哪个次数破坏局部上同调公式。再核对两项 Čech 复形张量后的各项，说明逐项基变换仍成立为何不足以保证上同调基变换。
\end{optionalexercise}
\begin{solution}
在 \(R\) 上，\(H_I^0(R)=0\)、\(H_I^1(R)=R_x/R\)，其余为零。乘 \(x\) 在 \(R_x/R\) 上满射，所以 \((R_x/R)\otimes_RB=(R_x/R)/x(R_x/R)=0\)。左边全部为零。\(IB=0\)，故 \(\Gamma_{IB}\) 是恒等函子，\(H_{IB}^0(B)=B\)，所有正次为零；公式在零次失败。

原 Čech 复形 \(R\to R_x\) 的张量为 \(B\to0\)，因为 \(R_x\otimes_RB=B_x=0\)。它也正是由 \(x\) 在 \(B\) 中的像构造的 Čech 复形，逐项比较并没有失败。但取上同调后得到 \(B\)，与先取原复形的零次上同调再张量所得的零模不同。不平坦张量没有保持原单射 \(R\to R_x\) 的核，这正是差别所在。
\end{solution}

\begin{optionalexercise}\label{b3:ex:22-remove-finite-torsion}
设 \((R,\mathfrak m)\) 为交换 Noether 局部环，\(M\) 为有限生成模，\(T=\Gamma_{\mathfrak m}(M)\)。证明 \(T\) 有有限长度，且 \(\Gamma_{\mathfrak m}(M/T)=0\)。若 \(M/T\ne0\)，推出它的深度至少为一。
\end{optionalexercise}
\begin{solution}
有限生成模 \(M\) 中的升链 \(0:_M\mathfrak m\subseteq0:_M\mathfrak m^2\subseteq\cdots\) 最终稳定，故 \(T\) 被某个 \(\mathfrak m^a\) 零化；其幂滤过的相邻商为有限维剩余域空间，所以长度有限。若 \(m+T\) 被 \(\mathfrak m^b\) 零化，则 \(\mathfrak m^bm\subseteq T\)，再乘 \(\mathfrak m^a\) 得零。因此 \(m\in T\)，所求商的挠子模为零。商非零时，深度与局部上同调的首次非零次数相等，而零次已消失，故深度至少一。
\end{solution}

\begin{optionalexercise}\label{b3:ex:22-dvr-dual-maps}
设 \(k\) 为域。令 \(R=k[\![t]\!]\)、\(E=k(\!(t)\!)/k[\![t]\!]\)、\(D_R=\operatorname{Hom}_R(-,E)\)，并设 \(n\ge1\) 为整数。把 \(D_R(R/(t^n))\) 写成 \(E\) 的具体子模，并求商映射 \(R/(t^{n+1})\to R/(t^n)\) 的 Matlis 对偶。
\end{optionalexercise}
\begin{solution}
同态由一的像决定，该像必须被 \(t^n\) 零化。因此
\[
D_R(R/(t^n))=(0:_Et^n)=t^{-n}R/R,
\]
其基为 \(t^{-1},\ldots,t^{-n}\) 的类。商映射的对偶是把同一个一的像看作被 \(t^{n+1}\) 零化的元素，所以是自然包含
\(t^{-n}R/R\hookrightarrow t^{-(n+1)}R/R\)。
若分别用 \(R/(t^n)\to t^{-n}R/R\)、\(a\mapsto at^{-n}+R\) 作坐标，则这个包含对应 \(a\mapsto ta\)。因此有限商的逆系统被对偶变为带乘 \(t\) 映射的正向系统。
\end{solution}

\begin{optionalexercise}\label{b3:ex:22-matlis-presentation}
设 \((R,\mathfrak m)\) 为交换完备 Noether 局部环，\(\kappa=R/\mathfrak m\)、\(E=E_R(\kappa)\)、\(D_R=\operatorname{Hom}_R(-,E)\)。有限生成的 \(R\) 模 \(M\) 有展示 \(R^a\xrightarrow{A}R^b\to M\to0\)。用矩阵 \(A\) 表示 \(D_R(M)\)。再设 \(k\) 为域，在 \(R=k[\![t]\!]\) 上计算 \(D_R(R/(t^2)\oplus R)\)。
\end{optionalexercise}
\begin{solution}
Matlis 对偶反变正合，所以
\[
0\longrightarrow D_R(M)\longrightarrow E^b\xrightarrow{A^{\mathsf T}}E^a
\]
正合，所求为转置矩阵作用的核。第二问中得到
\(D_R(R/(t^2)\oplus R)=t^{-2}R/R\oplus E\)。
第一直和项长度二，第二项 Artin 却非有限长度；因此有限生成模的对偶不能一律称作有限长度模。
\end{solution}

\begin{optionalexercise}\label{b3:ex:22-completion-in-double-dual}
设 \(k\) 为域，\(R=k[t]_{(t)}\)、\(E=k(t)/R\)、\(D_R=\operatorname{Hom}_R(-,E)\)，并取整数 \(q\ge1\)。对 \(M=R/(t^q)\oplus R\)，计算 \(D_R(D_R(M))\) 及自然求值映射的余核。解释有限长度直和项与自由直和项在非完备环上为何有不同的双对偶表现。
\end{optionalexercise}
\begin{solution}
\(E\) 是剩余域的内射包络。\Cref{b3:lem:matlis-finite-length}的有限长度求值同构把第一直和项的双对偶识别为自身，内射包络自同态的计算则把 \(D_R(D_R(R))\) 识别为 \(\widehat R=k[\![t]\!]\)。有限直和与两次对偶相容，因此
\[
D_R(D_R(M))\cong R/(t^q)\oplus k[\![t]\!].
\]
在这个识别下，求值映射在第一项上为恒等，在第二项上为完备化嵌入 \(R\hookrightarrow k[\![t]\!]\)，故余核为 \(k[\![t]\!]/R\)，非零。有限长度模已经只依赖某个有限阶商，完备化不会增加这项的数据；自由直和项则允许无限阶相容系数，所以双对偶返回其完备化。
\end{solution}

\begin{optionalexercise}\label{b3:ex:22-finite-length-local-duality}
设 \((S,\mathfrak n)\) 为 \(n\) 维交换完备 Noether 正则局部环，\(I\subsetneq S\) 为理想，\(R=S/I\) 为 \(d\) 维 CM 局部环。记 \(\mathfrak m=\mathfrak n/I\)、\(\kappa=R/\mathfrak m\)、\(D_R=\operatorname{Hom}_R(-,E_R(\kappa))\)，并取典范模 \(\omega_R=\operatorname{Ext}_S^{n-d}(R,S)\)。若 \(L\) 为有限长度 \(R\) 模，证明
\[
\operatorname{Ext}_R^j(L,\omega_R)=0\quad(j\ge0,\ j\ne d),\qquad
\operatorname{Ext}_R^d(L,\omega_R)\cong D_R(L).
\]
比较后一模与 \(L\) 的长度。
\end{optionalexercise}
\begin{solution}
\(L\) 有有限长度，故存在 \(q\) 使 \(\mathfrak m^qL=0\)。极大理想 \(\mathfrak m\) 的每个生成元 \(x_i\) 都满足 \(x_i^qL=0\)；局部化后 \(x_i\) 可逆，所以 \(L_{x_i}=0\)。因此 Čech 复形只有次数零的 \(L\)，其零次上同调为 \(L\)，正次上同调为零。局部对偶于是给出所列公式。Matlis 对偶正合且将剩余域送到同构的剩余域；对 \(L\) 的组成列逐项取对偶，因子数不变，故 \(\ell_R(D_R(L))=\ell_R(L)\)。
\end{solution}

\begin{optionalexercise}\label{b3:ex:22-quadric-canonical}
设 \(k\) 为域。在 \(S=k[\![x,y,z]\!]\) 中令 \(R=S/(xy-z^2)\)。计算其典范模并证明 \(R\) 为 Gorenstein 环；再说明它不是正则局部环。此结论是否要求 \(\operatorname{char}k\ne2\)？
\end{optionalexercise}
\begin{solution}
环境环为三维正则局部整环，\(f=xy-z^2\ne0\) 为非零因子，故商为二维 CM 环。两项消解 \(0\to S\xrightarrow{f}S\to R\to0\) 的对偶在次数一的余核为 \(R\)，所以 \(\omega_R=\operatorname{Ext}_S^1(R,S)\cong R\)，由典范模判据为 Gorenstein 环。关系属于极大理想的平方，因此商环的 \(\mathfrak m/\mathfrak m^2\) 仍为三维，而环维数为二，故不正则。所有论证对任意特征成立，没有除以二的步骤。
\end{solution}

\begin{optionalexercise}\label{b3:ex:22-artinian-pairing}
设 \(k\) 为域。令 \(A=k[u,v]/(u^2,v^3)\)，取有序基 \(1,u,v,uv,v^2,uv^2\)。写出配对 \((a,b)\mapsto\lambda(ab)\) 的矩阵，其中 \(\lambda\) 提取 \(uv^2\) 的系数。由此确定 \(A\) 的基座及 Gorenstein 性。
\end{optionalexercise}
\begin{solution}
基中元素的互补对分别是 \(1,uv^2\)、\(u,v^2\)、\(v,uv\)，因此矩阵为反对角线上全为一、其余为零的 \(6\times6\) 矩阵，行列式为 \((-1)^{15}\ne0\)。被 \(u,v\) 同时零化的元素恰为 \(k uv^2\)：乘 \(u\) 为零先去掉不含 \(u\) 的各项，再乘 \(v\) 为零只留下最高 \(v\) 次数。故基座一维，环为 Gorenstein；配对直接实现 \(A\cong\operatorname{Hom}_k(A,k)\)。
\end{solution}

\begin{optionalexercise}\label{b3:ex:22-three-generators-gorenstein}
设 \(k\) 为域。令
\[
A=k[x,y,z]/(xy,xz,yz,x^2-y^2,x^2-z^2).
\]
证明 \(A\) 是作为 \(k\) 向量空间五维的局部 \(k\) 代数，求其基座，判断其 Gorenstein 性。
\end{optionalexercise}
\begin{solution}
写 \(s=x^2=y^2=z^2\)。所有三次单项式为零：含不同变量的乘积由前三个关系为零，而 \(x^3=xy^2=0\)，其余同理。次数零、一、二分别有维数一、三、一，故基为 \(1,x,y,z,s\)。正次理想幂零，商为 \(k\)，所以它是唯一极大理想。若 \(a=a_0+a_1x+a_2y+a_3z+a_4s\) 被 \(x,y,z\) 同时零化，则先得 \(a_0=0\)，再由三个乘积分别得 \(a_1=a_2=a_3=0\)。因此基座为 \(ks\)，由基座判据为 Gorenstein 环。此计算没有根据“有几个写出的方程”判定完全交性质。
\end{solution}

\begin{optionalexercise}\label{b3:ex:22-cm-not-gorenstein-positive-dimension}
设 \(k\) 为域，令 \(R=k[\![u,v,w]\!]/(u^2,uv,v^3)\)、\(\mathfrak m=(u,v,w)R\)。计算 \(\dim R\)、\(\operatorname{depth}R\) 和 \(\dim_k\operatorname{Ext}_R^1(k,R)\)，并判断 CM 性与 Gorenstein 性。
\end{optionalexercise}
\begin{solution}
作为 \(k[\![w]\!]\) 模，\(R\) 以 \(1,u,v,v^2\) 为自由基，故 \(w\) 为非零因子。根理想为 \((u,v)\)，所以 \(\dim R=1\)，深度也为一，\(R\) 为 CM 环。商 \(B=R/wR\) 的基仍为 \(1,u,v,v^2\)，其基座为 \(ku\oplus kv^2\)：乘 \(u\) 为零迫使常数项为零，再乘 \(v\) 为零迫使一次 \(v\) 项为零。正则元素换环给出
\[
\operatorname{Ext}_R^1(k,R)\cong\operatorname{Hom}_B(k,B)
=\operatorname{Soc}(B),
\]
所以 CM 类型为二。\Cref{b3:thm:gorenstein-bass-criterion}的类型一判据于是排除 Gorenstein 性；也可由正则商保持 Gorenstein 性与 \(B\) 的二维基座得出相同结论。这里深度已经达到维数，失败来自类型而非 CM 条件。
\end{solution}

\begin{optionalexercise}\label{b3:ex:22-annihilator-duality}
设 \(k\) 为域，\(A\) 为 \(k\) 上的有限维交换局部 Gorenstein 代数，且剩余域为 \(k\)，\(J\subseteq A\) 为理想。证明
\[
\dim_k(0:_AJ)=\dim_k(A/J),\qquad
0:_A(0:_AJ)=J.
\]
\end{optionalexercise}
\begin{solution}
取正文中的非退化乘法配对 \(\lambda(ab)\)。\(J\) 的正交补恰是 \((0:_AJ)\)：一个方向直接成立；反之，若 \(\lambda(aJ)=0\) 但 \(aJ\ne0\)，非零理想 \(aJ\) 与一维基座相交，而 \(\lambda\) 在基座上非零，矛盾。非退化有限维配对给出正交补维数等于总维数减去 \(\dim_kJ\)，得到第一式；再取正交补恢复原子空间，得到第二式。理想条件用于使 \(aJ\) 仍为理想，不能把零化子结论照搬到任意线性子空间。
\end{solution}

\begin{optionalexercise}\label{b3:ex:22-longer-embedded-component}
设 \(k\) 为域。在 \(S=k[\![x,y]\!]\) 中令 \(R=S/(x^2,xy^2)\)、\(\mathfrak m=(x,y)R\)。求 \(H_{\mathfrak m}^0(R)\)、\(H_{\mathfrak m}^1(R)\)，以及 \(\operatorname{Ext}_S^1(R,S)\)、\(\operatorname{Ext}_S^2(R,S)\)。比较两组计算。
\end{optionalexercise}
\begin{solution}
每个元素唯一写成 \(a(y)+x(b_0+b_1y)\)。局部化于 \(y\) 后 \(x=0\)，所以
\[
H_{\mathfrak m}^0(R)=kx\oplus kxy,\qquad
H_{\mathfrak m}^1(R)=k(\!(y)\!)/k[\![y]\!].
\]
这里 \(\sqrt{(y)}=\mathfrak m\)，故两项 Čech 复形足够。另一方面，两生成元的全部关系由 \((-y^2,x)\) 生成，得到长度二的环境环自由消解。对偶复形中，\(-y^2a+xb=0\) 等价于 \((a,b)=(xc,y^2c)\)，前一像是此核的 \(x\) 倍，故
\[
\operatorname{Ext}_S^1(R,S)=S/(x),\qquad
\operatorname{Ext}_S^2(R,S)=S/(x,y^2).
\]
第二模长度二，与零次局部上同调的 Matlis 对偶长度相同；第一模对应一次局部上同调。比正文长度一的嵌入部分，这个例子还记录了非零作用 \(y\cdot x=xy\)，不只记录一个维数。
\end{solution}

\begin{optionalexercise}\label{b3:ex:22-cm-module-environment-ext}
设 \((S,\mathfrak n)\) 为 \(n\) 维交换完备 Noether 正则局部环，\(M\ne0\) 为有限生成的 \(S\) 模，\(r=\dim_SM\)。证明 \(M\) 为 CM 模当且仅当 \(\operatorname{Ext}_S^j(M,S)\) 仅在 \(j=n-r\) 非零。
\end{optionalexercise}
\begin{solution}
正则局部对偶把 \(\operatorname{Ext}_S^j(M,S)\) 识别为 \(D_S(H_{\mathfrak n}^{n-j}(M))\)。Matlis 对偶是忠实的：非零模的非零元素可由映入内射包络的同态检测，所以一个模为零当且仅当其对偶为零。因此 Ext 集中在 \(n-r\) 等价于局部上同调集中在 \(r\)。由深度首次非零和维数以上消失的两个定理，这又恰好等价于 \(\operatorname{depth}_SM=r\)。在两种情形中所指的唯一一项确实非零，而不只是说其他项消失。
\end{solution}

% !TeX root = ../../Book3.tex
% 纲要标题：## 第23章*　Hensel 性、非分歧与平展
% 纲要位置：计划纲要.md，第 2218 行。

\OptionalChapter{Hensel 性、非分歧与平展}{b3:ch:23}
\BookChapterNumber{b3:ch:23}

\begin{chapterabstract}
本章证明完备局部环中的单根及互素因式提升，比较 Hensel 化与完备化。平方零理想上的映射提升通向形式光滑的关系模判据、光滑映射的局部标准形与平展的刻画；尖点正规化的计算用于区分正则性、光滑性、非分歧性和平坦性。
\end{chapterabstract}

\begin{learninggoals}
\begin{itemize}
\item 用误差修正证明单根提升与互素因式提升，并辨认导数和互素条件的作用。
\item 根据泛性质区分 Hensel 化与完备化，给出有限阶商相同而环不同的例子。
\item 由平方零提升建立关系模的分裂，并把有限展示转化为邻域上的有限方程组。
\item 证明平展的几种刻画等价，区分平坦性与微分模消失这两个条件。
\item 计算尖点正规化的微分模、特殊纤维及原点外的逆映射，说明光滑性为何必须连同底环一起讨论。
\end{itemize}
\end{learninggoals}

方程 \(T^2-3=0\) 在模 \(11\) 下有根 \(5\)。若寻找模 \(121\) 的根，写成 \(5+11u\)，代入后得到 \(2+10u\equiv0\pmod{11}\)，故取 \(u=2\)，得到 \(27^2-3=726=6\cdot121\)。修正量可以唯一解出，是因为导数 \(2T\) 在原根处模 \(11\) 可逆。一次这样的运算还不是完整提升：要反复作下去，并在极限中找到元素，需要环能容纳相容的逐阶近似。

\ProseParagraphBreak

同样的思路也适用于提升多项式的互素因子：误差逐步进入更高的理想幂，而导数可逆改由因子的互素性承担。若把待提升的对象换成环同态，问题变成平方零扩张上的存在与唯一性；微分模测量其中的一阶自由度。把求根与映射提升放在一起看，才能区分 Hensel 性、形式光滑性和平展性所要求的不同条件。

\ProseParagraphBreak

\chapref{b3:ch:13}的完备化、\chapref{b3:ch:14}的微分与形式提升，以及\chapref{b3:ch:12}的平坦性是本章的主要基础；局部性质比较还使用\chapref{b3:ch:18}的正则性。第14章已经给出形式光滑、形式非分歧、形式平展和有限展示条件下的定义，本章证明这些条件之间的联系。

\ProseParagraphBreak
单根与因式提升、关系模判据、局部标准形、平展的等价刻画及尖点计算均在正文证明。Hensel 化的存在及基本性质、Noether 局部形式光滑蕴含平坦、完美基域上的正则与光滑比较分别给出证明概要，并注明省略的技术步骤及出处。Hensel 化的构造只用于\secref{b3:sec:2301}，局部平坦性定理用于一般标准模型，正则与光滑的比较用于最后一节。

\ProseParagraphBreak
Hensel 因式提升见松村英之\cite[Theorem 8.3, p. 58]{Matsumura1989CommutativeRingTheory}，Hensel 化的构造与泛性质见永田雅宜\cite[\S43, (43.3), (43.5), (43.8)--(43.10), pp. 180--183]{Nagata1975LocalRings}。形式光滑与关系模的联系可读松村英之\cite[Theorem 25.2, pp. 194--195]{Matsumura1989CommutativeRingTheory}；本章借用的局部平坦性定理及正则性比较，分别见同书\cite[Theorem 28.9, p. 222; Theorem 30.3 and Remark 2, pp. 233--235]{Matsumura1989CommutativeRingTheory}。

% !TeX root = ../../Book3.tex
% 纲要标题：### 23.1 Hensel 引理与 Hensel 化
% 纲要位置：工作记录/计划纲要.md:2703

\section{Hensel 引理与 Hensel 化}
\label{b3:sec:2301}

模去极大理想后的根或因式分解，只有满足可逆的线性修正条件才可逐阶提升。完备性保证这些修正有极限，Hensel 化则只加入使相应代数提升成立所需的元素。

% 纲要内容（仅作写作计划，尚非教材正文）：
% * 单根提升、互素因式提升与多变量 Jacobian 隐函数形式；有限精度 Hensel 算法
% * 完备分离局部环的 Hensel 性
% * Hensel 化的泛性质；采用 Nagata 的构造与有限单根邻域描述，直接展开有限阶商、忠实平坦性及 Noether 性
% * Hensel 化与完备化的区别

\subsection{单根提升与因式分解提升}
\label{b3:subsec:2301:01}

设 \((R,\mathfrak m)\) 是局部环，\(k=R/\mathfrak m\)。把多项式的系数约化到 \(k\)，常会得到一个容易求解的方程。问题是：约化方程的解，能否成为原方程的解的第一项？如果约化根是单根，导数便可以用来修正误差；如果它是重根，误差未必能消去，也可能有不止一种消去方法。

\begin{theorem}[单根的 Hensel 提升]\label{b3:thm:hensel-simple-root}
设 \((R,\mathfrak m)\) 为交换含幺局部环，且自然映射 \(R\to\varprojlim_n R/\mathfrak m^n\) 是同构。设 \(f\in R[T]\)，\(\bar a\in k=R/\mathfrak m\) 满足
\[
\bar f(\bar a)=0,\qquad \bar f'(\bar a)\ne0.
\]
则存在唯一的 \(a\in R\)，使 \(f(a)=0\) 且 \(a\bmod\mathfrak m=\bar a\)。这里不要求 \(f\) 首一。
\end{theorem}
\begin{proof}
任取 \(\bar a\) 的提升 \(a_0\)。因为 \(f'(a_0)\) 模 \(\mathfrak m\) 非零，它是单位。递归令
\[
a_{j+1}=a_j-\frac{f(a_j)}{f'(a_j)}.
\]
每次修正都在 \(\mathfrak m\) 中，故所有 \(a_j\) 具有相同的剩余类，\(f'(a_j)\) 始终可逆。多项式恒等式
\(f(u+v)=f(u)+f'(u)\cdot v+v^2\cdot h(u,v)\) 给出：若 \(f(a_j)\in\mathfrak m^{2^j}\)，则
\[
a_{j+1}-a_j\in\mathfrak m^{2^j},\qquad
f(a_{j+1})\in\mathfrak m^{2^{j+1}}.
\]
初始误差属于 \(\mathfrak m\)，所以归纳成立。由完备性，\((a_j)\) 收敛到 \(a\in R\)。有限次加法与乘法对 \(\mathfrak m\) 进拓扑连续，故 \(f(a)\in\bigcap_n\mathfrak m^n=0\)，并且 \(a\) 仍提升 \(\bar a\)。

若 \(a,b\) 都是这样的根，则
\[
0=f(a)-f(b)=(a-b)\bigl(f'(b)+(a-b)c\bigr)
\]
其中 \(c\in R\)。括号内的元素模 \(\mathfrak m\) 等于 \(\bar f'(\bar a)\)，所以是单位，因而 \(a=b\)。这个唯一性论证本身不需要完备性。
\end{proof}

定理的递推式也能计算指定精度的根。为明确每一步都可在有限时间内完成，以下取整数系数和 \(p\) 进精度；输出是模 \(p^N\) 的剩余类，而不是在有限步内写出整个 \(p\) 进整数。

\begin{algorithm}[htbp]
\caption{Hensel 单根有限精度提升}
\label{b3:alg:hensel-root-finite-precision}
\begin{algorithmsteps}{素数 \(p\)、\(f\in\mathbb Z[T]\)、正整数 \(N\)，以及 \(a_1\in\{0,\ldots,p-1\}\)，满足 \(f(a_1)\equiv0\pmod p\)、\(f'(a_1)\not\equiv0\pmod p\)。}{唯一的 \(a_N\in\mathbb Z/p^N\mathbb Z\)，满足 \(f(a_N)=0\) 且 \(a_N\equiv a_1\pmod p\)。}
\State 置 \(a=a_1\)、\(m=1\)。
\While{\(m<N\)}
\State 置 \(t=\min\{2m,N\}\)，用扩展 Euclid 算法求 \(u\equiv f'(a)^{-1}\pmod{p^t}\)。
\State 计算 \(a\equiv a-f(a)u\pmod{p^t}\)，取任一整数代表，并置 \(m=t\)。
\EndWhile
\State \Return \(a\bmod p^N\)。
\end{algorithmsteps}
\end{algorithm}

\cref{b3:alg:hensel-root-finite-precision}的每轮满足 \(f(a)\equiv0\pmod{p^m}\) 且 \(a\equiv a_1\pmod p\)；导数始终与 \(p\) 互素，所以模 \(p^t\) 的逆元可用整数 Euclid 运算求出。上面证明中的误差平方估计表明更新后 \(f(a)\equiv0\pmod{p^t}\)。精度指数每轮加倍，最后一轮可截在 \(N\)，故至多经过 \(\lceil\log_2N\rceil\) 轮。输出的唯一性由定理的单根结论给出。

\begin{example}\label{b3:exm:hensel-root-five-adic}
在 \(\mathbb Z_5\) 中，\(T^2+1\) 的约化根 \(2\) 是单根。故恰有一个根同余于 \(2\pmod5\)。写作 \(2+5c\)，模 \(25\) 的方程为 \(5+20c\equiv0\)，得 \(c\equiv1\pmod5\)，所以这个根模 \(25\) 等于 \(7\)。从约化根 \(3\) 得到另一个根，它是前一个的负数。这里是在指定剩余类内唯一，并不是说多项式在环中只能有一个根。
\end{example}

\begin{definition}\label{b3:def:henselian-local-ring}
设 \((R,\mathfrak m)\) 为交换含幺局部环，记 \(k=R/\mathfrak m\)。若对每个首一 \(f\in R[T]\) 及每个分解
\[
\bar f=\bar g\bar h\quad\text{于 }k[T],
\]
其中 \(\bar g,\bar h\) 首一且互素，都存在首一 \(g,h\in R[T]\)，满足
\(f=gh\)、\(g\bmod\mathfrak m=\bar g\)、\(h\bmod\mathfrak m=\bar h\)，并且次数分别等于 \(\deg\bar g,\deg\bar h\)，则称 \(R\) 为 \term[Hensel-huan]{Hensel 环}[Henselian local ring]，也称 \(R\) 具有 Hensel 性。
\end{definition}

在这个定义中取 \(\bar g=T-\bar a\)，便得到首一多项式的单根提升：另一个因子与它互素，恰好表示 \(\bar f'(\bar a)\ne0\)。因式提升还可以保留一整组根，而不要求它们分别属于剩余域。下面直接证明完备性足以保证这种较一般的提升。

单变量的导数条件推广到方程组时，应要求 Jacobian 方阵可逆。下面的 Hensel 隐函数形式还保留相对于给定理想的精度。

\begin{lemma}[Hensel 隐函数形式]\label{b3:lem:henselian-jacobian-lifting}
设 \((A,\mathfrak m)\) 为交换 Noether Hensel 局部环，\(I\subseteq\mathfrak m\) 为理想，\(h\ge1\)，\(F_1,\ldots,F_h\in A[U_1,\ldots,U_h]\)。若 \(b\in A^h\) 满足
\[
F(b)\in IA^h,
\qquad
\det\left(\frac{\partial F_i}{\partial U_j}(b)\right)\in A^\times,
\]
则唯一存在 \(a\in A^h\)，使 \(F(a)=0\) 且 \(a-b\in IA^h\)。
\end{lemma}
\begin{proof}
存在性采用 Hensel 性的多变量隐函数判据：约化的零点处 Jacobian 可逆时，该零点能够提升；此判据与因式提升定义等价，而且极大理想可换成任意真理想，见 Artin\cite[(1.9), pp. 25--26]{Artin1969AlgebraicApproximation}。应用于约化点 \(b\bmod I\)，便得到所需零点。

若 \(a,a'\) 都满足条件，多项式差商给出一个 \(h\times h\) 矩阵 \(C(a,a')\)，使
\[
F(a)-F(a')=C(a,a')(a-a'),
\qquad
C(a,a')\equiv J_F(b)\pmod{\mathfrak m}.
\]
第二式使用 \(a\equiv a'\equiv b\pmod{\mathfrak m}\)。故 \(\det C(a,a')\) 为单位，矩阵可逆，而左边为零，因而 \(a=a'\)。
\end{proof}

\subsection{完备局部环的 Hensel 性}
\label{b3:subsec:2301:02}

\begin{theorem}[互素因式的 Hensel 提升]\label{b3:thm:hensel-factor-lifting}
设 \((R,\mathfrak m)\) 为交换含幺局部环，记 \(k=R/\mathfrak m\)，并且 \(R\cong\varprojlim_n R/\mathfrak m^n\)。设 \(f\in R[T]\) 首一，且
\(\bar f=\bar g\bar h\)，其中 \(\bar g,\bar h\in k[T]\) 首一互素，次数分别为 \(r,s\)。则存在唯一的首一分解 \(f=gh\)，其中 \(g,h\) 具有上述次数和约化像。特别地，每个完备分离局部环都是 Hensel 环。
\end{theorem}
\begin{proof}
可假设 \(r,s>0\)，因为零次的首一因子只能是 \(1\)。先在 \(k\) 上考虑线性映射
\[
L:k[T]_{<r}\oplus k[T]_{<s}\longrightarrow k[T]_{<r+s},
\qquad (u,v)\longmapsto\bar h u+\bar g v.
\]
若像为零，由 \(\bar g\mid\bar h u\) 及互素性可得 \(\bar g\mid u\)。次数限制迫使 \(u=0\)，继而 \(v=0\)。两端维数均为 \(r+s\)，故 \(L\) 为同构。这正是每一阶修正都可以且只能有一种方式完成的原因。

任取首一提升 \(g_1,h_1\)，次数固定为 \(r,s\)，于是 \(f-g_1h_1\in\mathfrak m R[T]\)。设已构造 \(g_n,h_n\)，使误差属于 \(\mathfrak m^nR[T]\)。在 \(\mathfrak m^n/\mathfrak m^{n+1}\) 上对 \(L\) 作张量积，仍为同构，故存在系数属于 \(\mathfrak m^n\) 的多项式 \(u_n,v_n\)，次数分别小于 \(r,s\)，使
\[
f-g_nh_n\equiv h_nu_n+g_nv_n\pmod{\mathfrak m^{n+1}R[T]}.
\]
令 \(g_{n+1}=g_n+u_n\)、\(h_{n+1}=h_n+v_n\)。乘积中另外出现的 \(u_nv_n\) 属于 \(\mathfrak m^{2n}R[T]\subseteq\mathfrak m^{n+1}R[T]\)，故新的误差提高一阶。修正不改变首项和次数。

固定次数使每个系数各自成为 Cauchy 序列。逐个取极限得到首一 \(g,h\)；分离性使 \(f-gh=0\)。若还有另一个分解 \(g'h'\)，起初 \(g'-g,h'-h\) 的系数在 \(\mathfrak m\) 中。假设已经在 \(\mathfrak m^n\) 中，则模 \(\mathfrak m^{n+1}\) 有
\[
\bar h(g'-g)+\bar g(h'-h)=0.
\]
由同一个 \(L\) 的单射性，两差的系数都在 \(\mathfrak m^{n+1}\) 中。归纳并使用分离性可得 \(g'=g,h'=h\)。
\end{proof}

上述证明在整数系数、指定 \(p\) 进精度时也给出有限算法。初始的模 \(p\) 互素分解须作为输入；求得高精度因子后，并不能据此省去模 \(p\) 因式分解这一步。

\begin{algorithm}[htbp]
\caption{Hensel 互素因式有限精度提升}
\label{b3:alg:hensel-factor-finite-precision}
\begin{algorithmsteps}{素数 \(p\)、正整数 \(N\)、首一 \(f\in\mathbb Z[T]\)，及首一互素的非常数多项式 \(\bar g,\bar h\in\mathbb F_p[T]\)，满足 \(f\bmod p=\bar g\bar h\)。}{唯一的首一 \(g_N,h_N\in(\mathbb Z/p^N\mathbb Z)[T]\)，其次数及模 \(p\) 约化分别与 \(\bar g,\bar h\) 相同，且 \(f\equiv g_Nh_N\pmod{p^N}\)。}
\State 记 \(r=\deg\bar g\)、\(s=\deg\bar h\)，任取相应次数的首一整数系数提升 \(g,h\)。
\For{\(m=1,\ldots,N-1\)}
\State 按系数计算 \(E=(f-gh)/p^m\bmod p\in\mathbb F_p[T]_{<r+s}\)。
\State 在 \(\mathbb F_p\) 上解线性方程 \(\bar h u+\bar g v=E\)，其中 \(\deg u<r\)、\(\deg v<s\)。
\State 任取 \(u,v\) 的整数系数提升 \(\tilde u,\tilde v\)，置 \(g\gets g+p^m\tilde u\)、\(h\gets h+p^m\tilde v\)，再将系数约化模 \(p^{m+1}\)。
\EndFor
\State \Return \(g\bmod p^N,\ h\bmod p^N\)。
\end{algorithmsteps}
\end{algorithm}

\cref{b3:alg:hensel-factor-finite-precision}每轮解的是证明中同一个可逆的 \(r+s\) 阶线性方程组；交叉修正项含 \(p^{2m}\)，所以新误差至少提高到模 \(p^{m+1}\) 为零。\(N=1\) 时无需循环，直接返回输入因子的模 \(p\) 约化。同样的因式提升可参阅松村英之的《Commutative Ring Theory》\cite[Theorem 8.3, p. 58]{Matsumura1989CommutativeRingTheory}。单根情形的 Newton 修正每次至少使误差的理想次幂加倍，而这里逐阶求解固定的线性方程组；两者都依靠一次近似处的可逆性。

\ProseParagraphBreak
互素条件不能省略。例如 \(T^2-5\) 在 \(\mathbb Z_5\) 中约化为 \(T\cdot T\)，但若有根 \(a\)，则 \(2v_5(a)=1\)，不可能成立。因此这个约化分解不能提升为两个首一一次因子。完备性保证误差修正的极限存在，却不能代替求解修正方程所需的可逆条件。

\subsection{Hensel 化与完备化}
\label{b3:subsec:2301:03}

完备化同时加入所有相容的 \(\mathfrak m\) 进近似。若只希望使上述互素因式能够提升，就不必加入这么多元素。描述这个较小扩张的合适方式，是规定它到其他 Hensel 局部环的映射。

\begin{definition}\label{b3:def:henselization}
设 \((R,\mathfrak m)\) 为交换含幺局部环，\(\iota:R\to R^{\mathrm h}\) 为到 Hensel 局部环的局部同态。若对任意 Hensel 局部环 \(S\) 和局部同态 \(u:R\to S\)，存在唯一的局部同态 \(v:R^{\mathrm h}\to S\)，使 \(v\iota=u\)，则称 \(\iota:R\to R^{\mathrm h}\) 为 \(R\) 的 \term[Hensel-hua]{Hensel 化}[Henselization]。相应的唯一分解为
\[
\begin{tikzcd}
R \arrow[r,"\iota"] \arrow[dr,"u"'] & R^{\mathrm h} \arrow[d,dashed,"v"] \\
& S .
\end{tikzcd}
\]
\end{definition}

这一定义立即给出唯一性：两个具有该性质的对象之间各有唯一的相容局部同态，复合由唯一性等于恒等映射。存在性则不是形式推论。永田在《Local Rings》中构造 Hensel 化，并证明其 Hensel 性、泛性质及有限邻域描述：任意有限组元素可放进一个带有指定剩余根的单根邻域\cite[\S43, (43.3), (43.5), (43.9), pp. 180--183]{Nagata1975LocalRings}。下面的概要采用这些构造结果，把有限邻域写成单根方程，并展开各阶商、平坦性和 Noether 性的推导。

\begin{theorem}[Hensel 化的存在]\label{b3:thm:henselization-existence}
设 \((R,\mathfrak m)\) 为交换含幺局部环，记 \(k=R/\mathfrak m\)。则 \(R\) 有 Hensel 化 \(R^{\mathrm h}\)。它的极大理想为 \(\mathfrak mR^{\mathrm h}\)，剩余域为 \(k\)，且 \(R^{\mathrm h}\) 在 \(R\) 上忠实平坦。若 \(R\) 为 Noether 环，则 \(R^{\mathrm h}\) 也为 Noether 环，而且对每个 \(n\ge1\)，自然映射
\[
R/\mathfrak m^n\longrightarrow
R^{\mathrm h}/\mathfrak m^nR^{\mathrm h}
\]
为同构。因此二者的相应完备化自然同构。
\end{theorem}

\begin{proofsketch}
把上述构造所得的有限单根邻域写成方程，它具有形式
\[
C=\bigl(R[T]/(f)\bigr)_{(\mathfrak m,\bar T)},
\qquad f\text{ 首一},\quad f(0)\in\mathfrak m,
\quad f'(0)\notin\mathfrak m,
\]
其中 \(\bar T\) 是 \(T\) 在商环中的像，指定的剩余根为零。任意指定单根都可经平移写成这个形式。模 \(\mathfrak m\) 后，\(\bar f=T\bar g\)，且 \(\bar g(0)\ne0\)；在 \((T)\) 处局部化时 \(\bar g\) 可逆，所以
\[
C/\mathfrak mC\cong
\bigl(k[T]/(T\bar g)\bigr)_{(T)}\cong k.
\]
因而 \(C\) 的极大理想为 \(\mathfrak mC\)。这些有限邻域及保持指定剩余点的局部映射组成滤过系统，其余极限为 \(R^{\mathrm h}\)。于是 \(R^{\mathrm h}\) 为局部环，极大理想为 \(\mathfrak mR^{\mathrm h}\)，剩余域仍为 \(k\)。由于 \(R[T]/(f)\) 是有限秩自由 \(R\)-模，局部化 \(C\) 在 \(R\) 上平坦，滤过余极限也平坦；这又是局部映射，故 \(R\to R^{\mathrm h}\) 忠实平坦。

构造定理保证这个余极限具有 Hensel 性。其泛性质可以从邻域的方程读出：若 \(u:R\to S\) 为到 Hensel 局部环 \(S\) 的局部同态，则 \(u(f)\) 模 \(S\) 的极大理想有单根零，故存在唯一的根 \(a\)，其剩余类为零。令 \(\bar T\mapsto a\)，便得到唯一的相容局部同态 \(C\to S\)。原局部化的每个分母模指定剩余点非零，因而在 \(S\) 中映为单位。这些同态由唯一性相容，取余极限便得到唯一的局部同态 \(R^{\mathrm h}\to S\)。

各阶商同构也可直接由单根方程得到。记 \(A_n=R/\mathfrak m^n\)。从 \(a_1=0\in A_1=k\) 出发，假设已在 \(A_r\) 中得到满足 \(f(a_r)=0\) 的根。映射 \(A_{r+1}\to A_r\) 的核
\[
K_r=\mathfrak m^r/\mathfrak m^{r+1}
\]
满足 \(K_r^2=0\)。任取 \(a_r\) 的提升 \(\tilde a_r\in A_{r+1}\)，则误差 \(f(\tilde a_r)\in K_r\)，而 \(f'(\tilde a_r)\) 是单位。Taylor 恒等式在平方零核上给出唯一的修正
\[
a_{r+1}=\tilde a_r-
\dfrac{f(\tilde a_r)}{f'(\tilde a_r)},
\qquad f(a_{r+1})=0.
\]
两个提升之差属于 \(K_r\)，根方程使该差乘以单位导数为零，所以这个根与初始提升的选择无关。归纳得到相容的根 \(a_n\in A_n\)，并由 \(\bar T\mapsto a_n\) 得到反向同态
\[
\beta_n:C/\mathfrak m^nC\longrightarrow A_n.
\]
这里分母仍映为单位，因为它们模 \(\mathfrak m\) 的值非零。若 \(\alpha_n:A_n\to C/\mathfrak m^nC\) 是自然同态，则 \(\beta_n\alpha_n=\mathrm{id}_{A_n}\)，因为 \(\beta_n\) 保持底环。另一方面，\(\alpha_n\beta_n\) 与 \(\mathrm{id}\) 都把 \(\bar T\) 映到模极大理想为零的根。在 \(C/\mathfrak m^nC\) 中沿其极大理想的平方零分层使用同一唯一性，二者在每一层一致，因而 \(\alpha_n\beta_n=\mathrm{id}\)。所以
\[
R/\mathfrak m^n\xrightarrow{\sim}C/\mathfrak m^nC.
\]
反向同态的唯一性使这些同构与邻域间的映射相容。对邻域取余极限，便得到
\(R/\mathfrak m^n\cong R^{\mathrm h}/\mathfrak m^nR^{\mathrm h}\)。这个商环计算本身不需要 \(R\) Noether。

现在设 \(R\) Noether。由永田的有限邻域描述，\(R^{\mathrm h}\) 中任意有限组元素均包含在一个子环 \(C=R[a]_{(\mathfrak m,a)}\) 中，其中 \(a\) 满足首一单根方程，且 \(R^{\mathrm h}\) 也是 \(C\) 的 Hensel 化\cite[(43.9), pp. 182--183]{Nagata1975LocalRings}。这里 \(C\) 是有限 \(R\) 代数 \(R[a]\) 的局部化，故 Noether；其极大理想为 \(\mathfrak mC\)。将已经得到的 Hensel 化忠实平坦性用于底环 \(C\)，可得 \(C\to R^{\mathrm h}\) 忠实平坦。若 \(b\in\bigcap_n\mathfrak m^nR^{\mathrm h}\)，取容纳它的 \(C\)，忠实平坦收缩给出
\[
b\in\bigcap_n\mathfrak m^nC=0,
\]
最后一步是 \(C\) 中的 Krull 交定理。因此 \(R^{\mathrm h}\) 分离。由各阶商同构，\(R^{\mathrm h}\) 的完备化及每个 \(C\) 的完备化都自然等于 \(\widehat R\)。

设 \(J=(b_1,\ldots,b_s)R^{\mathrm h}\)，且 \(c\in J\widehat R\cap R^{\mathrm h}\)。在共同邻域 \(C\) 中容纳 \(c,b_1,\ldots,b_s\)。由于 \(C\) Noether，它到完备化 \(\widehat R\) 的映射忠实平坦，故
\[
c\in(b_1,\ldots,b_s)\widehat R\cap C
=(b_1,\ldots,b_s)C\subseteq J.
\]
因此每个有限生成理想都等于其完备化扩张的收缩。任一有限生成理想升链在 Noether 环 \(\widehat R\) 中扩张后稳定，收缩后原升链也稳定；有限生成理想满足升链条件，便说明所有理想都有限生成，即 \(R^{\mathrm h}\) Noether。这一理想收缩论证正是永田书中 (43.10) 的证明机制\cite[p. 183]{Nagata1975LocalRings}。
\end{proofsketch}

对 Noether 局部环，由 \(\widehat R\) 的 Hensel 性及泛性质，得到相容的局部同态
\[
R\longrightarrow R^{\mathrm h}\longrightarrow\widehat R.
\]
前者只需解决 Hensel 提升问题，后者还允许任意阶的相容极限。不能仅凭两环有相同的各阶商，就断言两环相同；还必须检查它们自身是否完备。

\begin{example}\label{b3:exm:henselization-not-completion}
令 \(R=\mathbb Q[t]_{(t)}\)，则 \(\widehat R=\mathbb Q[\![t]\!]\)。多项式 \(F(X)=X^2-X-t\) 模 \(t\) 有单根 \(0\)，但在 \(\mathbb Q(t)\) 中没有根：若有根，则 \(1+4t\) 是有理函数的平方；然而它在 \(t=-1/4\) 处的零点阶数为一，平方的零点阶数必为偶数。因此 \(R\) 不是 Hensel 环。

\ProseParagraphBreak
在 Hensel 化中，这个根存在；它在完备化中的展开以
\[
-t+t^2-2t^3+5t^4+\cdots
\]
开始。这里 \(R^{\mathrm h}\) 可以看作 \(\widehat R\) 的子环：前者为 Noether 局部环，它到自身完备化的映射由 Krull 交定理为单射，而这个完备化正是 \(\widehat R\)。进一步，由永田书中 (43.9)，Hensel 化的每个元素都代数于 \(\mathbb Q(t)\)，所以 \(R^{\mathrm h}\) 是可数集；\(\mathbb Q[\![t]\!]\) 却不可数，因为它包含所有系数只取 \(0,1\) 的无穷序列。因此
\[
R\subsetneq R^{\mathrm h}\subsetneq\widehat R.
\]
这个例子同时区分了局部化、Hensel 化与完备化，而不是仅仅赋予同一环三个名称。
\end{example}

% !TeX root = ../../Book3.tex
% 纲要标题：### 23.2 平方零提升与形式性质
% 纲要位置：计划纲要.md，第 2710 行。

\section{平方零提升与形式性质}
\label{b3:sec:2302}

把误差放入平方零理想后，多项式方程的修正成为一个线性问题。提升的存在与唯一性分别约束关系模和微分模；有限展示另外保证这些条件能由有限组方程描述。

% 纲要内容（仅作写作计划，尚非教材正文）：
%
% * 平方零扩张的非空提升集合由导子加法群自由传递地作用；存在性与差的参数化分别判断
% * 形式光滑的关系模左逆判据；微分模自由仍可不形式光滑的正特征反例
% * 形式光滑、形式非分歧和形式平展的复合、基变换与目标局部化
% * 代数有限展示的作用及其与模有限展示的区别
%
% ---
%

\subsection{平方零扩张上的提升问题}
\label{b3:subsec:2302:01}

\chapref{b3:ch:14}已经把导子解释为平方零方向的变化。这里进一步追问：一个满足方程的近似点能否修正成精确点？令 \(A\to B\) 为交换含幺环同态，\(C\) 为交换 \(A\) 代数，\(J\subseteq C\) 为满足 \(J^2=0\) 的理想。给定 \(A\) 代数同态 \(u:B\to C/J\)，要求寻找 \(v:B\to C\)，使它与商映射的复合为 \(u\)。这就是本节的提升问题。

\ProseParagraphBreak
假如 \(B=A[X_1,\ldots,X_n]/(f_1,\ldots,f_s)\)，先任取 \(u(\bar X_i)\) 在 \(C\) 中的提升 \(c_i\)。这些元素未必满足关系；只知道每个误差 \(f_j(c)\) 属于 \(J\)。再以 \(h_i\in J\) 修正时，平方零条件使 Taylor 公式精确地变成
\[
f_j(c+h)=f_j(c)+\sum_i\frac{\partial f_j}{\partial X_i}(c)h_i.
\]
高次修正项全部为零。因此本来非线性的方程，变成一个关于误差的线性方程组。

\begin{lemma}\label{b3:lem:square-zero-lifts-difference}
设 \(A\to B\) 为交换含幺环同态，\(C\) 为交换 \(A\) 代数，\(J\subseteq C\) 为平方零理想，\(u:B\to C/J\) 为 \(A\) 代数同态，且至少存在一个提升 \(v:B\to C\)。由 \(u\) 使 \(J\) 成为 \(B\) 模。则全部提升恰为
\[
v+D,\qquad D\in\operatorname{Der}_A(B,J)
\cong\operatorname{Hom}_B(\Omega_{B/A},J).
\]
特别地，两个提升的差由微分模控制，而提升是否存在还需要另作判断。
导子的加法群在这个非空提升集合上自由且传递地作用；选定 \(v\) 才给出了参数化所用的原点。
\end{lemma}
\begin{proof}
\(C/J\) 对 \(J\) 的作用不依赖代表元，因为 \(J^2=0\)。若 \(v'\) 也是提升，则 \(D=v'-v\) 取值于 \(J\)，在 \(A\) 上为零；展开 \(v'(bb')\) 并删去两个 \(J\) 元素之积，得到
\[
D(bb')=u(b)D(b')+u(b')D(b).
\]
故 \(D\) 为导子。反之，同一个等式证明 \(v+D\) 保持乘法，也保持加法、单位与 \(A\) 的像，所以给出提升。最后一个同构是微分模的泛性质。
\end{proof}

存在一个提升以后，每个导子给出唯一的修正，两个提升之间也有唯一的差；在尚无提升时，这个导子模不保证方程有解。\(B=A[X]\) 时，任取一个变量像都能提升，而全部提升之间相差一个 \(J\) 中的元素。若 \(B=A[T]/(f)\) 且 \(f'(t)\) 可逆，则误差由唯一的修正 \(-f(c)/f'(c)\) 消去。前者有自由的变化方向，后者连这种方向也没有。

\subsection{形式光滑、形式非分歧与形式平展}
\label{b3:subsec:2302:02}

按照\cref{b3:def:formal-infinitesimal-properties}，对每个平方零提升问题都存在提升、至多有一个提升、恰有一个提升，分别称为形式光滑、形式非分歧、形式平展。本节沿用这些定义，不附加 Noether 条件。一般幂零理想 \(J^N=0\) 的问题可以沿 \(C/J,C/J^2,\ldots,C/J^N\) 逐层处理，因为每一步的核平方为零。

\begin{proposition}\label{b3:prop:formal-properties-stability}
交换含幺环同态的形式光滑性、形式非分歧性及形式平展性，分别在复合、任意基变换和目标环局部化下保持。
\end{proposition}
\begin{proof}
对复合 \(A\to B\to D\)，先将 \(B\to C/J\) 提升为 \(B\to C\)，从而使 \(C\) 成为 \(B\) 代数，再提升 \(D\) 的映射，便得存在性。若两个复合提升存在，形式非分歧性先使它们在 \(B\) 上相同，再使它们在 \(D\) 上相同，便得唯一性。

基变换 \(B'=A'\otimes_AB\) 的 \(A'\) 代数同态等同于保持既定 \(A'\) 作用的 \(B\) 的 \(A\) 代数同态。提升后由张量积泛性质延拓，存在性和唯一性都保留。最后，若 \(B\to S^{-1}B\)，一个在 \(C/J\) 中成为单位的元素，其任意提升在 \(C\) 中也为单位：若 \(cd=1+j\)、\(j\in J\)，则 \((1+j)^{-1}=1-j\)。因此对 \(B\) 的提升自动延拓到 \(S^{-1}B\)，且这种延拓唯一。
\end{proof}

\cref{b3:thm:formally-unramified-differentials}已证明，形式非分歧恰好等价于 \(\Omega_{B/A}=0\)。形式光滑需要处理的却是关系误差。对于 \(B=A[X_1,\ldots,X_n]/I\)，\(I/I^2\) 记录一阶关系，\(B^n\) 记录变量的修正方向，微分映射 \(\delta:I/I^2\to B^n\) 记录这些方向在各关系中的系数。若它有左逆 \(\tau\)，每个误差同态 \(\lambda:I/I^2\to J\) 就能扩为 \(\lambda\tau:B^n\to J\)，从而用变量修正消去误差。只研究余核 \(\Omega_{B/A}\) 是否投射，还不能保证这种扩张。

\begin{theorem}[形式光滑的关系模判据]\label{b3:thm:formal-smooth-conormal-splitting}\label{b3:thm:free-differentials-not-smooth}
设 \(A\) 为交换含幺环，\(P=A[X_1,\ldots,X_n]\)，\(I\subseteq P\) 为理想，\(B=P/I\)。令
\[
\delta:I/I^2\longrightarrow B^n,
\qquad \bar f\longmapsto
\left(\overline{\frac{\partial f}{\partial X_1}},\ldots,
\overline{\frac{\partial f}{\partial X_n}}\right).
\]
则 \(B\) 在 \(A\) 上形式光滑，当且仅当 \(\delta\) 有 \(B\) 线性的左逆。此时关系正合列
\[
0\longrightarrow I/I^2\xrightarrow{\delta}B^n
\longrightarrow\Omega_{B/A}\longrightarrow0
\]
分裂，特别地 \(\Omega_{B/A}\) 为有限生成投射 \(B\) 模。

这个后果不能反推形式光滑性。具体地，若 \(k\) 为特征二域，\(B_0=k[X]/(X^2)\)，则 \(\Omega_{B_0/k}=B_0\,\mathrm{d}\bar X\) 为秩一自由模，但 \(B_0\) 在 \(k\) 上不形式光滑，因而不光滑。
\end{theorem}
\begin{proof}
若形式光滑，将恒等映射 \(B\to B\) 提升为 \(s:B\to P/I^2\)。设 \(q:P\to P/I^2\)、\(\pi:P\to B\) 为自然映射，则
\(D=q-s\pi\) 取值于 \(I/I^2\)，并且由平方零条件是 \(A\) 导子。它在 \(f\in I\) 上取值 \(\bar f\)。令 \(\tau:B^n\to I/I^2\) 把第 \(j\) 个基向量送到 \(D(X_j)\)，则导子的多项式求导公式给出 \(\tau\delta=1\)。

反之，设有这样的 \(\tau\)。对任意 \(u:B\to C/J\)、\(J^2=0\)，任取变量像的提升，得到 \(w:P\to C\)。由 \(w(I)\subseteq J\)，限制映射给出 \(B\) 线性映射
\(\lambda:I/I^2\to J\)：乘以 \(P\) 中元素时，其作用只取决于在 \(B\) 中的像；而 \(I^2\) 映为零。设 \(h_j=(\lambda\tau)(e_j)\)，将第 \(j\) 个变量像改为 \(w(X_j)-h_j\)。对每个 \(f\in I\)，平方零 Taylor 公式使新像等于
\[
\lambda(\bar f)-\lambda\tau\delta(\bar f)=0.
\]
于是修正后的映射经过 \(B\)，构成所需提升。正合列的右半部分来自微分模的关系展示，左逆则保证左端单射及分裂。

对于最后的反例，取 \(P_0=k[X]\)、\(I_0=(X^2)\)。乘以 \(X^2\) 给出 \(B_0\cong I_0/I_0^2\)，所以关系模非零且自由秩一。但特征二使 \(\mathrm{d}(X^2)=0\)，故它到 \(B_0\,\mathrm{d}X\) 的映射为零，没有左逆；其余核 \(\Omega_{B_0/k}\) 仍是自由模。也可直接看提升失败：置 \(C=k[T]/(T^3)\)、\(J=(T^2)C\)，则 \(J^2=0\)、\(C/J\cong B_0\)。把 \(\bar X\) 送到 \(\bar T\) 的恒等识别若有提升，就须把 \(\bar X\) 送到 \(\bar T+h\)、\(h\in J\)；其平方却恒为非零的 \(\bar T^2\)，不能满足关系。
\end{proof}

这一证明的核心是由提升 \(B\to P/I^2\) 得到关系模的分裂；可对照松村英之\cite[Theorem 25.2, pp. 194--195]{Matsumura1989CommutativeRingTheory}。\(\Omega_{B/A}\) 投射只保证右端的满射 \(B^n\to\Omega_{B/A}\) 分裂，关系模判据还要求左边的 \(\delta\) 单射并有左逆；特征二的反例正是在这个左端失败。同一证明适用于 \(P\) 的局部化：变量像以外的要求只是将指定元素变成单位，而平方零提升自动保持这一点。

\subsection{有限展示条件的作用}
\label{b3:subsec:2302:03}

本卷约定：有限展示且形式光滑的映射称为光滑映射；有限型且形式非分歧的映射称为非分歧映射。平展映射目前定义为有限展示、平坦且 \(\Omega_{B/A}=0\)；下一节将证明它等价于有限展示且形式平展。这里说的是交换环同态 \(A\to B\) 的有限展示，即只要求有限个代数生成元和有限条关系；它与 \(B\) 作为 \(A\) 模有限展示是不同的条件。

\begin{proposition}\label{b3:prop:finite-presentation-filtered-maps}
设 \(A\) 为交换含幺环，\(B\) 为有限展示交换 \(A\) 代数，\((C_i)\) 为有向系统的交换 \(A\) 代数，\(C=\varinjlim_i C_i\)。则自然映射
\[
\varinjlim_i\operatorname{Hom}_{A\text{-alg}}(B,C_i)
\longrightarrow\operatorname{Hom}_{A\text{-alg}}(B,C)
\]
为双射。
\end{proposition}
\begin{proof}
写 \(B=A[X_1,\ldots,X_n]/(f_1,\ldots,f_s)\)。一个到 \(C\) 的映射只需指定有限个变量像，故它们在某个共同的 \(C_i\) 中已有代表。每个 \(f_j\) 在极限中为零，因而在某个较后的阶段为零。关系也只有有限条，再取一个共同阶段便得到到该阶段的映射，证明满射性。若两个阶段的映射在极限中相同，则有限个生成元的像分别在更后阶段相同；再次取共同阶段，两映射就相同，证明单射性。
\end{proof}

这个命题中，有限个代数生成元使变量像能够落在同一个阶段，有限条关系使它们确实在较后的同一阶段定义一个同态，而再比较有限个生成元的像便控制了两个同态的相等。下一节的局部标准形还会利用有限生成的关系理想，把一点处的关系等式扩大到一个主开邻域；标准模型的有限多个系数则允许把计算放到有限型整数代数上。因此有限展示同时给出有限阶段、邻域和系数环中的有限数据，形式提升性质本身不提供这些有限性。

\begin{example}\label{b3:exm:formal-without-finite-presentation}
设 \(A\) 为交换含幺环。无限变量多项式环 \(A[X_1,X_2,\ldots]\) 在 \(A\) 上形式光滑，因为可以逐个选取变量像的提升。但当 \(A\ne0\) 时，它不是有限型 \(A\) 代数：任意有限组多项式只涉及有限个变量，不能生成另外的变量。

\ProseParagraphBreak
另一个例子是 \(\overline{\mathbb Q}/\mathbb Q\)。每个有限子扩张均可分，可取本原元并用\cref{b3:prop:standard-etale-lift}得到形式平展性。对整个代数闭包的任一平方零问题，逐个有限子扩张都有唯一提升；唯一性使这些提升在交叠处一致，从而合成整个代数闭包的唯一提升。因此此扩张形式平展，却不是有限展示：有限个代数生成元只能生成有限次扩张。有限展示不能从“每次小范围提升都唯一”中省略。
\end{example}

% !TeX root = ../../Book3.tex
% 纲要标题：### 23.3 光滑与平展
% 纲要位置：计划纲要.md，第 2717 行。

\section{光滑与平展}
\label{b3:sec:2303}

可逆 Jacobian 子式允许独立修正一部分坐标，其余坐标仍可自由变化。标准模型使这一计算成为局部方程组，而平坦性与唯一提升还需要各自的代数论证。

% 纲要内容（仅作写作计划，尚非教材正文）：
%
% * 标准光滑模型的自由变量、真实关系与 Jacobian 子式；Nakayama 排除未列出的关系
% * 采用 Noether 局部环的极大理想进形式光滑判据，核对离散提升到连续提升；以有限系数子环推广平坦性
% * 对角幂等元的模分裂及平坦性转移；有限展示下平展的等价刻画，平坦性消去额外商理想
% * 平方覆盖的函数域次数阻止 Zariski 开同构；固定系数坐标取幂的自由秩与代数闭域纤维长度，比较绝对 Frobenius 并区别点集和无穷小性质
%
% ---
%

\subsection{光滑映射的标准局部模型}
\label{b3:subsec:2303:01}

\begin{definition}\label{b3:def:standard-smooth-model}
设 \(A\) 为交换含幺环，\(n\ge0\)、\(0\le r\le n\)，并取 \(f_1,\ldots,f_r,g\in A[X_1,\ldots,X_n]\)。令
\[
B=\bigl(A[X_1,\ldots,X_n]/(f_1,\ldots,f_r)\bigr)_g.
\]
如果 Jacobian 矩阵 \((\partial f_i/\partial X_j)\) 的某个 \(r\) 阶子式 \(\Delta\) 在 \(B\) 中为单位，就称这个展示为一个\term[biaozhun-guanghua-moxing]{标准光滑模型}[standard smooth presentation]。当 \(r=n\) 时称为\term[biaozhun-pingzhan-moxing]{标准平展模型}[standard étale presentation]。\(r=0\) 时空子式取为 \(1\)，得到多项式环的局部化。
\end{definition}

可逆的 \(r\) 阶子式允许用 \(r\) 个坐标修正 \(r\) 条独立关系，余下 \(n-r\) 个坐标仍能自由变化。下一个命题将在微分模中把这个数表现为基向量的数目；\(r=n\) 时，全部变量都参与修正，不再留下无穷小变化方向。

\begin{proposition}\label{b3:prop:standard-smooth-lifting}
设 \(A\) 为交换含幺环，\(n\ge0\)、\(0\le r\le n\)，并取 \(f_1,\ldots,f_r,g\in A[X_1,\ldots,X_n]\)。令 \(B=(A[X_1,\ldots,X_n]/(f_1,\ldots,f_r))_g\)。若 Jacobian 矩阵 \((\partial f_i/\partial X_j)\) 的一个 \(r\) 阶子式在 \(B\) 中可逆，则 \(A\to B\) 光滑。若可逆子式由前 \(r\) 列组成，则 \(\Omega_{B/A}\) 自由，以 \(\mathrm{d}X_{r+1},\ldots,\mathrm{d}X_n\) 为基。特别地，\(r=n\) 时 \(A\to B\) 形式平展。
\end{proposition}
\begin{proof}
重排变量后，可设给定可逆子式由前 \(r\) 列组成。有限展示显然成立，因为逆元可增加一个生成元及关系 \(gY-1\)。对 \(u:B\to C/J\)、\(J^2=0\)，任取变量像的提升 \(c_i\)。\(g(c)\) 与 \(\Delta(c)\) 都是单位，因为它们模 \(J\) 是单位。将后 \(n-r\) 个变量保持不动，令前 \(r\) 个变量的修正 \(h_j\in J\) 解线性方程
\[
\sum_{j=1}^r\frac{\partial f_i}{\partial X_j}(c)h_j=-f_i(c),
\qquad 1\le i\le r.
\]
系数矩阵可逆，故有唯一这样的修正。平方零 Taylor 公式使全部方程精确成立，局部化所需的单位也仍是单位，故得到提升。

微分模的关系矩阵恰为这个 Jacobian 矩阵。当 \(0<r<n\) 时，将它写为 \((J_0\mid J_1)\)，其中 \(J_0\) 是可逆的前 \(r\) 列。将关系左乘 \(J_0^{-1}\)，得到
\[
\begin{pmatrix}\mathrm{d}X_1\\\vdots\\\mathrm{d}X_r\end{pmatrix}
=-J_0^{-1}J_1
\begin{pmatrix}\mathrm{d}X_{r+1}\\\vdots\\\mathrm{d}X_n\end{pmatrix}.
\]
于是前 \(r\) 个微分均由其余微分表示。若这些关系的线性组合只涉及其余微分，则其前 \(r\) 个系数都为零；归一后的左侧系数矩阵是单位阵，故这个线性组合的系数只能全为零。因此其余微分之间没有非零关系，证明所述自由性。\(r=0\) 时没有关系，全部 \(\mathrm{d}X_j\) 成为自由基；\(r=n>0\) 时可逆系数矩阵使全部微分为零。因此 \(r=n\) 时总有 \(\Omega_{B/A}=0\)，由形式非分歧判据及已证的形式光滑性，得到形式平展。
\end{proof}

标准模型不仅给出一类例子。有限展示使任意光滑代数都能在每一点附近写成这样的形式。取有限多项式展示 \(B=P/I\)，证明须从关系模的局部基选出真正的关系 \(f_i\)，再用分裂的微分映射得到可逆 Jacobian 子式，最后用 Nakayama 引理证明这些 \(f_i\) 在相应局部化中确实生成整个 \(I\)。Jacobian 检验关系的一阶独立性，最后一步才排除尚未列出的实际关系。

\begin{theorem}[光滑映射的局部标准形]\label{b3:thm:smooth-local-standard-form}
设 \(A\) 为交换含幺环，\(B\) 为有限展示交换 \(A\) 代数。则 \(A\to B\) 光滑，当且仅当对每个 \(\mathfrak q\in\operatorname{Spec}B\)，存在 \(b\notin\mathfrak q\)，使 \(B_b\) 在 \(A\) 上具有标准光滑模型。
\end{theorem}
\begin{proof}
先设 \(B\) 形式光滑。写 \(B=P/I\)，其中 \(P=A[X_1,\ldots,X_n]\)，\(I\) 有限生成。由\cref{b3:thm:formal-smooth-conormal-splitting}，\(I/I^2\) 是 \(B^n\) 的直和项，故为有限生成投射模。在给定 \(\mathfrak q\) 的某个主开邻域上，它自由；从 \(I\) 的一组有限生成元中可选取 \(f_1,\ldots,f_r\)，使其像成为这一邻域上的基。更具体地，先在剩余域上选出基，再用局部自由性及相应行列式非零，缩小邻域使它们确为基。

设 \(Q\subseteq P\) 为 \(\mathfrak q\) 的逆像。因为 \(I/I^2\to B^n\) 分裂单射，\(\mathrm{d}f_1,\ldots,\mathrm{d}f_r\) 在剩余域上仍独立，所以它们的 Jacobian 矩阵有一个 \(r\) 阶子式 \(\Delta\notin Q\)。还要证明这些 \(f_i\) 在邻域上生成 \(I\)，而不只是模 \(I^2\) 生成。令 \(F=(f_1,\ldots,f_r)\)。在局部环 \(P_Q\) 中，
\[
I_Q=F_Q+I_Q^2.
\]
有限生成模 \(N=I_Q/F_Q\) 因而满足 \(N=I_QN\)。由于 \(I_Q\subseteq QP_Q\)，Nakayama 引理给出 \(N=0\)。模 \(I/F\) 有限生成，所以它在 \(Q\) 处为零意味着存在一个 \(h\notin Q\)，使 \(I_h=F_h\)：对有限组生成元分别选择零化它们的分母，再取乘积即可。因此
\[
B_{\bar h\bar\Delta}
\cong\bigl(P/(f_1,\ldots,f_r)\bigr)_{h\Delta},
\]
这就是包含 \(\mathfrak q\) 的标准光滑邻域。这里没有假设 \(A\) 为 Noether 环；所用有限性来自 \(I\) 有限生成。

反之，设这样的邻域覆盖 \(\operatorname{Spec}B\)。仍取固定的有限展示 \(B=P/I\)。微分模 \(\Omega_{B/A}\) 有限展示，且由\cref{b3:prop:standard-smooth-lifting}在这些邻域上有限自由，故由\cref{b3:thm:finite-presentation-flat-projective}为有限生成投射模。另一方面，在每个邻域中 \(B_b\) 形式光滑，关系模判据对局部化的多项式环仍适用，故固定映射 \(I/I^2\to B^n\) 在该邻域上单射。其核处处局部化为零，所以全局单射。于是
\(0\to I/I^2\to B^n\to\Omega_{B/A}\to0\)
正合。因末项投射，此列分裂，第一箭头有左逆。再用关系模判据，得到 \(B\) 形式光滑。
\end{proof}

\subsection{平展映射的平坦性与非分歧性}
\label{b3:subsec:2303:02}

第14章对首一单变量模型，已经用商环的自由基直接证明平坦性。一般标准光滑模型不一定以这种单变量形式出现；从提升性质推出平坦性，需要一个更深的局部结果。对 Noether 局部环同态 \((A,\mathfrak m)\to(B,\mathfrak n)\)，松村英之\cite[\S28, pp. 213--215; Theorem 28.9, p. 222]{Matsumura1989CommutativeRingTheory}将 \(\mathfrak n\)-进形式光滑性刻画为平坦性与闭纤维的相应形式光滑性。该等价定理的证明较深，松村书也转引其原始证明。这里采用其中推出平坦性的方向；下面的概要核对本书的形式光滑定义如何满足其拓扑假设，所采用的局部平坦性判据本身不在此推导。

\begin{theorem}[Noether 局部形式光滑蕴含平坦]\label{b3:thm:local-formal-smooth-flat-external}
设 \((A,\mathfrak m)\to(B,\mathfrak n)\) 为交换 Noether 局部环之间的局部同态。若 \(B\) 在 \(A\) 上形式光滑，则 \(B\) 是平坦 \(A\) 模。
\end{theorem}

\begin{proofsketch}
先说明离散提升性质为什么满足松村定理的拓扑假设。令 \(C\) 为交换 \(A\) 代数，\(J\subseteq C\) 为平方零理想，并给定到离散环的连续映射 \(u:B\to C/J\)，即 \(u(\mathfrak n^v)=0\) 对某个 \(v\ge1\) 成立。本书的形式光滑性给出提升 \(w:B\to C\)，而
\[
w(\mathfrak n^v)\subseteq J,
\qquad w(\mathfrak n^{2v})\subseteq J^2=0.
\]
因此提升也连续。这正是 \(B\) 在 \(A\) 上 \(\mathfrak n\)-进形式光滑的含义。

现在记 \(k=A/\mathfrak m\)、\(B_0=B/\mathfrak mB\)、\(\mathfrak n_0=\mathfrak n/\mathfrak mB\)。松村的局部判据断言，\(B\) 在 \(A\) 上 \(\mathfrak n\)-进形式光滑，当且仅当它在 \(A\) 上平坦，且闭纤维 \(B_0\) 在 \(k\) 上 \(\mathfrak n_0\)-进形式光滑\cite[Theorem 28.9, p. 222]{Matsumura1989CommutativeRingTheory}。前一段已验证其左侧条件，故可用这一判据得到平坦性。闭纤维条件是该判据给出的另一个结论，无须另加到本定理的假设中。
\end{proofsketch}

\begin{theorem}\label{b3:thm:smooth-flat}
设 \(A\to B\) 为交换含幺环之间的光滑同态，则 \(B\) 在 \(A\) 上平坦。
\end{theorem}
\begin{proof}
标准模型只用到有限多个系数。先把这些系数放到一个 Noether 子环上，逐点应用前面的局部平坦性定理，再将得到的平坦模型基变换回 \(A\)，就能处理任意底环。

先处理标准模型 \(B=(A[X_1,\ldots,X_n]/(f_1,\ldots,f_r))_g\)，设其可逆子式为 \(\Delta\)。再使 \(\Delta\) 可逆不改变 \(B\)，故可把局部化元素写成 \(g\Delta\)。令 \(A_0\) 为由 \(\mathbb Z\) 在 \(A\) 中的像以及 \(f_i,g\) 的有限多个系数生成的子环。它是有限型 \(\mathbb Z\) 代数，因而为 Noether 环。\(\Delta\) 由 \(f_i\) 求导、取行列式得到，其系数也在 \(A_0\) 中。定义
\[
B_0=\bigl(A_0[X_1,\ldots,X_n]/(f_1,\ldots,f_r)\bigr)_{g\Delta}.
\]
这是 Noether 环；\(\Delta\) 已在此模型中显式可逆，且有限展示与局部化的基变换给出 \(B\cong A\otimes_{A_0}B_0\)。标准模型的提升证明给出 \(A_0\to B_0\) 形式光滑。对每个素理想 \(\mathfrak q\subseteq B_0\)，令 \(\mathfrak p=\mathfrak q\cap A_0\)。先沿 \(A_0\to(A_0)_{\mathfrak p}\) 作基变换，得到形式光滑的同态
\[
(A_0)_{\mathfrak p}\longrightarrow
C=(A_0)_{\mathfrak p}\otimes_{A_0}B_0.
\]
因为 \(A_0\setminus\mathfrak p\) 与 \(\mathfrak q\) 不相交，\(\mathfrak q\) 在 \(C\) 中扩张为素理想 \(\mathfrak qC\)。再对目标环局部化，得到仍然形式光滑的局部同态
\[
(A_0)_{\mathfrak p}\longrightarrow
C_{\mathfrak qC}\cong(B_0)_{\mathfrak q}.
\]
前一局部定理因此给出 \((B_0)_{\mathfrak q}\) 在 \((A_0)_{\mathfrak p}\) 上平坦；而 \((A_0)_{\mathfrak p}\) 是平坦 \(A_0\) 模，故 \((B_0)_{\mathfrak q}\) 在 \(A_0\) 上也平坦。

这足以证明 \(B_0\) 在 \(A_0\) 上平坦。事实上，对任意 \(A_0\) 模单射 \(U\to V\)，张量后映射的核是一个 \(B_0\) 模；把它在每个 \(\mathfrak q\) 处局部化，得到与平坦模 \((B_0)_{\mathfrak q}\) 张量后的核，均为零。由模的局部零检测，这个核为零。最后由平坦性在基变换下保持，\(B=A\otimes_{A_0}B_0\) 在 \(A\) 上平坦。

一般光滑代数由\cref{b3:thm:smooth-local-standard-form}在目标上被这样的标准模型覆盖。同样对张量映射的核逐个邻域作局部化，就得到全局平坦性。
\end{proof}

为了从“平坦且微分为零”得到唯一提升，还须把底环上的平坦性转到一个候选标准平展环上。对有限型 \(A\) 代数 \(C\)，微分模为零使对角理想 \(K=\ker(C\otimes_AC\to C)\) 满足 \(K/K^2=0\)。有限型保证 \(K\) 有限生成，于是 \(K=K^2\) 使它由一个幂等元生成。其互补幂等元将任意 \(C\) 模 \(M\) 从 \(C\otimes_AM\) 中分裂出来，正是下面引理转移平坦性的机制。

\ProseParagraphBreak
先在 \(C=A\times A\) 中看这个分裂。以 \(\mu:C\otimes_AC\to C\) 为乘法映射，令 \(e_1=(1,0)\)、\(e_2=(0,1)\)，并取
\[
z=e_1\otimes e_1+e_2\otimes e_2
\]
则 \(\mu(z)=1\) 使 \(m\mapsto e_1\otimes e_1m+e_2\otimes e_2m\) 成为乘法映射 \(C\otimes_AM\to M\) 的截面，而 \((c\otimes1)z=(1\otimes c)z\) 保证这个截面为 \(C\) 线性。两个等式分别完成复合为恒等与标量作用相容这两件事；下面构造一般情形中的同一个 \(z\)。

\begin{lemma}\label{b3:lem:unramified-diagonal-splitting}
设 \(A\) 为交换含幺环，\(C\) 为有限型交换 \(A\) 代数，且 \(\Omega_{C/A}=0\)。若 \(M\) 为一个 \(C\) 模并且在 \(A\) 上平坦，则 \(M\) 在 \(C\) 上也平坦。
\end{lemma}
\begin{proof}
令 \(D=C\otimes_AC\)，\(\mu:D\to C\) 为乘法，\(K=\ker\mu\)。若 \(c_1,\ldots,c_l\) 生成 \(C\)，则 \(K\) 由有限个 \(c_i\otimes1-1\otimes c_i\) 生成：模去这些差以后，两份 \(C\) 的像相同，商正是 \(C\)。映射 \(c\mapsto c\otimes1-1\otimes c\bmod K^2\) 为导子，并由其泛性质给出
\(K/K^2\cong\Omega_{C/A}\)。也可直接用这些生成元及乘积的求导关系验证这个同构。因此 \(K=K^2\)。

有限生成且幂等的理想由一个幂等元生成。为说明这一点，取生成元列 \(a\)，由 \(K=K^2\) 写 \(a=Ea\)，其中矩阵 \(E\) 的各项属于 \(K\)。伴随矩阵给出 \(\det(1-E)K=0\)。写 \(\det(1-E)=1-e\)、\(e\in K\)，则 \((1-e)e=0\)，即 \(e^2=e\)，而每个 \(x\in K\) 满足 \(x=ex\)，所以 \(K=(e)\)。令 \(z=1-e\)，则
\[
\mu(z)=1,\qquad(c\otimes1-1\otimes c)z=0\quad(c\in C).
\]
写 \(z=\sum_i u_i\otimes v_i\)。映射
\[
s:M\longrightarrow C\otimes_AM,
\qquad m\longmapsto\sum_i u_i\otimes v_i m
\]
由第二个等式为 \(C\) 线性，且乘法映射 \(C\otimes_AM\to M\) 与它的复合为恒等映射。故 \(M\) 是 \(C\otimes_AM\) 的直和项。后者由基变换在 \(C\) 上平坦，直和项也平坦，结论成立。
\end{proof}

\begin{theorem}[平展的等价刻画]\label{b3:thm:etale-equivalent-characterizations}\label{b3:thm:closed-point-unramified-not-etale}
设 \(A\) 为交换含幺环，\(B\) 为有限展示交换 \(A\) 代数。下列条件等价：
\begin{equivalences}
\item\label{b3:item:etale-flat-unramified} \(B\) 在 \(A\) 上平坦，且 \(\Omega_{B/A}=0\)，即 \(A\to B\) 平展；
\item\label{b3:item:etale-formal} \(A\to B\) 形式平展；
\item\label{b3:item:etale-standard} 每个 \(\mathfrak q\in\operatorname{Spec}B\) 都有主开邻域，在该邻域上 \(B\) 具有标准平展模型。
\end{equivalences}
其中平坦性不能仅由微分消失代替：对任意域 \(k\)，\(A=k[t]\to B=A/(t)=k\) 满足 \(\Omega_{B/A}=0\)，是有限展示的非分歧映射，但不平坦、不形式光滑，因而不平展。
\end{theorem}
\begin{proof}
先从条件\itemref{b3:item:etale-formal}出发。形式平展包含形式光滑性，由局部标准形定理得到标准光滑邻域 \(D(b)\)，其中包含给定的 \(\mathfrak q\)。形式非分歧又使 \(\Omega_{B/A}=0\)，而标准模型给出
\[
0=\Omega_{B_b/A}\cong B_b^{\,n-r}.
\]
由于 \(\mathfrak q\in D(b)\)，环 \(B_b\) 有素理想，因而 \(B_b\ne0\)。非零环上的自由模只有在秩为零时才为零，故 \(n=r\)，模型标准平展，得到条件\itemref{b3:item:etale-standard}。

再设条件\itemref{b3:item:etale-standard}成立。标准平展邻域由局部标准形定理使 \(B\) 光滑，且微分模处处为零，因而全局为零。这给出形式光滑与形式非分歧，合起来即为条件\itemref{b3:item:etale-formal}。光滑蕴含平坦又给出条件\itemref{b3:item:etale-flat-unramified}。

最后从条件\itemref{b3:item:etale-flat-unramified}构造标准平展邻域。先用微分消失选出一组 Jacobian 可逆的关系，得到一个候选标准平展环；原代数仍可能是它的商，平坦性将消去这些尚未列出的关系。写 \(B=P/I\)、\(P=A[X_1,\ldots,X_n]\)，其中 \(I\) 有限生成。给定 \(\mathfrak q\)，关系正合列和 \(\Omega_{B/A}=0\) 说明 \(I/I^2\to B^n\) 满射。模去 \(\mathfrak q\) 后，这些生成元的梯度向量张成 \(\kappa(\mathfrak q)^n\)，从中选一组基，便得 \(n\) 个关系 \(f_1,\ldots,f_n\)，使其 Jacobian 行列式满足 \(\bar\Delta\notin\mathfrak q\)。令
\[
C=\bigl(P/(f_1,\ldots,f_n)\bigr)_\Delta,
\qquad B'=B_{\bar\Delta}=C/J.
\]
这里 \(J\) 由其余有限个关系生成，\(C\) 为标准平展模型，故 \(\Omega_{C/A}=0\)。\(B'\) 在 \(A\) 上平坦，因为它是平坦 \(A\) 代数 \(B\) 的局部化。由\cref{b3:lem:unramified-diagonal-splitting}，\(B'\) 在 \(C\) 上平坦。

用平坦 \(C\) 模 \(B'\) 张量单射 \(J\to C\)，得到单射 \(J/J^2\to B'\)；这个映射把 \(j\) 送到其在 \(C/J\) 中的零像，故本身为零，只能有 \(J/J^2=0\)，即 \(J=J^2\)。将所选 \(\mathfrak q\) 扩张到 \(B'\)，令 \(\mathfrak r\subseteq C\) 为其逆像，则 \(J\subseteq\mathfrak r\)。在 \(C_{\mathfrak r}\) 中，有限生成模 \(J_{\mathfrak r}\) 满足
\(J_{\mathfrak r}=J_{\mathfrak r}^2\)，由 Nakayama 引理为零。因此存在 \(h\notin\mathfrak r\)，使 \(J_h=0\)。于是 \(B'_h=C_h\)，这是包含给定点的标准平展邻域。这样条件\itemref{b3:item:etale-flat-unramified}也推出条件\itemref{b3:item:etale-standard}。

最后的闭点反例已在\cref{b3:exm:unramified-not-etale}证明：相对导子全为零，乘 \(t\) 的单射张量 \(k\) 后成为零映射，而 \(k\to k[t]/(t^2)\) 的 \(k[t]\) 代数提升不能使 \(t\) 的两个像相容。在上面的候选模型构造中，这个例子取 \(C=A=k[t]\)、\(J=(t)\)，仍有 \(J/J^2\cong k\ne0\)。因此微分消失本身没有删掉这条关系，平坦性正是所缺的条件。
\end{proof}

这样便得到了第14章提出的定义比较。有限展示在选择有限组关系、把一点处的消失扩大到邻域，以及形成标准模型时都有实际作用；平坦性则在最后一个证明中消去了多余的商理想。只要求 \(\Omega_{B/A}=0\) 不足以做到这一点，\(k[t]\to k\) 正是已经算过的反例。

\ProseParagraphBreak
对交换含幺环同态 \(A\to B\)，上述形式性质与有限性条件的关系汇总于\cref{b3:tab:formal-geometric-finiteness}。其中的局部模型均指在 \(\operatorname{Spec}B\) 的每一点附近成立的展示。

\begin{table}[htbp]
\centering
\caption{形式性质、有限性条件与局部模型}
\label{b3:tab:formal-geometric-finiteness}
\begin{tabularx}{\textwidth}{@{}l l X@{}}
\toprule
形式性质 & 附加有限性 & 所得性质与局部模型 \\
\midrule
形式光滑 & 有限展示 & 光滑，提升存在；局部具有标准光滑模型，且必平坦 \\
形式非分歧 & 有限型 & 非分歧，提升至多唯一；\(\Omega_{B/A}=0\)，但不一定平坦 \\
形式平展 & 有限展示 & 平展，提升存在且唯一；局部具有标准平展模型 \\
\bottomrule
\end{tabularx}
\end{table}

平展一行也可以用有限展示、平坦及 \(\Omega_{B/A}=0\) 表示。这里的有限展示与有限型都是指 \(B\) 作为 \(A\) 代数的有限性条件。

\subsection{局部同胚直觉与平展的代数条件}
\label{b3:subsec:2303:03}

在复数几何中，可逆 Jacobian 容易使人联想到逆函数定理。这个直觉有助于理解“相对零维而没有无穷小变化”，却不能把平展直接定义成 Zariski 拓扑中的局部同构，更不能只检查点集映射。

\begin{example}\label{b3:exm:etale-not-zariski-isomorphism}
设 \(A=\mathbb C[s,s^{-1}]\)、\(B=\mathbb C[t,t^{-1}]\)，映射由 \(s\mapsto t^2\) 给出。写成
\(B\cong A[T]/(T^2-s)\)，其中 \(T\) 自动可逆，导数 \(2T\) 为单位。故此映射平展，且 \(B\) 为秩二自由 \(A\) 模。

\ProseParagraphBreak
然而不能在 \(\operatorname{Spec}B\)、\(\operatorname{Spec}A\) 中分别取非空 Zariski 开集，使给定谱映射的限制成为同构。这样的相容同构必诱导函数域同构，迫使 \(\mathbb C(s)\to\mathbb C(t)\)、\(s\mapsto t^2\) 为同构。但在函数域 \(\mathbb C(s)\) 的 \(s=0\) 离散赋值 \(v_0\) 下，\(v_0(s)=1\)，而非零平方的赋值必为偶数，故 \(T^2-s\) 在 \(\mathbb C(s)\) 上不可约，扩张次数为二。给定映射在函数域上已经有次数二的障碍，缩小非空开集不能消去它。这里用的是函数域上的赋值，\(s=0\) 不是 \(\operatorname{Spec}A\) 中的点；复解析拓扑中可以局部选定平方根，采用的是另一种邻域。
\end{example}

\begin{example}\label{b3:exm:frobenius-homeomorphism-not-etale}
设 \(k\) 为特征 \(p>0\) 的代数闭域，考虑 \(k[s]\to k[t]\)、\(s\mapsto t^p\)。目标以 \(1,t,\ldots,t^{p-1}\) 为自由基，所以映射有限且平坦。在素谱上，它把泛点送到泛点，并把闭点 \((t-a)\) 送到 \((s-a^p)\)；代数闭域上 \(a\mapsto a^p\) 为双射。有限整映射是闭映射，因此这个素谱映射为同胚。

\ProseParagraphBreak
这里保持 \(k\) 中系数不变，只对坐标取 \(p\) 次幂。它与对每个元素 \(r\) 都取 \(r^p\) 的绝对 Frobenius 一般不同；后者同时改变系数，且在素谱上逐点恒等。附录 C 的\secref{b3:sec:C02}将明确比较这两种环同态。

\ProseParagraphBreak
其相对微分模却为
\[
\Omega_{k[t]/k[s]}
\cong k[t]\,\mathrm{d}t/(\mathrm{d}(t^p))=k[t]\,\mathrm{d}t\ne0.
\]
故它不是非分歧映射，也不是平展映射。

\ProseParagraphBreak
进一步，在闭点 \(s=a\) 上，由代数闭性可唯一选取 \(b\in k\) 使 \(b^p=a\)。纤维代数为
\[
k[t]\otimes_{k[s]}k
\cong k[t]/(t^p-a)
=k[t]/((t-b)^p)
\cong k[\varepsilon]/(\varepsilon^p).
\]
它以 \(1,\varepsilon,\ldots,\varepsilon^{p-1}\) 为 \(k\) 基，维数为 \(p\)。这个局部 Artin 代数的剩余域仍是 \(k\)，所以其组成列长度也为 \(p\)，等于原有限自由扩张的秩。各闭纤维都只有一个底层点，却带有非零幂零元 \(\varepsilon\)；这种厚度与平坦性共存，而这份纤维代数在 \(k\) 上的微分模仍为非零的 \(k[\varepsilon]/(\varepsilon^p)\,\mathrm{d}\varepsilon\)。点集及普通拓扑看不到这部分无穷小信息。
\end{example}

因此，对给定代数映射，平展性应当由有限展示、平坦性、微分模，或由已经证明等价的标准模型来检验。只知道两个环都正则，或者两个素谱同胚，都不够。

% !TeX root = ../../Book3.tex
% 纲要标题：### 23.4 局部性质的比较与反例
% 纲要位置：计划纲要.md，第 2238 行。

\section{局部性质的比较与反例}
\label{b3:sec:2304}

比较环同态的局部性质时，必须同时写清底环、所考察的素点与纤维。尖点正规化把微分模非零、特殊纤维长度改变和原点外同构分成三个可以分别核验的计算。

% 纲要内容：
%
% * 尖点正规化在尖点处的局部计算
% * 光滑性、正则性与可分性的比较
% * 代数几何中几种拓扑的简介
% * 在特征零下令 \(A=k[t^2,t^3]\subseteq B=k[t]\)，计算 \(\Omega_{B/A}\cong B/(t)\,dt\)，故正规化在尖点处并非非分歧
% * 计算尖点纤维 \(B/(t^2,t^3)B\cong k[t]/(t^2)\)；其长度为2而双有理映射的泛秩为1，用有限生成平坦模在局部环上自由的结论证明尖点处不平坦
% * 在去掉尖点后通过 \(t=t^3/t^2\) 验证该映射为同构；分别说明“原点处非分歧失败”“原点处平坦失败”和“原点外为同构”三项结论
%
% ---
%

\subsection{尖点正规化的局部计算}
\label{b3:subsec:2304:01}

本节固定特征零域 \(k\)，考察平面尖点及其参数化
\[
A=k[t^2,t^3]\subseteq B=k[t].
\]
这是一个尤其适合比较局部性质的例子：参数环很简单，映射有限且双有理，但在尖点处既不非分歧，也不平坦。后三个小节分别计算微分模、特殊纤维和去掉尖点后的环，以区分这三项判断。

\begin{proposition}\label{b3:prop:cusp-normalization-setup}
设 \(k\) 为特征零域，\(A=k[t^2,t^3]\)、\(B=k[t]\)。则
\[
A\cong k[x,y]/(y^2-x^3),\qquad B=A+At,
\qquad\operatorname{Frac}(A)=\operatorname{Frac}(B)=k(t).
\]
并且 \(B\) 是 \(A\) 在这个分式域中的整闭包。尖点对应 \(\mathfrak m=(t^2,t^3)\subseteq A\)，其上方唯一的素理想为 \(\mathfrak n=(t)\subseteq B\)。
\end{proposition}
\begin{proof}
映射 \(k[x,y]\to B\)、\(x\mapsto t^2,y\mapsto t^3\) 的像为 \(A\)，核包含 \(y^2-x^3\)。模去此关系后，每个多项式都能写成 \(F(x)+y\cdot G(x)\)。代入后，\(F(t^2)\) 只含偶次幂，\(t^3\cdot G(t^2)\) 只含至少三次的奇次幂，不能相互抵消；所以核恰为上述主理想。

\(t\) 满足首一方程 \(T^2-t^2=0\)，且所有次数至少二的单项式都在 \(A\) 中，因此 \(B=A+At\) 为有限整扩张。又有 \(t=t^3/t^2\in\operatorname{Frac}(A)\)，故分式域相同。\(B=k[t]\) 为整闭整环；若分式域中元素整于 \(A\)，它也满足系数在 \(B\) 中的同一个首一方程，故属于 \(B\)。这证明整闭包的断言。

若 \(\mathfrak q\subseteq B\) 位于 \(\mathfrak m\) 上方，则 \(t^2\in\mathfrak q\)，从而 \(t\in\mathfrak q\)，所以 \(\mathfrak q=(t)\)。反向收缩显然为 \(\mathfrak m\)。
\end{proof}

局部环 \(A_{\mathfrak m}\) 维数为一，而 \(\mathfrak m/\mathfrak m^2\) 以 \(t^2,t^3\) 的像为基，维数为二，故不正则。相反，\(B_{\mathfrak n}=k[t]_{(t)}\) 是离散赋值环，因而正则。正规化改善了目标坐标环本身的奇点，但这种改善并不意味着原来的环同态光滑。

\subsection{光滑性、正则性与可分性}
\label{b3:subsec:2304:02}

正则性属于局部环本身：比较 \(\dim R\) 与 \(\dim_{R/\mathfrak m}\mathfrak m/\mathfrak m^2\) 就可作出判断。光滑性属于一个带有底环的同态。同一个域看作自身上的代数总是光滑，换成一个较小而不完美的基域以后却可能不光滑。

\begin{theorem}[完美基域上的比较]\label{b3:thm:perfect-field-regular-smooth}
设 \(k\) 为完美域，\(B\) 为有限型交换 \(k\) 代数，\(\mathfrak q\in\operatorname{Spec}B\)。则 \(B_{\mathfrak q}\) 为正则局部环，当且仅当存在 \(b\notin\mathfrak q\)，使 \(k\to B_b\) 光滑。
\end{theorem}

\begin{proof}
写 \(B=k[X_1,\ldots,X_n]/I\)，令 \(Q\) 为 \(\mathfrak q\) 在多项式环中的逆像，\(P=k[X_1,\ldots,X_n]_Q\)、\(\mathfrak a=QP\)、\(\kappa=\kappa(\mathfrak q)\)。多项式环的各个素点局部环正则：\cref{b3:prop:polynomial-global-bound}给出其剩余域的有限投射维数，再用\cref{b3:thm:regular-local-global-dimension}即可。设 \(h=\dim P\)，则 \(\dim_\kappa\mathfrak a/\mathfrak a^2=h\)。

先证明这里需要的可分超越基存在性。域 \(\kappa\) 在 \(k\) 上有限生成，记 \(s=\operatorname{trdeg}_k\kappa\)。特征零时，任取超越基后剩余的有限代数扩张可分。特征 \(p>0\) 时，先任取超越基 \(u_1,\ldots,u_s\)，令 \(E=k(u_1,\ldots,u_s)\)、\(m=[\kappa:E]\)。完美性使 \(E^p=k(u_1^p,\ldots,u_s^p)\)，指数都小于 \(p\) 的 \(u\) 单项式构成 \(E/E^p\) 的基：多项式按指数模 \(p\) 分组，分式则先把分母改写成 \(p\) 次幂即可。因此 \([E:E^p]=p^s\)。Frobenius 给出两扩张 \(\kappa/E\) 与 \(\kappa^p/E^p\) 的同构，故
\[
[\kappa:\kappa^p]
=\frac{[\kappa:E][E:E^p]}{[\kappa^p:E^p]}
=p^s.
\]
对有限扩张 \(\kappa/\kappa^p\) 应用\cref{b3:lem:exponent-one-differentials}的逐次选元构造，得到 \(t_1,\ldots,t_s\)，使 \(\kappa=\kappa^p(t_1,\ldots,t_s)\)，且所有 \(t^\alpha\)、\(0\le\alpha_i<p\)，构成 \(\kappa/\kappa^p\) 的基。这些 \(t_i\) 对 \(k\) 代数独立。否则取总次数最小的非零关系 \(F(t)=0\)；因 \(k\) 完美，可以写成
\[
F(T)=\sum_{0\le\alpha_i<p}T^\alpha G_\alpha(T)^p,
\qquad G_\alpha\in k[T_1,\ldots,T_s].
\]
上述基的线性无关性使各 \(G_\alpha(t)=0\)，而至少一个 \(G_\alpha\) 非零，其次数严格小于 \(F\)，矛盾。于是 \(t_i\) 为超越基，\(\kappa/k(t_1,\ldots,t_s)\) 有限；等式 \(\kappa=k(t_1,\ldots,t_s)\kappa^p\) 再由同一引理给出该扩张的微分模为零，\cref{b3:thm:finite-field-differentials-separable}遂保证它可分。\(s=0\) 时选空基，同一计算仍适用。

现在证明微分映射
\[
\mathfrak a/\mathfrak a^2\longrightarrow
\Omega_{P/k}\otimes_P\kappa
\]
单射。取刚得到的可分超越基，由本原元素定理写 \(\kappa=k(t_1,\ldots,t_s)[z]\)，其中 \(z\) 的最小多项式可分。将 \(t_i\) 提升到 \(P/\mathfrak a^2\)，其非零有理函数都为单位，因而得到 \(k(t_1,\ldots,t_s)\) 的提升。对 \(z\) 的任一提升，用最小多项式的非零导数作一次平方零 Newton 修正，就使其精确满足方程。由此得到 \(\kappa\to P/\mathfrak a^2\) 的 \(k\) 代数截面。于是
\(P/\mathfrak a^2\cong\kappa\oplus(\mathfrak a/\mathfrak a^2)\)，后一个直和项平方为零；投影到该项是 \(k\) 导子，并在 \(\mathfrak a/\mathfrak a^2\) 上为恒等。这给出微分映射的左逆，证明单射。

若 \(R=B_{\mathfrak q}=P/IP\) 正则，令 \(e=\dim R\)、\(c=h-e\)。余切空间公式
\(\mathfrak m_R/\mathfrak m_R^2\cong\mathfrak a/(\mathfrak a^2+IP)\)
使 \(IP\) 在 \(\mathfrak a/\mathfrak a^2\) 中的像有维数 \(c\)。选 \(f_1,\ldots,f_c\in I\) 使其像为基；在正则局部环 \(P\) 中，这些元素可补成一组正则参数。因此 \(P/(f_1,\ldots,f_c)P\) 正则，维数为 \(e\)，并满射到同维数的 \(R\)。这个满射的核为零：源为正则局部整环，任何非零素理想都使其商的维数至少降低一。故 \(IP=(f_1,\ldots,f_c)P\)。

刚证的单射又使 \(\mathrm{d}f_i\) 在 \(\Omega_{P/k}\otimes_P\kappa=\kappa^n\) 中独立，故其 Jacobian 有一个 \(c\) 阶子式 \(\Delta\notin Q\)。理想 \(I/(f_1,\ldots,f_c)\) 有限生成，且在 \(Q\) 处为零；可选择 \(g\notin Q\) 使它在倒置 \(g\) 后为零。于是 \(B_{\bar g\bar\Delta}\) 是标准光滑模型，由\cref{b3:prop:standard-smooth-lifting}光滑。

反过来，光滑代数在给定点附近有标准模型。其方程的微分在 \(\kappa^n\) 中线性独立，由同一个单射，它们在 \(\mathfrak a/\mathfrak a^2\) 中也独立，故构成正则参数的一部分。商去这些参数仍为正则局部环，得到 \(B_{\mathfrak q}\) 正则。

上述可分剩余域下的余切空间比较，也见 Stacks Project, Tags 00TU、00TV（\url{https://stacks.math.columbia.edu/tag/00TU}，\url{https://stacks.math.columbia.edu/tag/00TV}）；与几何正则性的比较见松村英之\cite[Theorem 28.7, pp. 219--220; Theorem 30.3 and Remark 2, pp. 233--235]{Matsumura1989CommutativeRingTheory}。
\end{proof}

\begin{example}\label{b3:exm:regular-not-smooth-imperfect}
设 \(k=\mathbb F_p(s)\)，\(L=k[u]/(u^p-s)\)。多项式 \(T^p-s\) 不可约，故 \(L\) 是域，自身作为零维局部环正则。但是 \(L\) 在 \(k\) 上不光滑。为直接看到提升失败，令
\[
C=k[T]/\bigl((T^p-s)^2\bigr),\qquad
J=(T^p-s)C.
\]
则 \(J^2=0\)、\(C/J\cong L\)。若恒等映射 \(L\to C/J\) 可以提升，\(u\) 的像必为 \(T+h\)、\(h\in J\)。然而特征 \(p\) 下
\((T+h)^p-s=T^p-s+h^p=T^p-s\ne0\)，因为 \(h^p=0\)。关系不成立，矛盾。

\ProseParagraphBreak
这个例子已在第14章出现；这里的新比较是 \(L\) 的正则性与 \(k\to L\) 的不光滑性可以同时成立。基变换后
\(L\otimes_kL\cong L[\varepsilon]/(\varepsilon^p)\)
也解释了差别：正则的原环在不可分基变换后产生幂零元。
\end{example}

上述完美性条件不能简单改写成“光滑点的剩余域必须可分于基域”。例如在 \(k=\mathbb F_p(s)\) 上，多项式代数 \(k[T]\) 光滑，但闭点 \((T^p-s)\) 的剩余域正是这个不可分扩张。光滑性研究整个代数映射的无穷小行为，并不要求每个闭点本身都是一个平展的 \(k\) 点。

\subsection{代数几何中的不同拓扑}
\label{b3:subsec:2304:03}

Zariski 邻域只使用环的局部化；平展映射则容许加入在给定点附近唯一可提升的代数解。为了使这些不同的局部研究方式各有相应的覆盖概念，代数几何使用比普通开集覆盖更一般的约定。这里只介绍仿射环所描述的三种覆盖，不展开层与上同调理论。

\begin{definition}\label{b3:def:affine-etale-fppf-covers}
设 \(A\) 为交换含幺环，\((A\to B_i)_{i\in I}\) 为一族交换 \(A\) 代数。若每个 \(\mathfrak p\in\operatorname{Spec}A\) 都是某个 \(B_i\) 中素理想的收缩，就称这一族在素谱上联合满射。
\begin{enumerate}
\item 对 \(f_i\in A\)（\(i\in I\)），若形如 \(A\to A_{f_i}\) 的族联合满射，则称之为仿射的 Zariski 基本开覆盖；其条件等价于 \((f_i\mid i\in I)=A\)。
\item 若每个 \(A\to B_i\) 平展且这一族联合满射，则称这一族为一个\term[pingzhan-fugai]{平展覆盖}[étale covering]。
\item 若每个 \(A\to B_i\) 平坦、有限展示且这一族联合满射，则称这一族为一个 \term[fppf-fugai]{fppf 覆盖}[fppf covering]。
\end{enumerate}
\end{definition}

字母 fppf 来自法语 \emph{fidèlement plat et de présentation finie}。对单个映射，平坦并在素谱上满射等价于忠实平坦，所以这个名称与第12章的代数条件相符。基本开覆盖是平展覆盖，平展覆盖也是 fppf 覆盖；但这些覆盖的每一项不必是原素谱的一个开子集。

\ProseParagraphBreak
例如平方映射 \(\mathbb C[s,s^{-1}]\to\mathbb C[t,t^{-1}]\) 是单个平展覆盖，因为它有限自由且素谱满射，却不是基本开嵌入。特征 \(p\) 下的映射 \(k[s]\to k[t]\)、\(s\mapsto t^p\) 则是 fppf 覆盖，而不是平展覆盖。因此在这些覆盖中说“局部存在解”，可能需要先扩张环；它不等同于在原来的 Zariski 开集中寻找解。

\ProseParagraphBreak
一般的平展拓扑与 fppf 拓扑把这样的覆盖作为基本资料，交叠由张量积描述。第12章的忠实平坦下降说明，扩张环以后得到的模及同构，只有满足交叠上的相容条件才来自原环。这也是为什么采用更丰富的覆盖后，仍须记录如何把局部对象拼回去。

\subsection{尖点正规化的微分模与非分歧性}
\label{b3:subsec:2304:04}

\begin{proposition}\label{b3:prop:cusp-relative-differentials}
设 \(k\) 为特征零域，\(A=k[t^2,t^3]\subseteq B=k[t]\)。则作为 \(B\) 模，
\[
\Omega_{B/A}\cong B/(t)\,\mathrm{d}t.
\]
因此 \(A\to B\) 在 \(\mathfrak n=(t)\) 处不是非分歧映射。
\end{proposition}
\begin{proof}
对 \(k\to A\to B\) 的微分传递正合列为
\[
B\otimes_A\Omega_{A/k}\longrightarrow\Omega_{B/k}
\longrightarrow\Omega_{B/A}\longrightarrow0.
\]
\(\Omega_{B/k}=B\,\mathrm{d}t\)，而 \(A\) 由 \(t^2,t^3\) 在 \(k\) 上生成，所以第一箭头的像由
\(\mathrm{d}(t^2)=2t\,\mathrm{d}t\)、\(\mathrm{d}(t^3)=3t^2\,\mathrm{d}t\) 生成。特征零使 \(2\) 可逆，故像为 \((t)\,\mathrm{d}t\)，得到所述商模。在 \((t)\) 处局部化后，这个模仍同构于 \(k\,\mathrm{d}t\)，不为零，所以该点不满足形式非分歧判据，也不满足非分歧条件。
\end{proof}

这一计算说明参数化在尖点处有一个不能被底环中的函数检测的一阶方向：对平方零参数 \(\varepsilon\)，代入 \(t=\varepsilon\) 后，\(t^2,t^3\) 都变成零，但 \(t\) 本身未必为零。因此底环只看到同一个点，上方却存在不同的无穷小变化。

\subsection{尖点纤维的长度与非平坦性}
\label{b3:subsec:2304:05}

\begin{proposition}\label{b3:prop:cusp-fiber-nonflat}
设 \(k\) 为特征零域，\(A=k[t^2,t^3]\subseteq B=k[t]\)，\(\mathfrak m=(t^2,t^3)\)、\(\mathfrak n=(t)\)。则特殊纤维为
\[
B\otimes_A A/\mathfrak m
\cong B/(t^2,t^3)B
\cong k[t]/(t^2),
\]
其长度为二；而一般纤维为 \(k(t)\)，维数为一。因此局部映射 \(A_{\mathfrak m}\to B_{\mathfrak n}\) 不平坦。
\end{proposition}
\begin{proof}
特殊纤维的同构直接由商与张量积相容得到；\(1,t\) 构成其 \(k\) 基，且唯一极大理想 \((t)\) 的平方为零，所以长度为二。令 \(K=\operatorname{Frac}(A)=k(t)\)。由于 \(t=t^3/t^2\)，有 \(B\otimes_AK=K\)，故一般秩为一。

注意 \(B\otimes_AA_{\mathfrak m}\) 是有限 \(A_{\mathfrak m}\) 代数，其极大理想只能位于 \(\mathfrak m\) 上方，而那里唯一的素理想是 \(\mathfrak n\)。因此这个有限代数是局部环，并等于 \(B_{\mathfrak n}\)：\(B\) 中 \(\mathfrak n\) 外元素在它里面全为单位。

若 \(B_{\mathfrak n}\) 在 \(A_{\mathfrak m}\) 上平坦，则它作为有限生成平坦模，由\cref{b3:lem:finite-flat-local-free}必须自由，设秩为 \(r\)。与 \(K\) 张量给出 \(r=1\)。再与剩余域 \(k\) 张量，便应得到一维 \(k\) 空间，却与刚算出的二维纤维矛盾。因此局部映射不平坦。
\end{proof}

这里比较纤维长度时，有限生成平坦模在局部环上自由是必不可少的论证。单说“纤维大小发生变化”不足以证明一般映射不平坦；必须把这种变化与平坦性在当前有限条件下的具体后果联系起来。

\subsection{尖点外的同构与局部性质的比较}
\label{b3:subsec:2304:06}

\begin{proposition}\label{b3:prop:cusp-isomorphism-off-origin}
设 \(k\) 为特征零域，\(A=k[t^2,t^3]\subseteq B=k[t]\)。则
\[
A_{t^2}=B_{t^2}=k[t,t^{-1}].
\]
因此正规化映射在去掉尖点后为同构；尖点之外的局部映射都平展且光滑。
\end{proposition}
\begin{proof}
在 \(A_{t^2}\) 中，\(t=t^3/t^2\)，而 \(t\) 的逆元为 \(t/t^2\)，故 \(k[t,t^{-1}]\subseteq A_{t^2}\)。反向包含显然，\(B_{t^2}\) 也等于这个环。又若 \(A\) 的素理想包含 \(t^2\)，则由 \((t^3)^2=(t^2)^3\) 也包含 \(t^3\)，从而等于极大理想 \(\mathfrak m\)。所以 \(D(t^2)\) 正是尖点的补集。环同构保持全部平方零提升且作为模自由秩一，故局部映射平展和光滑。
\end{proof}

现在可以分别读出三个结论。在原点，\(\Omega_{B/A}\) 不为零，所以非分歧性失败；特殊纤维长度二而一般秩一，所以平坦性失败；在原点外，显式公式 \(t=t^3/t^2\) 给出同构。这些结论来自三个不同的计算，不能以“正规化去掉了奇点”一句话相互替代。

\ProseParagraphBreak
另一方面，\(B_{(t)}\) 本身正则，\(k\to B\) 也光滑，但 \(A_{\mathfrak m}\to B_{(t)}\) 不光滑，因为光滑必平坦。这个例子再次说明：讨论光滑时必须写清底环；讨论非分歧时必须保留无穷小信息；讨论平坦时则必须考察模的关系是否在张量之后仍被正确保留。


\begin{chaptersummary}
Hensel 提升首先是一个误差修正过程：单根条件使一元导数可逆，互素条件使控制因式修正的线性映射可逆，完备性再使各阶近似收敛。Hensel 化用泛性质只加入保证这些代数提升所需的元素；即使它与原环有相同的各阶商，也未必等于完备化。

\ProseParagraphBreak
平方零提升把关系的误差变成线性问题。形式光滑使关系模嵌入自由微分模并分裂；有限展示进一步把这个分裂转化为邻域上的标准方程组。一般平坦性的证明还采用了明确列出的 Noether 局部定理。加上微分模消失以后，平展可以等价地描述为唯一提升，或局部标准平展模型。

\ProseParagraphBreak
尖点正规化把这些区别落实为计算：相对微分模为 \(k\,\mathrm{d}t\)，特殊纤维为长度二的 \(k[t]/(t^2)\)，原点外则由 \(t=t^3/t^2\) 得到同构。正则的参数环、双有理性以及点集上的行为，都不能代替对给定环同态的光滑、非分歧和平坦条件逐项检验。
\end{chaptersummary}

\BookExercises

\begin{optionalexercise}\label{b3:ex:23-five-adic-root}
在 \(\mathbb Z_5\) 中求 \(T^2+1\) 同余于 \(2\pmod5\) 的根模 \(125\) 的值，并解释所得剩余类为何能继续唯一提升。
\end{optionalexercise}
\begin{solution}
先写 \(a=2+5c\)，模 \(25\) 得 \(5+20c\equiv0\)，故 \(c\equiv1\pmod5\)，即 \(a\equiv7\pmod{25}\)。再写 \(a=7+25d\)，模 \(125\) 得 \(50+350d\equiv0\)，即 \(2+4d\equiv0\pmod5\)，所以 \(d\equiv2\pmod5\)。因此所求值为 \(57\pmod{125}\)，确有 \(57^2+1=3250=26\cdot125\)。导数 \(2a\) 始终模 \(5\) 非零，单根提升定理保证存在唯一的 \(5\) 进根具有这个初始剩余类。
\end{solution}

\begin{optionalexercise}\label{b3:ex:23-series-root-factors}
在 \(\mathbb Q[\![t]\!]\) 中求 \(X^2-X-t\) 同余于 \(0\pmod t\) 的根到 \(t^4\) 为止，并写出完整多项式的两个首一因子。解释为何这些因子不都属于 \(\mathbb Q[t]_{(t)}[X]\)。
\end{optionalexercise}
\begin{solution}
设根为 \(a=c_1t+c_2t^2+c_3t^3+c_4t^4+\cdots\)。比较 \(a^2-a-t=0\) 的各阶系数，依次得到
\(c_1=-1\)、\(c_2=c_1^2=1\)、\(c_3=2c_1c_2=-2\)、\(c_4=2c_1c_3+c_2^2=5\)。另一个根为 \(1-a\)，所以准确的分解是 \((X-a)(X-(1-a))\)，其中 \(a\) 表示完整的 Hensel 根。若 \(a\in\mathbb Q(t)\)，则判别式 \(1+4t\) 为有理函数的平方，与它的简单零点矛盾。因此两个一次因子的系数不能都来自原局部环。
\end{solution}

\begin{optionalexercise}\label{b3:ex:23-repeated-root-henselian}
设 \(k\) 为域，\(R=k[\varepsilon]/(\varepsilon^2)\)。证明 \(R\) 为 Hensel 环，而 \(T^2\) 模极大理想的重根 \(0\) 有多个提升。由此说明 Hensel 性没有取消单根条件。
\end{optionalexercise}
\begin{solution}
\(R\) 为局部环，极大理想 \((\varepsilon)\) 的平方为零，所以其进拓扑的逆极限从第二阶起稳定，\(R\) 完备分离，因而 Hensel。对每个 \(c\in k\)，元素 \(c\varepsilon\) 都同余于零且平方为零，至少 \(0\) 与 \(\varepsilon\) 不同。导数 \(2T\) 在约化根处为零，故唯一单根提升定理不适用。
\end{solution}

\begin{optionalexercise}\label{b3:ex:23-cubic-hensel-factor}
在 \(\mathbb Z_7[T]\) 中，将 \(T^3-T+7\) 分解为一个约化为 \(T\) 的首一次因子和一个约化为 \(T^2-1\) 的首二次因子，并求前一因子的常数项模 \(49\)。
\end{optionalexercise}
\begin{solution}
\(T\) 与 \(T^2-1\) 模 \(7\) 互素，所以唯一存在这样的分解。设同余于零的根为 \(a\)，则
\[
T^3-T+7=(T-a)\bigl(T^2+aT+a^2-1\bigr),
\]
因为余项为 \(a^3-a+7=0\)。写 \(a=7c\)，模 \(49\) 的根方程为 \(-7c+7\equiv0\)，所以 \(a\equiv7\pmod{49}\)，首一次因子的常数项为 \(-7\pmod{49}\)。二次因子的两个剩余根还可分别提升，但本题的因式提升不需要先选择它们。
\end{solution}

\begin{optionalexercise}\label{b3:ex:23-henselization-universal}
仅用 Hensel 化的泛性质证明：若局部环 \(R\) 已经 Hensel，则 \(R\to R^{\mathrm h}\) 是同构；两个 Hensel 化之间的相容同构唯一。再解释为什么相同的有限阶商不能证明一个环已经完备。
\end{optionalexercise}
\begin{solution}
若 \(R\) Hensel，恒等映射给出唯一局部同态 \(v:R^{\mathrm h}\to R\)，使 \(v\iota=1_R\)。两个从 \(R^{\mathrm h}\) 到自身并延拓 \(\iota\) 的局部同态 \(1\) 与 \(\iota v\) 由唯一性相同，故互逆。对两个 Hensel 化，分别使用泛性质得到相容映射，再用同样的复合论证即得唯一同构。至于有限阶商，正文中的 \(\mathbb Q[t]_{(t)}\)、它的 Hensel 化及 \(\mathbb Q[\![t]\!]\) 有相同的各阶商，但前两者没有包含全部相容级数；完备性要求环到这些商的逆极限为同构，单独比较每个商没有证明满射性。
\end{solution}

\begin{optionalexercise}\label{b3:ex:23-node-lifting-obstruction}
设 \(k\) 为域，\(B=k[x,y]/(xy)\)。用一个平方零提升问题证明 \(k\to B\) 不形式光滑。
\end{optionalexercise}
\begin{solution}
取 \(C=k[\varepsilon]/(\varepsilon^3)\)、\(J=(\varepsilon^2)\)。定义 \(u:B\to C/J\) 将 \(x,y\) 都送到 \(\varepsilon\)，这是同态，因为乘积在 \(C/J\) 中为零。任意提升都必须把二者分别送到 \(\varepsilon+a\varepsilon^2\)、\(\varepsilon+b\varepsilon^2\)，而乘积恒为非零的 \(\varepsilon^2\)。关系不能成立，因此没有提升。失败发生在同时满足关系的要求，而不是变量像自身不能选取提升。
\end{solution}

\begin{optionalexercise}\label{b3:ex:23-free-differentials-not-smooth}
设 \(k\) 为特征二域，\(B=k[x]/(x^2)\)。证明 \(\Omega_{B/k}\) 为秩一自由模，但 \(B\) 不形式光滑。指出关系模分裂判据中的哪一步失败。
\end{optionalexercise}
\begin{solution}
微分关系为 \(\mathrm{d}(x^2)=2x\,\mathrm{d}x=0\)，所以 \(\Omega_{B/k}=B\,\mathrm{d}x\)。取 \(P=k[X]\)、\(I=(X^2)\)，则 \(I/I^2\) 为以 \(X^2\) 的像为基的秩一自由 \(B\) 模，但到 \(B\,\mathrm{d}X\) 的映射为零，没有左逆。也可取 \(C=k[T]/(T^4)\)、\(J=(T^2)\)，恒等映射 \(B\to C/J\) 的任何候选提升把 \(x\) 送到 \(T+h\)、\(h\in J\)，其平方为 \(T^2\ne0\)，故没有提升。这表明微分模投射只是形式光滑的一个后果，不能独自代替关系模的分裂。
\end{solution}

\begin{optionalexercise}\label{b3:ex:23-etale-stability}
证明有限展示平展映射在复合和任意基变换下保持平展。证明中应说明有限展示条件为何也得到保持。
\end{optionalexercise}
\begin{solution}
先使用平展等价于有限展示且形式平展。形式平展在复合与基变换下保持已在正文证明。基变换把有限个生成元与有限条关系的系数映到新底环，仍给出有限展示。对复合 \(A\to B\to C\)，取 \(B\) 的有限生成元和关系，再添上 \(C\) 在 \(B\) 上的有限生成元；后者有限条关系所含的 \(B\) 系数都能写成前一组生成元的多项式。把这些关系一并写在 \(A\) 上，得到有限展示。因此复合与基变换仍满足平展的等价条件。
\end{solution}

\begin{optionalexercise}\label{b3:ex:23-hyperbola-standard-model}
设 \(A\) 为交换含幺环，\(B=A[x,y]/(xy-1)\)。给出标准光滑模型，计算 \(\Omega_{B/A}\)，并判断当 \(A\ne0\) 时映射是否平展。
\end{optionalexercise}
\begin{solution}
环中 \(x,y\) 互逆，所以 Jacobian 行 \((y,x)\) 中的 \(x\) 已为单位。以变量 \(y\) 对应的列为子式，得到相对维数一的标准光滑模型。微分关系为 \(y\,\mathrm{d}x+x\,\mathrm{d}y=0\)，故 \(\mathrm{d}y=-x^{-2}\mathrm{d}x\)，而 \(\Omega_{B/A}=B\,\mathrm{d}x\)。若 \(A\ne0\)，则 \(B\cong A[x,x^{-1}]\ne0\)，这个微分模不为零，所以映射光滑且平坦，却不平展。
\end{solution}

\begin{optionalexercise}\label{b3:ex:23-gaussian-integers-etale}
对 \(\mathbb Z\to\mathbb Z[i]\) 计算相对微分模，证明该映射平坦，并准确给出其平展的点。
\end{optionalexercise}
\begin{solution}
\(\mathbb Z[i]\cong\mathbb Z[T]/(T^2+1)\)，所以作为整数模以 \(1,i\) 为基，因而平坦。相对微分模为
\(\mathbb Z[i]/(2i)\,\mathrm{d}i=\mathbb Z[i]/(2)\,\mathrm{d}i\)，因为 \(i\) 可逆。故在 \(D(2)\) 上导数可逆，映射平展；在包含 \(2\) 的素理想处该模非零，不平展。\(2=-i(1+i)^2\)，唯一这样的素理想是 \((1+i)\)。这显示有限自由不自动非分歧，整数 \(2\) 必须保留在判断中。
\end{solution}

\begin{optionalexercise}\label{b3:ex:23-closed-point-unramified}
设 \(A=k[t]\)、\(B=A/(t)=k\)。证明 \(\Omega_{B/A}=0\)，但 \(B\) 不是平坦 \(A\) 模，并直接给出形式光滑性失败的平方零问题。
\end{optionalexercise}
\begin{solution}
对商代数，所有元素都来自 \(A\)，所以任意 \(A\) 导子为零，\(\Omega_{B/A}=0\)。\(A\) 上的单射乘法 \(A\xrightarrow{t}A\) 与 \(B\) 张量后成为零映射 \(B\to B\)，不是单射，故不平坦。取 \(C=A/(t^2)\)、\(J=(t)C\)，则 \(C/J=B\)。若恒等映射 \(B\to C/J\) 有 \(A\) 代数提升，\(t\) 在 \(B\) 中为零会迫使其在 \(C\) 中的像为零，但它实际非零。这说明唯一性条件可以成立而存在性条件失败。
\end{solution}

\begin{optionalexercise}\label{b3:ex:23-product-diagonal}
令 \(C=A\times A\)，其中 \(A\to C\) 为对角同态。用 \(e_1=(1,0)\)、\(e_2=(0,1)\) 写出正文对角分裂中的元素 \(z\)，并显式分裂任意 \(C\) 模的映射 \(C\otimes_AM\to M\)。
\end{optionalexercise}
\begin{solution}
取 \(z=e_1\otimes e_1+e_2\otimes e_2\)。乘法把它送到 \(e_1+e_2=1\)，而对每个 \(c=(a,b)\)，有 \((c\otimes1)z=(1\otimes c)z\)。所以截面为
\(m\mapsto e_1\otimes e_1m+e_2\otimes e_2m\)。它为 \(C\) 线性且复合乘法为恒等。对应的分解是 \(M=e_1M\oplus e_2M\)；若 \(M\) 在 \(A\) 上平坦，两个直和项均平坦，这也直接证明 \(M\) 在 \(A\times A\) 上平坦。
\end{solution}

\begin{optionalexercise}\label{b3:ex:23-finite-flat-quotient}
设 \(C\) 为交换含幺环，\(J\subseteq C\) 为有限生成理想。证明 \(C/J\) 在 \(C\) 上平坦，当且仅当 \(J=(e)\) 对某个幂等元 \(e\in C\) 成立。解释相应闭子谱为何也为开子谱。
\end{optionalexercise}
\begin{solution}
若商平坦，用它张量 \(J\hookrightarrow C\) 得到单射 \(J/J^2\to C/J\)，但映射为零，故 \(J=J^2\)。正文的行列式技巧给出 \(J=(e)\)、\(e^2=e\)。反之，\(C=eC\oplus(1-e)C\)，且 \(C/J\cong(1-e)C\) 为自由模 \(C\) 的直和项，故投射并平坦。对每个素理想，\(e(1-e)=0\) 且 \(e+(1-e)=1\)，所以恰有一个幂等因子属于该素理想。因此 \(V(e)=D(1-e)\)，既闭又开。
\end{solution}

\begin{optionalexercise}\label{b3:ex:23-frobenius-fiber-length}
设 \(k\) 为特征 \(p>0\) 的代数闭域，\(k[s]\to k[t]\) 由 \(s\mapsto t^p\) 给出。计算点 \(s=a\) 上的纤维及其长度，说明幂零纤维为何并不与平坦性矛盾。
\end{optionalexercise}
\begin{solution}
唯一取 \(b\in k\) 使 \(b^p=a\)，则纤维为
\(k[t]/(t^p-a)=k[t]/((t-b)^p)\cong k[\varepsilon]/(\varepsilon^p)\)，长度为 \(p\)。原扩张以 \(1,t,\ldots,t^{p-1}\) 为自由基，秩正是 \(p\)，所以特殊纤维长度与自由秩一致，完全符合平坦性。它有幂零元说明非分歧性失败；平坦性本身不要求纤维约化。
\end{solution}

\begin{optionalexercise}\label{b3:ex:23-cusp-other-characteristics}
撤去尖点计算中的特征零假设，对任意域 \(k\) 计算
\(\Omega_{k[t]/k[t^2,t^3]}\)。分别给出特征二、特征三及其他特征的答案，并判断原点外是否仍为同构。
\end{optionalexercise}
\begin{solution}
传递正合列仍给出
\(\Omega_{B/A}=B/(2t,3t^2)\,\mathrm{d}t\)。特征二时，\(2=0\)、\(3=1\)，所以商为 \(B/(t^2)\,\mathrm{d}t\)；特征三时，\(2\) 可逆，商为 \(B/(t)\,\mathrm{d}t\)；其他特征下也因 \(2\) 可逆得到 \(B/(t)\,\mathrm{d}t\)。在特征二中原点处的微分模长度为二，其余情形长度为一，均不为零。原点外的公式 \(t=t^3/t^2\) 不涉及特征，故同构结论始终成立。
\end{solution}

\begin{optionalexercise}\label{b3:ex:23-cusp-module-presentation}
设 \(A=k[t^2,t^3]\)、\(B=k[t]\)，记 \(x=t^2,y=t^3\)。证明作为 \(A\) 模，\(B\) 有展示
\[
A^2\xrightarrow{\begin{psmallmatrix}-y&-x^2\\x&y\end{psmallmatrix}}
A^2\xrightarrow{(1\quad t)}B\longrightarrow0.
\]
利用这个展示重算尖点纤维的维数。
\end{optionalexercise}
\begin{solution}
后一箭头满射，因为 \(B=A+At\)。两列确实给出关系：\(-y+xt=0\)、\(-x^2+yt=0\)。反之，若 \(a+bt=0\)、\(a,b\in A\)，则 \(bt\in A\)。\(A\) 的多项式恰好没有一次项，所以 \(b\) 的常数项必须为零，即 \(b\in(x,y)\)。写 \(b=xc+yd\)，便有 \(a=-yc-x^2d\)，因此 \((a,b)\) 是所列两列的线性组合，证明正合。模去 \((x,y)\) 后矩阵为零，张量积的右正合性给出纤维为 \(k^2\)，与 \(k[t]/(t^2)\) 的基 \(1,t\) 一致。
\end{solution}

\begin{optionalexercise}\label{b3:ex:23-inseparable-closed-point}
设 \(k=\mathbb F_p(s)\)，\(B=k[T]\)，\(\mathfrak q=(T^p-s)\)。证明 \(k\to B\) 光滑、\(B_{\mathfrak q}\) 正则，但剩余域 \(B/\mathfrak q\) 不可分于 \(k\)。解释这为何不与正文比较定理冲突。
\end{optionalexercise}
\begin{solution}
多项式环有任意变量像的提升，所以 \(k\to B\) 光滑。\(T^p-s\) 不可约，\(B_{\mathfrak q}\) 是主理想整环在非零素理想处的局部化，为 DVR，故正则。剩余域 \(L=k[u]\)、\(u^p=s\)，其极小多项式导数为零，故扩张纯不可分且次数为 \(p\)。正文没有断言光滑的每个点都具有可分剩余域；本例的整个多项式代数仍有一个可自由提升的变量，而商掉 \(T^p-s\) 所得的 \(k\to L\) 已经是另一个映射，后者不光滑。
\end{solution}

\begin{optionalexercise}\label{b3:ex:23-square-root-cover}
设 \(A=\mathbb C[s,s^{-1}]\)，\(B=A[T]/(T^2-s)\)。证明 \(\operatorname{Spec}B\to\operatorname{Spec}A\) 是平展覆盖，但在任何非空 Zariski 开集上都没有与底环相容的截面。说明“平展局部存在平方根”在这个例子中的含义。
\end{optionalexercise}
\begin{solution}
\(B\) 有自由基 \(1,T\)，且 \(T\) 为单位，所以导数 \(2T\) 可逆，映射平展。自由秩二的模忠实平坦，因而素谱映射满射，构成平展覆盖。若在某个非空开集上有截面，限制到一般点便给出 \(\mathbb C(s)\) 中的元素 \(r\)，满足 \(r^2=s\)。但平方的零点阶数为偶数，而 \(s\) 在零点有一阶零点，矛盾。所谓平展局部存在，是说先经过这个平展覆盖，即把底环换成 \(B\)，平方根 \(T\) 就已经在新环中；并不是说原来的某个非空 Zariski 开集已经含有这样的有理函数。
\end{solution}

% !TeX root = ../../Book3.tex
% 纲要标题：## 第24章*　概形与射影方法初步
% 纲要位置：计划纲要.md，第 2734 行。

\OptionalChapter{概形与射影方法初步}{b3:ch:24}
\BookChapterNumber{b3:ch:24}

\begin{chapterabstract}
本章把环的局部化组织成结构层，建立仿射概形与拟凝聚层的基本构造及模层对应。分次环给出 Proj，射影直线的仿射图用于计算扭曲层和上同调；纤维积描述局部对象的交叠，平展性、分离性、有限型与固有性则描述态射的进一步性质。
\end{chapterabstract}

\begin{learninggoals}
\begin{itemize}
\item 从分式黏合证明结构层的基本开截面与茎公式。
\item 在环同态、仿射概形态射、模与拟凝聚层之间准确转换。
\item 用开集及过渡同构构造概形和局部自由层，区分纤维与茎。
\item 从齐次环计算 Proj 的仿射图及扭曲层，完成射影直线的低次上同调计算。
\item 从仿射局部模型判断平展性，用对角判断分离性，检验有限型与普遍闭性，并核对 Čech 比较所需的交集条件。
\end{itemize}
\end{learninggoals}

把域 \(k\) 上的两份仿射直线各自的非零部分按相同坐标对应起来，两个原点仍各自保留。每一份直线上的多项式运算都清楚，交叠处也允许 \(1/x\)；难处在于说明这些局部函数怎样彼此限制、怎样黏合，以及所得空间是否还来自一个单一的环。为回答这些问题，需要在每个开集上记录可合法使用的函数，并规定它们在重叠处如何比较。空间的结构因而还须包含这些函数环及限制映射。

\ProseParagraphBreak

层把每个开集上的函数与开集之间的限制一并记录，仿射概形则从环的局部化建立这样的结构。模也能以相同方式给出随开集变化的截面。分次环则把齐次方程变成射影空间上的局部方程。射影直线的两张仿射图让我们可以具体比较不同坐标中的截面；它们的相容条件决定扭曲层的全局截面，交叠上无法消去的差则给出一次上同调。前述双原点例子还提醒我们，黏合成功与分离性是不同的问题。

\ProseParagraphBreak

本章主要需要\chapref{b3:ch:01}的素谱、\chapref{b3:ch:02}的模局部化、\chapref{b3:ch:11}的分次理论，以及\cref{b3:thm:finite-presentation-flat-projective}。层和概形在此从定义开始。\secref{b3:sec:2401}的结构层、\secref{b3:sec:2402}的仿射模层对应与黏合，以及\secref{b3:sec:2403}的 Proj 仿射图和射影直线扭曲层，组成相对独立的基础部分；局部自由而无非零全局截面的例子也在其中。\secref{b3:sec:2404}用纤维积与基变换讨论平展性、分离性和固有性，平展的局部判据来自\secref{b3:sec:2303}。\secref{b3:sec:2405}的层上同调使用复形与内射消解，并回用\cref{b3:lem:injective-cech-acyclic}中内射模的局部化与支集挠子性质。

\ProseParagraphBreak
可配合 Hartshorne 的《Algebraic Geometry》\cite[Chapter II, \S\S 1--2 and \S 5]{Hartshorne1977}学习层、仿射概形与模层，Proj 的仿射图见同书\cite[Chapter II, Proposition 2.5, pp. 76--77]{Hartshorne1977}。分式黏合已在\cref{b3:lem:localization-sheaf-gluing}证明；本章对仿射对应与射影空间闭子概形的固有性给出直接证明，最后一节写出 Čech 比较的双复形证明，仿射上同调消失则给出概要，并注明所采用的层论结果与省略的技术步骤。

% !TeX root = ../../Book3.tex
% 纲要标题：### 24.1 结构层
% 纲要位置：计划纲要.md，第 2736 行。

\section{结构层}
\label{b3:sec:2401}


% 纲要内容（仅作写作计划，尚非教材正文）：
%
% * 基本开限制的方向及非单射情形；环值、集合值和群值层采用同一黏合公理
% * 空覆盖、芽与剩余域值的区别；由局部分式族建立结构层，再证明基本开截面与茎公式
% * 局部同态保持指定点的消失条件；仿射态射对应、Noether 概形与拟紧性的条件
%
% ---
%

前面已经把环的素理想组织成拓扑空间，但拓扑本身无法区分约化点与加厚点。本节把各个基本开集上的局部化同时保留下来，并说明这些环如何限制、黏合。这种记录邻域中代数信息的方法，就是结构层。

\subsection{基本开集上的局部化}
\label{b3:subsec:2401:01}

在交换环 \(A\) 的素谱上，\cref{b3:thm:spec-localization}把 \(D(f)\) 识别为 \(\operatorname{Spec}A_f\)。把开集缩小以后，原来尚不可逆的元素可能变得可逆，所以截面应从较大开集限制到较小开集。具体地，若 \(D(g)\subseteq D(f)\)，则 \(f\) 在 \(A_g\) 中可逆：包含关系给 \(g^N=af\)，从而 \(a/g^N\) 是其逆。因此有唯一限制同态 \(A_f\rightarrow A_g\)，且这些同态与复合相容。

\ProseParagraphBreak
这种限制同态不一定单射，因为缩小开集也可能使原来的元素消失。例如，对域 \(k\)，令 \(A=k\times k\)、\(e=(1,0)\)，则 \(A\to A_e\cong k\) 是第一分量投影，它把非零的 \((0,1)\) 送到零。这里记录的是限制同态，不能把不同开集上的环一律放成包含关系。从相容的局部分式恢复整体元素所需的代数步骤，已由\cref{b3:lem:localization-sheaf-gluing}证明。

\ProseParagraphBreak
对 \(A_f\) 使用该引理，可以处理 \(D(f)\) 的基本开覆盖；任意这样的覆盖由\cref{b3:prop:spec-quasicompact}取有限子覆盖，未选中的开集再由唯一性核对。因此局部化既能限制，也能在交集相容时黏合。这正是下一小节层公理的代数来源。

\subsection{仿射概形的结构层}
\label{b3:subsec:2401:02}

先写本章所需的环值情形。随后将值换成集合或交换群时，限制与唯一黏合仍采用同一套公理，只是保留的代数运算不同。

\begin{definition}\label{b3:def:sheaf-rings}
设 \(X\) 为拓扑空间。对每个开集 \(U\subseteq X\) 给定交换环 \(\mathcal F(U)\)，对 \(V\subseteq U\) 给定保持单位的限制同态 \(\rho_{UV}:\mathcal F(U)\rightarrow\mathcal F(V)\)。若 \(\rho_{UU}=\operatorname{id}\)，且对 \(W\subseteq V\subseteq U\) 有 \(\rho_{VW}\rho_{UV}=\rho_{UW}\)，则称这些数据为环的\term[yuceng]{预层}[presheaf]。若再对任意开覆盖 \(U=\bigcup_iU_i\)，在两两交集上限制相同的 \(s_i\in\mathcal F(U_i)\) 都来自唯一的 \(s\in\mathcal F(U)\)，则称其为环的\term[ceng]{层}[sheaf]。空覆盖约定使 \(\mathcal F(\varnothing)\) 为零环。元素 \(s\in\mathcal F(U)\) 称为 \(U\) 上的截面，也记 \(\Gamma(U,\mathcal F)\)。
\end{definition}
\symidx{\(\Gamma(U,\mathcal F)\)：层 \(\mathcal F\) 在开集 \(U\) 上的截面}

\begin{definition}\label{b3:def:sheaf-sets-groups}
设 \(X\) 为拓扑空间。对每个开集 \(U\subseteq X\) 给定集合 \(\mathcal F(U)\)，对每个包含关系 \(V\subseteq U\) 给定限制映射 \(\rho_{UV}:\mathcal F(U)\rightarrow\mathcal F(V)\)，并记 \(s|_V=\rho_{UV}(s)\)。若 \(\rho_{UU}=\operatorname{id}\)，且对 \(W\subseteq V\subseteq U\) 有 \(\rho_{VW}\rho_{UV}=\rho_{UW}\)，则称这些数据为集合值预层。若再对任意开覆盖 \(U=\bigcup_iU_i\)，满足
\[
s_i|_{U_i\cap U_j}=s_j|_{U_i\cap U_j}
\quad\text{对所有}\;i,j\;\text{成立}
\]
的截面族 \(s_i\in\mathcal F(U_i)\) 都来自唯一的 \(s\in\mathcal F(U)\)，则称其为集合值层。若各 \(\mathcal F(U)\) 为交换群，且限制映射为群同态，则在相同公理下分别称为交换群值预层、交换群值层。对层，空集的空覆盖要求 \(\mathcal F(\varnothing)\) 恰有一个元素；在交换群值情形，这就是零群。
\end{definition}

环层忘掉乘法以后是交换群值层，再忘掉加法以后是集合值层，因为限制与唯一黏合的条件都保留下来。上一小节的 \(A_f\) 因而同时给出三种取值的局部模型；对 \(A\) 模 \(M\)，\cref{b3:lem:localization-sheaf-gluing}中的 \(M_f\) 也给出基本开集上的交换群值黏合模型。空基本开集 \(D(0)\) 上分别取零环、零群或单点集，与空覆盖的要求一致。

\ProseParagraphBreak

集合值层的态射是一组与限制相容的集合映射；对交换群值层、环层，则分别要求群同态、保持单位的环同态。这三种情形都可以谈截面、芽和茎：若两个截面分别定义在含 \(x\) 的开邻域上，并在某个更小的含 \(x\) 邻域上限制相等，就定义相同的芽。所有这样的芽组成茎 \(\mathcal F_x=\varinjlim_{x\in U}\mathcal F(U)\)。在交换群值或环值情形，先把代表截面限制到共同邻域，再作运算；限制保留运算，保证所得结果与代表截面的选择无关，因而得到茎上的交换群或环结构。芽只保留任意充分小邻域内的信息，所以不是把截面简单地“代入点”所得的数。
\symidx{\(\mathcal F_x\)：层 \(\mathcal F\) 在点 \(x\) 处的茎}

\ProseParagraphBreak
一般开集未必形如一个 \(D(f)\)，因而没有预先给定的单一局部化环。结构层先要求每一点附近都能用同一个分式表示，再把这些局部表示黏合到任意开集上；下面的基本开截面定理将说明，这个构造在 \(D(f)\) 上恰好返回熟悉的 \(A_f\)。

\begin{definition}\label{b3:def:affine-structure-sheaf}
设 \(A\) 为交换环，\(X=\operatorname{Spec}A\)。对开集 \(U\subseteq X\)，令 \(\mathcal O_X(U)\) 由所有族 \((s_{\mathfrak p})_{\mathfrak p\in U}\)、\(s_{\mathfrak p}\in A_{\mathfrak p}\) 组成，并要求每个点都有开邻域 \(V\subseteq U\) 及 \(a,b\in A\)，使所有 \(\mathfrak q\in V\) 都满足 \(b\notin\mathfrak q\) 及 \(s_{\mathfrak q}=a/b\)。配以逐点运算和通常限制，所得层称为素谱的\term[jiegou-ceng]{结构层}[structure sheaf]。
\end{definition}
\symidx{\(\mathcal O_X\)：仿射概形的结构层}

限制保持局部表示；相容的族可按点拼合，拼合后仍在原邻域上使用原分式，因而确实满足层公理。这里的值在局部环 \(A_{\mathfrak p}\) 中，而非仅在剩余域 \(\kappa(\mathfrak p)\) 中，这使结构层仍能记录幂零元。

\begin{theorem}\label{b3:thm:affine-sheaf-sections}
设 \(A\) 为交换环，\(X=\operatorname{Spec}A\)，\(f\in A\)，\(\mathfrak p\in X\)。自然映射给出
\[
\Gamma(D(f),\mathcal O_X)\cong A_f,\qquad
\mathcal O_{X,\mathfrak p}\cong A_{\mathfrak p}.
\]
特别地，\(\Gamma(X,\mathcal O_X)\cong A\)；这些同构与限制和局部化相容。
\end{theorem}
\begin{proof}
先识别茎与局部环，再证明基本开截面映射单射，最后用有限子覆盖上的分式黏合证明满射。三个步骤分别使用局部化的相等判据、局部零检测与清分母黏合。

先看茎。分式 \(a/b\)、\(b\notin\mathfrak p\)，在 \(D(b)\) 上给出截面，所以每个局部环元素都来自芽。若两个分式在 \(A_{\mathfrak p}\) 中相等，存在 \(c\notin\mathfrak p\) 清除它们差的分子，故它们在 \(D(b_1b_2c)\) 上相等。这证明芽与局部环一一对应。

自然映射 \(A_f\rightarrow\Gamma(D(f),\mathcal O_X)\) 是单射。事实上，若一个分式在每个 \(A_{\mathfrak p}\)、\(\mathfrak p\in D(f)\)，中为零，按分式零判据可用一组更小的基本开集使它为零；取有限子覆盖后，由\cref{b3:lem:localization-sheaf-gluing}的单射部分，在 \(A_f\) 中也为零。

给定 \(D(f)\) 上截面，将其局部表示邻域缩为基本开集 \(D(g_i)\)。若局部表达的分母为 \(b_i\)，则 \(D(g_i)\subseteq D(b_i)\) 保证 \(b_i\) 在 \(A_{g_i}\) 可逆，故该局部截面确实来自 \(A_{g_i}\)。有限子覆盖上的这些元素在交集一致；刚证的单射性把截面相等变成局部化中的相等。对 \(A_f\) 应用黏合引理，便得到唯一的 \(A_f\) 元素。所有映射都由同一分式给出，故自动与限制相容。
\end{proof}

例如，对域 \(k\)，在 \(\operatorname{Spec}k[x]\) 的原点 \((x)\) 处，分式 \(1/(1-x)\) 定义在含原点的 \(D(1-x)\) 上，给出茎 \(k[x]_{(x)}\) 中的一个芽。它不是全局多项式；但模去局部环的极大理想后，它的值为 \(1\in k\)。常数 \(1\) 与这个分式具有相同的剩余域值，却给出不同的芽，差为非零的 \(x/(1-x)\)。

\ProseParagraphBreak
\(\operatorname{Spec}k\) 与 \(\operatorname{Spec}k[\varepsilon]/(\varepsilon^2)\) 都只有一个点，结构层的全局截面却分别为 \(k\) 与 \(k[\varepsilon]/(\varepsilon^2)\)。因此相同的拓扑空间可以承载不同的代数结构。结构层的基本性质可对照 Hartshorne 的《Algebraic Geometry》\cite[Chapter II, Proposition 2.2, pp. 71--72]{Hartshorne1977}。

\subsection{局部环空间与态射}
\label{b3:subsec:2401:03}

\begin{definition}\label{b3:def:locally-ringed-space}
设 \(X\) 为拓扑空间，\(\mathcal O_X\) 为其上的交换环层。若每个茎 \(\mathcal O_{X,x}\) 都是局部环，则称 \((X,\mathcal O_X)\) 为\term[jubu-huan-kongjian]{局部环空间}[locally ringed space]。从 \((X,\mathcal O_X)\) 到 \((Y,\mathcal O_Y)\) 的态射由连续映射 \(f:X\rightarrow Y\) 及与限制相容的同态
\[
f^\#_V:\mathcal O_Y(V)\longrightarrow\mathcal O_X(f^{-1}V)
\]
组成，且要求每个诱导茎同态 \(f_x^\#:\mathcal O_{Y,f(x)}\rightarrow\mathcal O_{X,x}\) 为局部同态。若两局部环的极大理想分别记为 \(\mathfrak m_{f(x)}\)、\(\mathfrak m_x\)，这要求
\[
(f_x^\#)^{-1}(\mathfrak m_x)=\mathfrak m_{f(x)}.
\]
\end{definition}

一个芽在目标点的剩余域中为零，当且仅当其拉回在源点的剩余域中为零；在目标点可逆的芽，拉回后仍可逆。因此局部同态不仅给出邻域函数的拉回，也使这一拉回尊重所指定的点。

\begin{definition}\label{b3:def:scheme}
设 \((X,\mathcal O_X)\) 为局部环空间。若存在交换环 \(A\) 及局部环空间同构
\[
(X,\mathcal O_X)\cong(\operatorname{Spec}A,\mathcal O_{\operatorname{Spec}A}),
\]
则称 \(X\) 为\term[fangshe-gaixing]{仿射概形}[affine scheme]。若 \(X\) 的每点都有开邻域 \(U\)，使 \((U,\mathcal O_X|_U)\) 为仿射概形，则称 \(X\) 为\term[gaixing]{概形}[scheme]。概形态射就是局部环空间态射。
\end{definition}

在这个定义中，\(\mathcal O_X|_U\) 表示只在 \(U\) 的开集上取截面。由\cref{b3:thm:affine-sheaf-sections}，仿射概形的基本开集仍为仿射概形；一般开子概形则不必仿射。

\begin{definition}\label{b3:def:noetherian-scheme}
设 \(X\) 为概形。若 \(X\) 有由 Noether 环的素谱组成的仿射开覆盖，则称 \(X\) 为局部 Noether 概形；若这样的覆盖还可以取为有限覆盖，则称 \(X\) 为 Noether 概形。后一个条件等价于局部 Noether 且底空间拟紧。
\end{definition}

若 \(X\) 局部 Noether 且拟紧，从定义中的 Noether 仿射开覆盖抽取有限子覆盖，即得到 Noether 概形。反之，每个仿射谱拟紧，有限个这样的开集覆盖 \(X\) 时，任意开覆盖都可在每个仿射块上取有限子覆盖，再合成 \(X\) 的有限子覆盖，所以底空间拟紧。

\begin{theorem}\label{b3:thm:affine-scheme-hom}
设 \(A,B\) 为交换环。取全局截面给出自然双射
\[
\operatorname{Hom}_{\mathrm{Sch}}(\operatorname{Spec}B,\operatorname{Spec}A)
\cong\operatorname{Hom}_{\mathrm{Ring}}(A,B),
\]
并把态射复合对应于环同态的反向复合。
\end{theorem}
\begin{proof}
给定 \(\varphi:A\rightarrow B\)，点映射为 \(\mathfrak q\mapsto\varphi^{-1}(\mathfrak q)\)，连续性已由\cref{b3:prop:spectrum-functor}证明。分式映射 \(a/s\mapsto\varphi(a)/\varphi(s)\) 在相应基本开集上相容，所以给出层映射；茎映射的极大理想原像显然为 \(\varphi^{-1}(\mathfrak q)A_{\varphi^{-1}(\mathfrak q)}\)，故是局部同态。

反之，从概形态射 \((f,f^\#)\) 取全局截面得 \(\varphi:A\rightarrow B\)。对于 \(\mathfrak q\in\operatorname{Spec}B\)，一个 \(a\in A\) 属于 \(f(\mathfrak q)\)，等价于其芽在目标局部环 \(A_{f(\mathfrak q)}\) 的极大理想中。局部同态的极大理想原像条件使这又等价于其像在源局部环 \(B_{\mathfrak q}\) 的极大理想中，即 \(\varphi(a)\in\mathfrak q\)。于是 \(f(\mathfrak q)=\varphi^{-1}(\mathfrak q)\)。在茎上，分母一旦可逆，映射只能是上述分式映射；层态射由所有芽决定，故也唯一。这证明两种构造互逆，复合的断言由拉回次序得到。
\end{proof}

这项对应说明，仿射概形的全局截面环既能恢复整个仿射概形，又能确定它到另一个仿射概形的态射。结构层保留了各局部化，但在仿射情形，这些局部数据仍由同一个全局环控制；一般概形还要保留仿射图之间的黏合数据。

\ProseParagraphBreak
局部性是实质条件。若 \(R\) 为 DVR、\(K\) 为其分式域，把 \(\operatorname{Spec}K\) 的点送到 \(\operatorname{Spec}R\) 的闭点，并在环层上用 \(R\hookrightarrow K\)，虽可得到环空间态射，却不是概形态射。取统一化元 \(\pi\)：它在 \(R\) 的闭点处属于极大理想，拉回到 \(K\) 却成为单位，因而没有保留指定点处的消失条件。等价地，\(K\) 的零极大理想收缩为 \((0)\)，不是 \(R\) 的极大理想。真正由环包含 \(R\hookrightarrow K\) 诱导的概形态射应把唯一点送到一般点 \((0)\)。定理及此条件的辨析见\cite[Chapter II, Proposition 2.3 and Example 2.3.2, pp. 73--74]{Hartshorne1977}。

% !TeX root = ../../Book3.tex
% 纲要标题：### 24.2 模与层
% 纲要位置：计划纲要.md，第 2742 行。

\section{模与层}
\label{b3:sec:2402}

把模的元素局部化到每个素点后，还须要求这些局部元素由分式在邻域内相容地表示。由此得到的伴随层保留张量积与正合列，并把模的局部性质转为层的局部性质。\cref{b3:thm:affine-quasicoherent-equivalence}从全局模恢复仿射空间上的层，本节末的黏合命题则从相容的局部层构造全局对象。层化、逆像、拉回和直像在定理前给出定义；它们在仿射图上的表达也用于下一节的 Proj 构造。


% 纲要内容（仅作写作计划，尚非教材正文）：
% * 模对应的拟凝聚层；仿射模层等价及取全局截面的正合性
% * 局部有限秩自由层、线性过渡矩阵与代数向量丛；截面层及线性态射对应
% * 局部化相容性和黏合

\subsection{模对应的拟凝聚层}
\label{b3:subsec:2402:01}

环 \(A\) 在谱上形成结构层，\(A\) 模 \(M\) 也应在每个点保留 \(M_{\mathfrak p}\)。这里的局部元素不是标量函数，但同样可以在分母不消失的区域内表示为分式。

\begin{definition}\label{b3:def:associated-module-sheaf}
设 \(A\) 为交换环，\(M\) 为 \(A\) 模，\(X=\operatorname{Spec}A\)。对开集 \(U\subseteq X\)，令 \(\widetilde M(U)\) 为所有族 \((s_{\mathfrak p})_{\mathfrak p\in U}\)、\(s_{\mathfrak p}\in M_{\mathfrak p}\) 的集合，要求每个 \(\mathfrak p\in U\) 都有开邻域 \(V\subseteq U\) 以及 \(m\in M,f\in A\)，使所有 \(\mathfrak q\in V\) 都满足 \(f\notin\mathfrak q\) 及 \(s_{\mathfrak q}=m/f\)。配以限制与逐点模运算，所得 \(\mathcal O_X\) 模层称为 \(M\) 的伴随层。这里一个 \(\mathcal O_X\) 模层 \(\mathcal F\) 要求每个 \(\mathcal F(U)\) 为 \(\mathcal O_X(U)\) 模，且限制与标量限制相容。
\end{definition}

\cref{b3:lem:localization-sheaf-gluing}本来就是对模证明的。用同样的局部族构造和分母核验，得到
\begin{equation}\label{b3:eq:module-sheaf-localization}
\Gamma(D(f),\widetilde M)=M_f,\qquad
(\widetilde M)_{\mathfrak p}=M_{\mathfrak p},\qquad
\Gamma(X,\widetilde M)=M.
\end{equation}
模同态在各局部化上给出相容同态，因而诱导层态射。层的核可逐开集取核；余核则一般须把局部余类重新黏合，不能直接把每个开集上的余核当作层。等价地，层序列的正合性在每个茎上检验：一个芽属于像，表示在某个足够小的邻域中有原像，不要求全局已经选好原像。

\begin{definition}\label{b3:def:quasicoherent-sheaf}
设 \(X\) 为概形，\(\mathcal F\) 为 \(\mathcal O_X\) 模层。若存在仿射开覆盖 \(X=\bigcup_iU_i\)、\(U_i=\operatorname{Spec}A_i\)，以及 \(A_i\) 模 \(M_i\)，使 \(\mathcal F|_{U_i}\cong\widetilde{M_i}\)，则称 \(\mathcal F\) 为\term[niningju-ceng]{拟凝聚层}[quasi-coherent sheaf]。在 Noether 概形上，若还能选各 \(M_i\) 为有限生成模，则称 \(\mathcal F\) 为\term[ningju-ceng]{凝聚层}[coherent sheaf]。
\end{definition}

后一用语在本章只用于 Noether 概形，不把“有限生成”作为任意底空间上凝聚性的完整定义。拟凝聚性首先按某个仿射覆盖定义；下一定理说明在仿射空间上，它仍由一个全局模决定。

\ProseParagraphBreak
先说明定理中所用的层运算。对拓扑空间上的预层 \(\mathcal P\)，取所有局部由 \(\mathcal P\) 的同一截面表示的芽族，得到层 \(\mathcal P^+\)，称为其层化。限制和黏合逐芽进行，正如结构层的构造；\(\mathcal P\rightarrow\mathcal P^+\) 在茎上为同构。若 \(\mathcal P\rightarrow\mathcal G\) 的目标已经是层，各局部代表的像便能唯一黏合，故该映射唯一经过 \(\mathcal P^+\)。这样，模层 \(\mathcal F,\mathcal G\) 的张量积就是预层 \(U\mapsto\mathcal F(U)\otimes_{\mathcal O_X(U)}\mathcal G(U)\) 的层化，其茎为 \(\mathcal F_x\otimes_{\mathcal O_{X,x}}\mathcal G_x\)。

\ProseParagraphBreak
对连续映射 \(f:Y\rightarrow X\)，逆像层 \(f^{-1}\mathcal F\) 是预层
\[
U\longmapsto\varinjlim_{f(U)\subseteq V}\mathcal F(V)
\]
的层化，其中 \(U\subseteq Y\)、\(V\subseteq X\) 开，转移映射为缩小 \(V\) 的限制。它的茎为 \((f^{-1}\mathcal F)_y=\mathcal F_{f(y)}\)：两边的代表都是 \(f(y)\) 邻域上的截面，而连续性允许把该邻域拉回到 \(y\) 的邻域。若 \(f\) 是概形态射，就可进一步扩张标量，定义
\(f^*\mathcal F=\mathcal O_Y\otimes_{f^{-1}\mathcal O_X}f^{-1}\mathcal F\)，称为模层的拉回。直像层则直接定义为 \((f_*\mathcal G)(V)=\mathcal G(f^{-1}V)\)；开覆盖的原像仍为开覆盖，所以它已经满足层公理。

\ProseParagraphBreak
一个闭点可以直接区分逆像与拉回。设 \(\mathfrak m\) 为 \(A\) 的极大理想，\(i:\operatorname{Spec}\kappa(\mathfrak m)\to\operatorname{Spec}A\) 为相应态射。单点空间上的层由唯一的茎确定，故
\[
i^{-1}\widetilde M=M_{\mathfrak m},\qquad
i^*\widetilde M
=\kappa(\mathfrak m)\otimes_{A_{\mathfrak m}}M_{\mathfrak m}
\cong M_{\mathfrak m}/\mathfrak mM_{\mathfrak m}.
\]
第一式仍以 \(A_{\mathfrak m}\) 为标量环，保留点附近的芽；第二式通过剩余域扩张标量，只保留该点的纤维值。例如 \(A=k[x]\)、\(\mathfrak m=(x)\)、\(M=A\) 时，前者为 \(k[x]_{(x)}\)，后者为 \(k\)：逆像并没有自动把在点处消失的函数模去。

\begin{theorem}\label{b3:thm:affine-quasicoherent-equivalence}
设 \(A\) 为交换环，\(X=\operatorname{Spec}A\)。函子 \(M\mapsto\widetilde M\) 与 \(\mathcal F\mapsto\Gamma(X,\mathcal F)\) 给出 \(A\) 模和 \(X\) 上拟凝聚层的互逆等价。它们保持正合列；对任意 \(A\) 模 \(M,N\)，有
\[
\widetilde{M\otimes_A N}\cong
\widetilde M\otimes_{\mathcal O_X}\widetilde N.
\]
对任意交换 \(A\) 代数 \(B\)，由 \(A\rightarrow B\) 诱导的 \(f:\operatorname{Spec}B\rightarrow X\) 满足 \(f^*\widetilde M\cong\widetilde{B\otimes_AM}\)。
\end{theorem}
\begin{proof}
先证明拟凝聚层的截面可以清分母。将其定义中的覆盖细化为有限个环境谱 \(X\) 的基本开集 \(D(g_i)\)，使 \(\mathcal F|_{D(g_i)}=\widetilde{N_i}\)。具体地，若 \(D_X(g)\subseteq U=\operatorname{Spec}B\) 且 \(\mathcal F|_U=\widetilde N\)，令 \(b\in B\) 是全局函数 \(g\) 在 \(U\) 上的限制。一个点处的 \(g\) 可逆与其限制 \(b\) 可逆等价，所以 \(D_X(g)=D_U(b)\)，而限制层来自 \(N_b\)。这证明所需细化保持模的描述；有限性来自 \(X\) 的拟紧性。

若 \(s\in\Gamma(X,\mathcal F)\) 限制到 \(D(f)\) 后为零，则它在各 \(N_i\) 中的像被 \(f\) 的某个幂零化。取统一的幂后，由层的唯一性有 \(f^Ns=0\)。若 \(t\in\Gamma(D(f),\mathcal F)\)，则在每个 \(D(g_i)\cap D(f)\) 上，它是 \(N_i\) 的分式。乘统一的 \(f^N\) 可延拓为 \(D(g_i)\) 上截面 \(t_i\)。在每个交集 \(D(g_ig_j)\) 上，差 \(t_i-t_j\) 限制到 \(D(fg_ig_j)\) 为零。对该仿射交集使用刚证的零化结论，再乘统一的 \(f^L\)，各截面即相容，能够黏合为全局截面。

因此，令 \(M=\Gamma(X,\mathcal F)\)，自然映射
\(M_f\rightarrow\Gamma(D(f),\mathcal F)\) 既单又满。它们在基本开集上相容，遂给出 \(\widetilde M\cong\mathcal F\)。式\eqref{b3:eq:module-sheaf-localization}给出另一方向的复合为恒等。层态射取全局截面后，其在基本开集上的作用由分式唯一确定，所以态射也得到互逆对应。

局部化正合，故模的正合列在取伴随层后仍正合。为证明逆函子的正合性，给定拟凝聚层的短正合列
\[
0\longrightarrow\mathcal F'\longrightarrow\mathcal F
\longrightarrow\mathcal F''\longrightarrow0,
\]
令 \(M'=\Gamma(X,\mathcal F')\)、\(M=\Gamma(X,\mathcal F)\)、\(M''=\Gamma(X,\mathcal F'')\)。上面得到的自然同构把它们在每个基本开集 \(D(f)\) 上的截面映射识别为
\(M'_f\rightarrow M_f\rightarrow M''_f\)。再取茎，层列的正合性给出
\[
0\longrightarrow M'_{\mathfrak p}\longrightarrow M_{\mathfrak p}
\longrightarrow M''_{\mathfrak p}\longrightarrow0
\quad\text{对每个 }\mathfrak p\in\operatorname{Spec}A\text{ 成立}.
\]
因此全局模列 \(0\rightarrow M'\rightarrow M\rightarrow M''\rightarrow0\) 在各素点局部化后正合。局部化保持核、像及余核，故这条模列的各个同调模在全部素点为零，因而本身为零：若某个模中有非零元素，取包含其零化子的极大理想，就能在该处局部化后仍检测到它。这证明全局模列正合，也就得到逆函子的正合性。张量积公式在每个基本开集上就是局部化与张量积的相容式。

最后，拉回模层按“先取局部截面的逆像，再把标量由 \(f^{-1}\mathcal O_X\) 扩为 \(\mathcal O_{\operatorname{Spec}B}\)”定义，即
\(f^*\mathcal F=\mathcal O_{\operatorname{Spec}B}\otimes_{f^{-1}\mathcal O_X}f^{-1}\mathcal F\)；在仿射基本开集上代入分式，恰得到 \(B\otimes_AM\) 及其局部化。这证明最后一式。
\end{proof}

例如 \(\widetilde{A/I}\) 在点 \(\mathfrak p\) 的茎为 \(A_{\mathfrak p}/IA_{\mathfrak p}\)，故它的支集是 \(V(I)\)，且仍保留商环中可能存在的幂零元。上述等价的标准叙述及清分母引理见\cite[Chapter II, Lemma 5.3, Proposition 5.4 and Corollary 5.5, pp. 112--113]{Hartshorne1977}。

\subsection{局部自由层与向量丛}
\label{b3:subsec:2402:02}

\begin{definition}\label{b3:def:locally-free-sheaf}
设 \(X\) 为概形，\(\mathcal E\) 为 \(\mathcal O_X\) 模层。若每点都有开邻域 \(U\)，使 \(\mathcal E|_U\cong\mathcal O_U^r\)、\(r\) 为有限非负整数，则称 \(\mathcal E\) 为有限秩\term[jubu-ziyou-ceng]{局部自由层}[locally free sheaf]。秩可在不同连通分量改变；恒为一的情形称为可逆层。
\end{definition}

\begin{corollary}\label{b3:cor:vector-bundle-projective-module}
设 \(A\) 为交换环，\(X=\operatorname{Spec}A\)。在\cref{b3:thm:affine-quasicoherent-equivalence}中，有限生成投射 \(A\) 模恰对应有限秩局部自由层。
\end{corollary}
\begin{proof}
若 \(M\) 有限生成投射，\cref{b3:thm:finite-presentation-flat-projective}使其在每点的某个基本开邻域上为有限秩自由模，式\eqref{b3:eq:module-sheaf-localization}随即给出局部自由层。反之，局部自由层显然拟凝聚，写成 \(\widetilde M\)。取有限个基本开集使其平凡，则每个 \(M_{f_i}\) 为有限秩自由模。\cref{b3:prop:finite-presentation-open-cover}先保证 \(M\) 有限展示，再由同一局部自由判据得 \(M\) 投射。
\end{proof}

纤维与茎保存的信息不同。在点 \(x\) 处，纤维是
\(\mathcal E(x)=\mathcal E_x\otimes_{\mathcal O_{X,x}}\kappa(x)\)；
茎则是邻域数据的局部环模。一个有限生成模的所有纤维维数相同，仍不能在有幂零元的基底上保证局部自由。例如 \(A=k[\varepsilon]/(\varepsilon^2)\)、\(M=A/(\varepsilon)\) 只有一个一维纤维，却不自由：若自由，维数要求它为秩一，而其非零零化子 \((\varepsilon)\) 排除了这一点。

\ProseParagraphBreak
概形本身也可以沿开子概形黏合。给定 \(X_i\)、开集 \(U_{ij}\subseteq X_i\) 和同构 \(\alpha_{ij}:U_{ij}\rightarrow U_{ji}\)，要求 \(U_{ii}=X_i\)、\(\alpha_{ii}=1\)、\(\alpha_{ji}=\alpha_{ij}^{-1}\)，并且
\[
\alpha_{ij}(U_{ij}\cap U_{ik})=U_{ji}\cap U_{jk},\qquad
\alpha_{jk}\circ\alpha_{ij}=\alpha_{ik}
\quad\text{于 }U_{ij}\cap U_{ik}.
\]
先把不交并按这些同构取商；每个 \(X_i\) 到商空间是开嵌入，因为其开集的饱和在各分量仍开。再按相容族的方法黏合环层，茎在每个局部仍是原茎。因此商空间的每点有原来的仿射邻域，得到概形。这个构造同时拼合拓扑和结构层。

\ProseParagraphBreak
局部自由层的几何描述正是这种黏合的一个应用。对开子概形 \(U\) 和非负整数 \(r\)，在各仿射图 \(\operatorname{Spec}A\subseteq U\) 上取
\(\operatorname{Spec}A[T_1,\ldots,T_r]\)。沿 \(U\) 的交叠，用系数的限制及同名变量进行黏合，得到 \(U\) 上的平凡仿射空间，记为 \(\mathbb A_U^r\)；各图上的自然投影也黏合为 \(\mathbb A_U^r\rightarrow U\)。例如在基本开集上，所用限制只是
\(A[T_1,\ldots,T_r]\rightarrow A_f[T_1,\ldots,T_r]\)。共同细化仿射覆盖表明，构造与覆盖选择无关。

\begin{definition}\label{b3:def:algebraic-vector-bundle}
设 \(X\) 为概形，\(p:V\rightarrow X\) 为概形态射。若存在开覆盖 \(X=\bigcup_iU_i\)、有限非负整数 \(r_i\)，以及 \(U_i\) 上的同构
\[
\phi_i:p^{-1}(U_i)\xrightarrow{\sim}\mathbb A_{U_i}^{r_i},
\]
使每个非空交集 \(U_i\cap U_j\) 上有 \(r_i=r_j\)，且坐标变换 \(\phi_i\phi_j^{-1}\) 形如
\(v_i=G_{ij}v_j\)，其中 \(G_{ij}\) 是该交集上正则函数组成的可逆矩阵，则称 \(p\) 及这组相容的线性平凡化为 \(X\) 上的局部有限秩\term[daishu-xiangliangcong]{代数向量丛}[algebraic vector bundle]。两组平凡化之间的过渡仍为可逆线性变换时，视为同一个向量丛结构。向量丛态射是底空间 \(X\) 上的概形态射，要求它在双方同时平凡化的开集上由正则函数矩阵给出，不含常数项；两个丛的秩不同时，矩阵可以是矩形的。
\end{definition}

矩阵条件描述的是基底环上的线性变换，不能只检查每个剩余域纤维上的点映射；在含幂零元的基底上，后者未必检测到坐标表达中的全部项。纤维秩在每个平凡图上为 \(r_i\)，允许不同连通分量有不同秩。

\begin{proposition}\label{b3:prop:locally-free-vector-bundle-equivalence}
设 \(X\) 为概形。局部有限秩代数向量丛与有限秩局部自由 \(\mathcal O_X\) 模层互相对应，且这种对应保持态射的方向。具体地，向量丛 \(p:V\rightarrow X\) 的截面层定义为
\[
\mathcal S_V(U)=\{\,s:U\rightarrow V\text{ 为概形态射}\mid
p\circ s\text{ 是 }U\hookrightarrow X\,\}.
\]
由平凡化中的坐标加法和标量乘法，\(\mathcal S_V\) 成为有限秩局部自由模层；给定局部自由层 \(\mathcal E\)，可以构造截面层同构于 \(\mathcal E\) 的向量丛。两种构造互为逆，精确到与所给局部识别相容的同构。
\end{proposition}
\begin{proof}
在仿射平凡图 \(U_i=\operatorname{Spec}A_i\) 上，\(\mathbb A_{U_i}^{r_i}\) 的截面由保持 \(A_i\) 的环同态
\(A_i[T_1,\ldots,T_{r_i}]\rightarrow A_i\) 给出，恰由各 \(T_j\) 的像决定。因此截面就是 \(r_i\) 个正则函数。进一步限制到开集并黏合，得到 \(\mathcal S_V|_{U_i}\cong\mathcal O_{U_i}^{r_i}\)。线性过渡矩阵保持这些坐标的加法和标量乘法，使它们黏合为模层结构；概形态射的局部黏合也给出截面层的层公理。

反过来，取 \(\mathcal E\) 的仿射平凡覆盖，令 \(e_i\) 为第 \(i\) 个图上的一行基截面。在交集上写
\(e_j=e_iG_{ij}\)，则 \(G_{ij}G_{jk}=G_{ik}\)。以
\(v_i=G_{ij}v_j\) 把 \(\mathbb A_{U_j}^{r_j}\) 与 \(\mathbb A_{U_i}^{r_i}\) 的对应开子概形识别，余圈关系保证能黏合为 \(V\rightarrow X\)。一个相容坐标族 \((v_i)\) 同时给出该向量丛的截面和层 \(\mathcal E\) 的截面 \(e_iv_i\)，从而得到自然同构 \(\mathcal S_V\cong\mathcal E\)。更换基或细化覆盖，只改变这些局部表示，相容的局部同构给出总空间的同构。从原向量丛的截面层出发，取原平凡化对应的基，所还原的过渡矩阵也与原矩阵相同，因而恢复原向量丛。

最后，在同时平凡化的图上，模层态射由矩阵 \(H_i\) 给出；它在交集上的相容条件恰为
\(H_iG_{ij}=G'_{ij}H_j\)。同一矩阵定义从第一向量丛到第二向量丛的坐标映射，并由这个等式黏合为向量丛态射。反向把向量丛态射与截面复合，得到同一个矩阵的模层态射。局部矩阵也保持恒等映射与复合，故两个对应在态射上互为逆。
\end{proof}

这里的总空间在平凡图上的坐标代数是 \(\operatorname{Sym}(\mathcal E^\vee)\)，其中对偶层由
\(\mathcal E^\vee(U)=\operatorname{Hom}_{\mathcal O_U}(\mathcal E|_U,\mathcal O_U)\) 给出，所以其截面层是 \(\mathcal E\)。固定秩向量丛的局部平凡化与截面层构造见 Hartshorne\cite[Chapter II, Exercise 5.18, pp. 128--129]{Hartshorne1977}；该书对 \(\operatorname{Sym}(\mathcal E)\) 的总空间记法给出的截面层为 \(\mathcal E^\vee\)。下面射影直线上的非平凡可逆层将具体说明，局部基未必能够黏合为全局基。

\subsection{局部化相容性与黏合}
\label{b3:subsec:2402:03}

\begin{proposition}\label{b3:prop:sheaf-gluing}
设 \(X\) 为概形，\(X=\bigcup_iU_i\) 为开覆盖。给定各 \(U_i\) 上的 \(\mathcal O_{U_i}\) 模层 \(\mathcal F_i\) 及同构
\(\theta_{ij}:\mathcal F_j|_{U_i\cap U_j}\rightarrow\mathcal F_i|_{U_i\cap U_j}\)。若 \(\theta_{ii}=1\)，且 \(\theta_{ij}\theta_{jk}=\theta_{ik}\) 在每个三重交集上成立，则存在模层 \(\mathcal F\) 及相容同构 \(\mathcal F|_{U_i}\cong\mathcal F_i\)，并在保留这些同构的意义下唯一。若各 \(\mathcal F_i\) 拟凝聚或局部自由，则 \(\mathcal F\) 也分别具有该性质。
\end{proposition}
\begin{proof}
对开集 \(V\subseteq X\)，定义 \(\mathcal F(V)\) 为全部族
\[
s_i\in\mathcal F_i(V\cap U_i),\qquad
s_i=\theta_{ij}(s_j)\ \text{于 }V\cap U_i\cap U_j.
\]
限制和标量作用逐分量给出。局部相容族先在每个 \(U_i\) 内由层公理黏合，所得截面仍满足交集关系，因而 \(\mathcal F\) 是层。在 \(V\subseteq U_i\) 时，第 \(i\) 分量决定其余分量为 \(\theta_{ji}(s_i)\)；余圈关系保证相容，遂得所需同构。任何其他实现都按这些局部同构给出同一个相容族，得到唯一性。最后两项性质本来就是开局部条件。
\end{proof}

取两份 \(\operatorname{Spec}k[t]\)，沿 \(D(t)\) 用恒等映射黏合，得到带双原点的仿射直线。它的全局函数是
\[
\{(a(t),b(t))\in k[t]^2:a=b\text{ 于 }k[t,t^{-1}]\}
\cong k[t].
\]
若此概形仿射，全局函数环将由\cref{b3:thm:affine-scheme-hom}恢复整个概形；但两个原点在全局函数上的取值理想均为 \((t)\)，不可能在同一个素谱中成为两个不同点。故它不是仿射概形。这个例子说明，知道各仿射块和全局函数仍不足以丢掉黏合数据。

% !TeX root = ../../Book3.tex
% 纲要标题：### 24.3 Proj 构造
% 纲要位置：计划纲要.md，第 2748 行。

\section{Proj 构造}
\label{b3:sec:2403}

齐次坐标表示一个方向，局部图中的函数因此必须由相同次数的齐次元相除得到。取局部化环的零次部分实现这一要求，扭曲层则记录不同次数之间的转换。


% 纲要内容（仅作写作计划，尚非教材正文）：
%
% * 无关理想与标准图；先计算次数一例，再证明一般齐次素理想的零次对应
% * 射影厚点保留幂零结构；饱和化是各标准图理想的收缩交，无关挠子例正式计算
% * 扭曲层的移位符号与射影直线过渡；加权反例说明标准分次条件不能省略
% * Spec 与 Proj 的信息区别；后者还取局部零次部分，不能只说丢掉无关挠子
%
% ---
%

\subsection{齐次素理想}
\label{b3:subsec:2403:01}

若方程齐次，缩放坐标不会改变其零点条件。对域 \(k\) 上的标准分次环 \(k[x_0,\ldots,x_r]\)，正次数理想 \((x_0,\ldots,x_r)\) 对应仿射锥的原点，却不表示一个非零向量的方向。射影构造因此选取齐次素理想，并排除包含全部正次数元素的那些素理想。对一般正分次环 \(S\)，将其正次数理想记为 \(S_+\)；若 \(S_+\subseteq\mathfrak p\)，每个正次数齐次元都落在 \(\mathfrak p\) 中，这个点就不会出现在任何射影标准图上。

\begin{definition}\label{b3:def:proj}
设 \(S=\bigoplus_{n\ge0}S_n\) 为交换正分次环，\(S_+=\bigoplus_{n>0}S_n\)。定义
\[
\operatorname{Proj}S
=\{\,\mathfrak p\subset S:\mathfrak p\;\text{为齐次素理想且}\;
S_+\not\subseteq\mathfrak p\,\}.
\]
对齐次理想 \(I\)，令 \(V_+(I)=\{\mathfrak p:I\subseteq\mathfrak p\}\)，以这些集合为闭集。对正次数齐次元 \(f\)，记 \(D_+(f)=\operatorname{Proj}S\setminus V_+((f))\)。
\end{definition}
\symidx{\(\operatorname{Proj}S\)：分次环 \(S\) 的射影谱}
\symidx{\(V_+(I)\)、\(D_+(f)\)：射影谱中的闭集与基本开集}

零理想与单位理想分别给全空间和空集；理想和及乘积分别对应闭集的任意交和有限并，证明与素谱相同。各 \(D_+(f)\) 构成开集基，且 \(D_+(f)\cap D_+(g)=D_+(fg)\)。排除 \(S_+\) 的条件恰好保证每个点至少属于其中一个基本开集。

\ProseParagraphBreak
设 \(k\) 为域。在 \(k[x_0,\ldots,x_r]\) 的标准分次下，非零向量 \(a=(a_0,\ldots,a_r)\in k^{r+1}\) 给出齐次素理想
\[
\ker\bigl(k[x_0,\ldots,x_r]\longrightarrow k[t],\
x_i\longmapsto a_it\bigr).
\]
同比例向量给相同理想；选择 \(a_i\ne0\) 后，坐标比 \(a_j/a_i\) 又确定该点。这恢复通常的 \(k\) 值射影点。Proj 还包含非闭点及可能具有更大剩余域的闭点，不能只用 \(k\) 中向量穷尽一般域上的全部点。

\subsection{标准开集与射影空间}
\label{b3:subsec:2403:02}

先看次数一的模型。设 \(k\) 为域，在 \(S=k[x_0,x_1]\) 中倒置 \(x_0\) 后，令 \(t=x_1/x_0\)，则
\[
S_{x_0}=k[t][x_0,x_0^{-1}],\qquad (S_{x_0})_0=k[t].
\]
齐次方向在这张图上由坐标比 \(t\) 表示，零次环给出的就是一条仿射直线。更一般地，倒置次数一的齐次元 \(f\) 后，每个次数 \(e\) 的分量都等于 \(f^e(S_f)_0\)，齐次素理想因此由其零次部分确定。对于 \(\deg f=d>1\)，不一定能直接将每个齐次元变成零次元，但先取 \(d\) 次幂就能做到；下面的证明据此将齐次元分成单位与幂零元两种情形。

\begin{lemma}[齐次素理想的零次对应]\label{b3:lem:proj-homogeneous-prime-correspondence}
设 \(S=\bigoplus_{n\ge0}S_n\) 为交换含幺正分次环，\(f\in S_d\)、\(d>0\) 为齐次元。记 \(F=S_f\)、\(B=F_0\)。则
\[
\mathfrak p\longmapsto(\mathfrak pS_f)_0
\]
给出 \(S\) 中不含 \(f\) 的齐次素理想与 \(B\) 的素理想之间的双射。
\end{lemma}
\begin{proof}
局部化把 \(S\) 中不含 \(f\) 的齐次素理想与 \(F\) 的齐次素理想对应；\(F\) 具有 \(\mathbb Z\) 分次，且 \(f\) 是次数 \(d\) 的单位。对 \(\mathfrak q\in\operatorname{Spec}B\)，零次元素的局部化和取商仍保持分次，所以在 \(B\) 中收缩为 \(\mathfrak q\) 的齐次素理想，与
\[
C=F\otimes_B\kappa(\mathfrak q)\cong F_{B\setminus\mathfrak q}/\mathfrak qF_{B\setminus\mathfrak q}
\]
的齐次素理想对应。这里只需证明 \(C\) 恰有一个齐次素理想。其零次部分为域 \(\kappa(\mathfrak q)\)，而 \(f\) 仍为次数 \(d\) 的单位。

若 \(c\in C_e\) 齐次，则 \(c^df^{-e}\in C_0\)。若这个零次元非零，它是单位，从而 \(c\) 也是单位；若为零，则 \(c^d=0\)。因此每个齐次元或者可逆，或者幂零。令 \(J\) 为齐次幂零元生成的理想。我们有
\[
J=\sqrt{(0)},
\]
即 \(C\) 的幂零根。确实，\(J\subseteq\sqrt{(0)}\)；反过来，若 \(z\) 幂零，则 \(z\) 的最高非零齐次分量也幂零，因为 \(z^N\) 的最高次分量是该分量的 \(N\) 次幂。减去这个分量后仍为幂零元；对有限个齐次分量逐次重复，便得到 \(z\in J\)。这个论证不要求环 Noether。

\(J\) 是齐次理想，且 \(J\cap C_0=0\)，所以 \(C/J\ne0\)。在 \(C/J\) 中，每个非零齐次元都是单位：这样的分量可取齐次代表，而不在 \(J\) 中的齐次元不能幂零，故可逆。于是两个非零元素的最高非零齐次分量之积非零，商环是整环。因而 \(J\) 是素理想。任何齐次素理想都包含幂零根，又不能包含齐次单位；上述二择性遂迫使它等于 \(J\)。将 \(J\) 沿剩余域映射和局部化收缩，就得到唯一的、在 \(B\) 中收缩为 \(\mathfrak q\) 的齐次素理想。
\end{proof}

\begin{theorem}\label{b3:thm:proj-affine-charts}\label{b3:thm:fat-projective-point}
设 \(S\) 为交换正分次环，\(f\in S_d\)、\(d>0\) 为齐次元，记 \(S_{(f)}=(S_f)_0\)。则
\[
D_+(f)\cong\operatorname{Spec}S_{(f)}.
\]
这些仿射结构在交集上相容，因而唯一确定 \(\operatorname{Proj}S\) 的概形结构。这里 \(S_{(f)}\) 由分式 \(a/f^m\)、\(a\in S_{md}\) 组成，分子与分母的次数相同。

具体地，若 \(k\) 为域，\(R=k[x,y]/(x^2)\) 取标准分次，则
\[
\operatorname{Proj}R\cong\operatorname{Spec}\bigl(k[t]/(t^2)\bigr),
\qquad t=x/y.
\]
其底空间只有一个点，全局截面为 \(k[t]/(t^2)\)，这个局部 Artin 环的长度为二。
\end{theorem}
\begin{proof}
由\cref{b3:lem:proj-homogeneous-prime-correspondence}，映射
\[
\mathfrak p\longmapsto(\mathfrak pS_f)_0
\]
给出 \(D_+(f)\) 与 \(\operatorname{Spec}S_{(f)}\) 的双射。若 \(g\in S_e\)、\(e>0\) 齐次，则 \(g\notin\mathfrak p\) 等价于 \(g^d/f^e\notin(\mathfrak pS_f)_0\)。于是 \(D_+(f)\) 中的基本开集 \(D_+(fg)\) 对应 \(S_{(f)}\) 的基本开集 \(D(g^d/f^e)\)。这些集合也给出右侧的开集基：对任意 \(a/f^m\in S_{(f)}\)，若 \(m>0\)，相应的 \(a^d/f^{md}\) 正是 \((a/f^m)^d\)，二者定义同一个基本开集；若 \(m=0\)，可把 \(a\) 改写为 \((af)/f\)。因此得到同胚。

在交集上，局部化环满足
\[
(S_{(f)})_{g^d/f^e}\cong(S_{fg})_0.
\]
确实，在 \(S_f\) 中把 \(g^d/f^e\) 变成单位等价于把 \(g\) 变成单位；被倒置的元素次数为零，取零次部分与这次局部化交换。同样从 \(g\) 一侧得到相同环。三重交集上的映射都是分式的限制，故满足黏合条件。上一节的概形黏合构造给出结构层，而它在这些仿射块上的指定限制决定其唯一性。

在所给的 \(R\) 中，\(x\) 幂零，所以每个素理想都包含 \(x\)。避开正次数理想 \((x,y)\) 就意味着不含 \(y\)，故 \(D_+(y)\) 是整个 Proj。局部化后 \(R_y=k[y,y^{-1}][x]/(x^2)\)，其零次部分以 \(1,x/y\) 为 \(k\) 基，乘法关系为 \((x/y)^2=0\)。因此这一张图便给出所列的概形与全局截面。其唯一素理想为 \((t)\)，而链 \(0\subset(t)\subset k[t]/(t^2)\) 的两个因子均为剩余域 \(k\)，所以长度为二。一个射影点的底空间仍可承载非零幂零结构。
\end{proof}

Proj 的图上函数取自 \((S_f)_0\)：相同次数的分子、分母之比保留方向信息，而整个 \(S_f\) 还保留非零次数的齐次坐标。上例也表明，射影谱的概形结构包含点集计数看不见的厚度。一般构造可对照\cite[Chapter II, Proposition 2.5, pp. 76--77]{Hartshorne1977}。

\begin{definition}\label{b3:def:projective-space}
设 \(A\) 为交换环，\(r\ge0\) 为整数，多项式环 \(A[x_0,\ldots,x_r]\) 中各变量次数为一。概形
\[
\mathbb P_A^r=\operatorname{Proj}A[x_0,\ldots,x_r]
\]
称为 \(A\) 上的 \(r\) 维\term[sheying-kongjian]{射影空间}[projective space]。
\end{definition}
\symidx{\(\mathbb P_A^r\)：环 \(A\) 上的 \(r\) 维射影空间}

它由 \(r+1\) 个仿射开集覆盖：
\[
D_+(x_i)=\operatorname{Spec}
A[x_0/x_i,\ldots,\widehat{x_i/x_i},\ldots,x_r/x_i]
\cong\mathbb A_A^r.
\]
这些分式代数独立，因为齐次化把它们的多项式关系还原为原变量的关系。对 \(\mathbb P_k^1\)，两张图为 \(k[t]\)、\(k[u]\)，交集为 \(k[t,t^{-1}]\)，黏合规则 \(u=t^{-1}\)。这与双原点直线的恒等黏合不同：射影直线增加的是另一端的点，而非再复制原点。

\ProseParagraphBreak
若 \(I\subseteq S\) 为齐次理想，则 \(\operatorname{Proj}(S/I)\) 在每张图上由
\((S_f)_0\rightarrow((S/I)_f)_0\) 给出商环，因此是 \(\operatorname{Proj}S\) 的闭子概形。齐次方程由此不仅规定闭点的零集，也规定其上的商结构。与仿射闭子概形不同，同一个射影闭子概形可能由不同的齐次理想给出：被无关理想 \(S_+\) 的某个幂零化的分次部分，在每张 \(D_+(f)\) 上都会消失。在有限生成的标准分次下，这个差异可由饱和化精确描述。

\ProseParagraphBreak
例如，对域 \(k\)，标准分次环 \(k[x,y]\) 的两个理想 \((x^2,xy)\) 与 \((x)\) 不同，但倒置 \(x\) 后都变成单位理想，倒置 \(y\) 后又都由 \(x\) 生成。两张射影图看到相同的局部理想；饱和化将这些局部条件收缩回原环，下面的命题算出其共同结果为 \((x)\)。

\begin{proposition}[射影闭子概形与饱和化]\label{b3:prop:proj-saturation-standard}\label{b3:prop:proj-saturation-example}
设 \(S\) 为交换含幺正分次环，由有限个次数一元素 \(x_0,\ldots,x_r\) 作为 \(S_0\) 代数生成，并记 \(S_+=(x_0,\ldots,x_r)\)。对齐次理想 \(I\subseteq S\)，定义其 \(S_+\)-饱和化为
\[
I^{\mathrm{sat}}=I:S_+^\infty
=\bigcup_{N\ge0}\{\,a\in S:S_+^Na\subseteq I\,\}.
\]
则 \(I\) 与 \(I^{\mathrm{sat}}\) 给出 \(\operatorname{Proj}S\) 内同一个闭子概形。两个齐次理想 \(I,J\) 给出其中同一个闭子概形，当且仅当 \(I^{\mathrm{sat}}=J^{\mathrm{sat}}\)。

例如，对域 \(k\) 上的标准分次环 \(S=k[x,y]\)，取 \(I=(x^2,xy)\)、\(J=(x)\)。则 \(I\ne J\)，而 \(I^{\mathrm{sat}}=J\)，两者给出同一个射影闭子概形。差模 \(J/I\) 由 \(x\) 的类生成，且被 \(S_+=(x,y)\) 零化，故其射影伴随层为零。
\end{proposition}
\symidx{\(I^{\mathrm{sat}}\)：齐次理想关于无关理想的饱和化}
\begin{proof}
先对任意齐次理想 \(I\) 说明
\[
I^{\mathrm{sat}}
=\bigcap_{i=0}^r(I:x_i^\infty).
\]
这里 \(I:x_i^\infty\) 正是局部理想 \(I_{x_i}\) 收缩回 \(S\) 的结果，所以这条公式将每张标准图看到的理想条件重新求交。左侧的元素显然属于右侧。若 \(a\) 属于右侧，分别取正整数 \(N_i\)，使 \(x_i^{N_i}a\in I\)。取 \(N>\sum_i(N_i-1)\)，则每个由 \(x_0,\ldots,x_r\) 组成的 \(N\) 次乘积都含某个 \(x_i^{N_i}\)，所以 \(S_+^Na\subseteq I\)。这证明反向包含；有限个标准坐标正用于把逐图的零化指数合成一个共同指数。

若 \(a\in I^{\mathrm{sat}}\)，则在每个 \(S_{x_i}\) 中有 \(a\in I_{x_i}\)，因而 \((I^{\mathrm{sat}})_{x_i}=I_{x_i}\)。取零次部分后，两理想在每张图上给出相同的商环，故给出同一个闭子概形。

若 \(I,J\) 给出同一个闭子概形，则每张图中的局部理想 \((I_{x_i})_0\) 与 \((J_{x_i})_0\) 相同。对齐次元 \(a\in I_e\)，次数一条件使 \(a/x_i^e\) 属于 \((I_{x_i})_0\)，于是 \(a\in J_{x_i}\)。对所有齐次分量应用此论证，得到 \(I_{x_i}\subseteq J_{x_i}\)；交换两理想可得相等。由刚才的收缩公式，\(I^{\mathrm{sat}}=J^{\mathrm{sat}}\)。反向由二者均与自己的饱和化给出相同闭子概形立即得到。

对最后的例子，\(x\notin(x^2,xy)\)，故 \(I\ne J\)，但 \(S_+x\subseteq I\)，于是 \(J\subseteq I^{\mathrm{sat}}\)。另一方面，\(J\) 已饱和：若 \(S_+^Na\subseteq J\)，模掉 \(x\) 后有 \(y^N\bar a=0\) 于整环 \(k[y]\)，故 \(a\in J\)。由 \(I\subseteq J\)，得到 \(I^{\mathrm{sat}}\subseteq J^{\mathrm{sat}}=J\)，所以 \(I^{\mathrm{sat}}=J\)。逐图看，倒置 \(x\) 后二者都变成单位理想，倒置 \(y\) 后二者都由 \(x\) 生成。差模 \(J/I\) 中 \(x\) 的类同时被 \(x,y\) 零化，在两张图上都局部化为零，因此其伴随层消失。
\end{proof}

\subsection{分次模与射影几何}
\label{b3:subsec:2403:03}

设 \(M\) 为分次 \(S\) 模。在 \(D_+(f)\) 上把它取为 \(S_{(f)}\) 模 \((M_f)_0\) 的伴随层；上一定理中的限制同构对模也成立，所以这些层黏合为 \(\operatorname{Proj}S\) 上的拟凝聚层，仍记 \(\widetilde M\)。当同时讨论 Spec 与 Proj 时，必须注明所在空间：前者使用全部 \(M_f\)，后者仅用其次数零部分。

\begin{definition}\label{b3:def:twisting-sheaf}
设 \(S\) 为交换正分次环，\(X=\operatorname{Proj}S\)，\(n\in\mathbb Z\)。令 \(S(n)\) 为次数平移后的分次 \(S\) 模，即 \(S(n)_j=S_{n+j}\)。对每个正次数齐次元 \(f\in S\)，在仿射开集 \(D_+(f)=\operatorname{Spec}(S_f)_0\) 上取与 \((S(n)_f)_0\) 伴随的层；这些局部层黏合为 \(X\) 上的层，定义
\[
\mathcal O_X(n)=\widetilde{S(n)}.
\]
\(\mathcal O_X(1)\) 称为 Serre 的\term[niuqu-ceng]{扭曲层}[twisting sheaf]；对 \(X\) 上的 \(\mathcal O_X\) 模层 \(\mathcal F\)，记 \(\mathcal F(n)=\mathcal F\otimes_{\mathcal O_X}\mathcal O_X(n)\)。
\end{definition}
\symidx{\(\mathcal O_X(n)\)、\(\mathcal F(n)\)：扭层与模层的扭转}

本书的平移 \(S(n)_j=S_{n+j}\) 使 \((S(n)_f)_0=(S_f)_n\)。在次数一元素 \(f\) 的图上，\(\mathcal O(-1)\) 因而以 \(f^{-1}\) 为基，\(\mathcal O(1)\) 则以 \(f\) 为基。下面的标准分次条件使这些图覆盖整个 Proj，并保证两层互为对偶。

\begin{proposition}\label{b3:prop:standard-twists-invertible}\label{b3:prop:weighted-twist-counterexample}
设 \(S\) 为交换正分次环，且由 \(S_1\) 作为 \(S_0\) 代数生成，\(X=\operatorname{Proj}S\)。则对任意整数 \(n,a,b\) 及分次 \(S\) 模 \(M\)，\(\mathcal O_X(n)\) 为可逆层，且有
\[
\widetilde{M(n)}\cong\widetilde M(n),\qquad
\mathcal O_X(a)\otimes\mathcal O_X(b)\cong\mathcal O_X(a+b).
\]
特别地，\(\mathcal O_X(-1)^\vee\cong\mathcal O_X(1)\)。次数一生成条件不能省略：若 \(k\) 为域，\(S'=k[x]\) 取 \(\deg x=2\)，\(X'=\operatorname{Proj}S'\)，则
\[
X'\cong\operatorname{Spec}k,\qquad
\mathcal O_{X'}(1)=0,\qquad
\mathcal O_{X'}(2)\cong\mathcal O_{X'}.
\]
\end{proposition}
\begin{proof}
次数一元素的 \(D_+(f)\) 覆盖 \(X\)。在这种图上，\(S(n)_{(f)}=(S_f)_n\)，乘以可逆元 \(f^n\) 给出 \(S_{(f)}\cong(S_f)_n\)，故局部秩为一。相同映射给
\((M_f)_0\otimes_{S_{(f)}}(S_f)_n\cong(M_f)_n\)。
这些同构在交集上都由乘法给出，所以黏合。取 \(M=S(a)\) 得最后一式。

对 \(S'\) 的例子，只有齐次素理想 \((0)\) 避开 \(S'_+=(x)\)，所以 \(D_+(x)=X'\)。有 \(S'_x=k[x,x^{-1}]\)，其零次部分为 \(k\)，奇数次部分全部为零，次数二部分则为 \(kx\)。应用扭曲层定义，得到 \(X'\)、\(\mathcal O_{X'}(1)\)、\(\mathcal O_{X'}(2)\) 的三项结论。零层在这个非空概形上不可能可逆；而 \(\mathcal O(1)\otimes\mathcal O(1)=0\) 也不能同构于 \(\mathcal O(2)\)。
\end{proof}

设 \(k\) 为域，\(n\in\mathbb Z\)。在 \(\mathbb P_k^1\) 上，令 \(t=x_1/x_0\)、\(u=x_0/x_1\)。\cref{b3:tab:projective-line-twist-charts}列出 \(\mathcal O(n)\) 在两张仿射图中的局部基，及它们在交集上的关系。

\begin{table}[htbp]
\centering
\caption{射影直线的两图与扭曲层局部基}
\label{b3:tab:projective-line-twist-charts}
\begin{tabularx}{\linewidth}{@{}lXX@{}}
\toprule
开集 & 坐标环 & \(\mathcal O(n)\) 的局部基与过渡 \\
\midrule
\(U_0=D_+(x_0)\) & \(k[t]\) & \(e_0=x_0^n\) \\
\(U_1=D_+(x_1)\) & \(k[u]\)，\(u=t^{-1}\) & \(e_1=x_1^n\) \\
\(U_0\cap U_1\) & \(k[t,t^{-1}]\) & \(e_1=t^ne_0\) \\
\bottomrule
\end{tabularx}
\end{table}

这些局部基是线丛纤维的基，不是两张图上必须取同一数值的普通函数。对 \(\mathcal O(-1)\)，在 \(U_0\) 上把基 \(e_0=x_0^{-1}\) 识别为向量 \((1,t)\)，在 \(U_1\) 上把 \(e_1=x_1^{-1}\) 识别为 \((u,1)\)。交叠上恰有
\[
(u,1)=t^{-1}(1,t),\qquad e_1=t^{-1}e_0,
\]
所以这两个局部自由秩一子模黏合为平凡秩二向量丛中的一个线丛。在 \(k\) 有理点 \([a:b]\) 上，它的纤维就是 \(k(a,b)\subseteq k^2\)；在任意点也可用该点的剩余域作同样描述。每个点所表示的直线就是这个线丛的纤维。负扭曲因此有具体的几何对象，转换函数记录的是同一条直线的两种选基。

\ProseParagraphBreak
以交集上的 \(e_0\) 为基，局部截面 \(a(t)e_0\) 与 \(b(t^{-1})e_1\) 的相容条件是 \(a(t)=t^n b(t^{-1})\)。因此全局截面是下列映射的核：
\begin{equation}\label{b3:eq:projective-line-twist-cech-map}
\begin{aligned}
\delta_n:k[t]\oplus k[t^{-1}]&\longrightarrow k[t,t^{-1}],\\
(a,b)&\longmapsto a-t^n b.
\end{aligned}
\end{equation}
当 \(n\ge0\) 时，\(a\) 正是次数不超过 \(n\) 的多项式，截面空间的维数为 \(n+1\)；当 \(n<0\) 时只有零截面。特别地，\(\mathcal O(-1)\) 在每点局部自由，却没有非零全局截面，而结构层至少有非零截面 \(1\)，故两者不可能同构。这是在概形上看到的局部秩一自由与全局平凡性的区别；\cref{b3:cor:vector-bundle-projective-module}将仿射可逆层对应到秩一投射模，在 Noether 正规整环情形便是\cref{b3:prop:picard-class-group}中的 Picard 群对象。层上同调还将解释同一映射 \(\delta_n\) 的余核。

\ProseParagraphBreak
与仿射情形不同，\(M\mapsto\widetilde M\) 在 Proj 上不能由不带次数的全局截面反演。这一步先局部化，零化那些在标准图上消失的部分，再只取局部模的零次分量。例如 \(S/S_+\) 在每个 \(D_+(f)\) 上局部化为零，故其伴随层为零；当 \(S_0\ne0\) 时，这个分次模本身却不为零。零次筛选也会丢失信息：在加权例子的 \(S'\) 中，\(S'(1)\) 没有被 \((x)\) 的幂零化的非零元素，但 \(\widetilde{S'(1)}=0\)。因此一般分次下不能把全部信息损失归结为无关理想挠子。\Cref{b3:tab:spec-proj-sheaf-comparison}把这些运算与闭子概形的饱和化放在一起比较；仿射列中的 \(M\) 是 \(A\) 模，射影列中的 \(M\) 是分次 \(S\) 模。扭曲层的局部计算见\cite[Chapter II, Propositions 5.11--5.12, pp. 116--117]{Hartshorne1977}。

\begin{table}[htbp]
\centering
\caption[Spec 与 Proj 的模层比较]{Spec 与 Proj 的模层比较。射影列中关于理想的结论假设 \(S\) 为有限生成标准分次 \(S_0\) 代数。}
\label{b3:tab:spec-proj-sheaf-comparison}
\begin{tabularx}{\linewidth}{@{}lXX@{}}
\toprule
比较对象 & \(\operatorname{Spec}A\) & \(\operatorname{Proj}S\) \\
\midrule
基本开集上的模 & \(M_f\) & \((M_f)_0\) \\
闭子概形的理想 & 确定原理想 \(I\subseteq A\) & 确定 \(I^{\mathrm{sat}}\)，未必确定 \(I\) \\
全局截面 & \(\Gamma(\operatorname{Spec}A,\widetilde M)\cong M\) & 一般不能恢复 \(M\)；\(S/S_+\) 对应零层 \\
\bottomrule
\end{tabularx}
\end{table}

% !TeX root = ../../Book3.tex
% 纲要标题：### 24.4 纤维积、分离性与固有性
% 纲要位置：计划纲要.md，第 2755 行。

\section{纤维积、分离性与固有性}
\label{b3:sec:2404}

闭集在投影下未必仍闭：双曲线 \(xy=1\) 向 \(x\) 轴的投影漏掉原点。射影空间为何能补足这种缺陷？这个问题需要对所有基变换同时考察；纤维积给出改变基底的几何对象，分离性要求对角在相对乘积中成为闭子概形，而固有性还要求有限型及任意基变换后的闭性。


% 纲要内容（仅作写作计划，尚非教材正文）：
%
% * 纤维积的泛性质图、仿射目标由全局函数控制、存在性与张量积模型
% * 概形纤维的底空间与点逆像相同，结构层还保留厚度和剩余域；平方映射的分歧纤维
% * 仿射局部定义的平展态射及其换基、复合稳定性；开浸入与实数到复数的例子
% * 闭浸入、分离、有限型及固有性分别说明；普遍闭性不能由映射自身闭代替
% * 任意基底的射影固有性：空纤维、高次齐次分量在邻域消失、无关理想幂零，以及乘积图中的对角商环
% * Hartshorne 赋值判据与射影固有性引文的 Noether 范围
%
% ---
%

\subsection{纤维积与分离性}
\label{b3:subsec:2404:01}

前面构造了概形和具体黏合。进一步研究态射时，需要同时比较一个空间与所有可能的基变换。为此先说明纤维积；它在仿射环上就是已经熟悉的张量积。

\begin{definition}\label{b3:def:scheme-fiber-product}\label{b3:def:scheme-theoretic-fiber}
设 \(X,Y,T\) 为概形，并给定态射 \(f:X\to T\)、\(g:Y\to T\)。若概形 \(P\) 带有投影 \(p_X:P\to X\)、\(p_Y:P\to Y\)，满足 \(fp_X=gp_Y\)，并且对任意概形 \(Z\) 及态射 \(a:Z\to X\)、\(b:Z\to Y\)，只要 \(fa=gb\)，就存在唯一态射 \(h:Z\to P\) 满足 \(p_Xh=a\)、\(p_Yh=b\)，则称 \(P\) 为\term[xianwei-ji]{纤维积}[fiber product]，记为 \(X\times_TY\)。其泛性质由下图表示：
\[
\begin{tikzcd}[column sep=large,row sep=large]
Z \arrow[drr,bend left=18,"a"]
  \arrow[ddr,bend right=18,"b"']
  \arrow[dr,dashed,"{\exists!\,h}"] & & \\
& P \arrow[r,"p_X"] \arrow[d,"p_Y"']
  \arrow[dr,phantom,"\lrcorner",very near start]
& X \arrow[d,"f"] \\
& Y \arrow[r,"g"'] & T
\end{tikzcd}
\]
等价地，若 \(\operatorname{Hom}\) 表示概形态射集合，对每个 \(Z\) 都有自然双射
\[
\operatorname{Hom}(Z,P)
\cong\operatorname{Hom}(Z,X)
\times_{\operatorname{Hom}(Z,T)}\operatorname{Hom}(Z,Y).
\]
投影 \(X\times_TY\rightarrow Y\) 称为 \(X\rightarrow T\) 沿 \(Y\rightarrow T\) 的基变换。

对概形态射 \(u:X\to Y\) 及点 \(y\in Y\)，由剩余域给出的自然态射 \(\operatorname{Spec}\kappa(y)\to Y\) 作基变换，所得 \(\kappa(y)\) 上概形
\[
X_y=X\times_Y\operatorname{Spec}\kappa(y)
\]
称为 \(u\) 在 \(y\) 处的\term[gaixing-xianwei]{概形纤维}[scheme-theoretic fiber]。
\end{definition}
\symidx{\(X\times_TY\)：概形在 \(T\) 上的纤维积}
\symidx{\(X_y\)：态射在点 \(y\) 处的概形纤维}

\begin{proposition}\label{b3:prop:scheme-fiber-products-exist}\label{b3:prop:scheme-hom-affine-target}\label{b3:prop:scheme-ramified-fibers}
对任意概形 \(Z\) 与交换环 \(A\)，取全局截面给出自然双射
\[
\operatorname{Hom}_{\mathrm{Sch}}(Z,\operatorname{Spec}A)
\cong\operatorname{Hom}_{\mathrm{Ring}}(A,\Gamma(Z,\mathcal O_Z)).
\]
概形的纤维积存在，并在保留投影的意义下唯一。对交换环同态 \(A\rightarrow B\)、\(A\rightarrow C\)，有
\[
\operatorname{Spec}B\times_{\operatorname{Spec}A}\operatorname{Spec}C
\cong\operatorname{Spec}(B\otimes_AC).
\]
具体地，设 \(k\) 为代数闭域，\(\operatorname{char}k\ne2\)，令 \(u:\operatorname{Spec}k[t]\to\operatorname{Spec}k[x]\) 对应于 \(x\mapsto t^2\)。对任意 \(a\in k\)，闭点 \((x-a)\) 上的概形纤维为 \(\operatorname{Spec}\bigl(k[t]/(t^2-a)\bigr)\)。\(a=0\) 时它为一个长度二的非约化点，\(a\ne0\) 时它为两个约化点；各闭纤维的总长度均为二。
\end{proposition}
\begin{proof}
先证明到仿射目标的态射公式。任一 \(v:Z\to\operatorname{Spec}A\) 的层同态在全局截面上给出 \(A\to\Gamma(Z,\mathcal O_Z)\)。反之，给定 \(\alpha:A\to\Gamma(Z,\mathcal O_Z)\)，选 \(Z\) 的仿射开覆盖 \(U_i=\operatorname{Spec}R_i\)。将 \(\alpha\) 与截面限制复合，得到 \(A\to R_i\)，再用\cref{b3:thm:affine-scheme-hom}得到 \(v_i:U_i\to\operatorname{Spec}A\)。在 \(U_i\cap U_j\) 的任一更小仿射开集上，两份环同态都来自 \(\alpha\) 的同一个限制，故 \(v_i,v_j\) 相同。它们于是唯一黏合为 \(v\)。在每张仿射图上使用原来的双射，可见这两种构造互逆；限制与复合相容，也证明关于 \(Z,A\) 的自然性。

现在取 \(T=\operatorname{Spec}A\)、\(X=\operatorname{Spec}B\)、\(Y=\operatorname{Spec}C\)。对任意概形 \(Z\)，两份到 \(X,Y\) 的态射若在 \(T\) 上相同，便对应于两份到 \(\Gamma(Z,\mathcal O_Z)\) 的环同态在 \(A\) 上相同。张量积的泛性质将它们唯一合成为 \(B\otimes_AC\to\Gamma(Z,\mathcal O_Z)\)。刚证的公式再将它转成唯一的 \(Z\to\operatorname{Spec}(B\otimes_AC)\)，证明仿射纤维积公式，包括任意 \(Z\) 的检验。

若 \(X\times_TY\) 已存在，则 \(X\) 的开子概形 \(U\) 在第一投影下的原像具有 \(U\times_TY\) 的泛性质。假设 \(X\) 有开覆盖 \(X_i\)，且各 \(X_i\times_TY\) 存在；在重叠处，这些空间的开子概形都表示同一个纤维积，故由泛性质唯一同构。唯一性还保证三重交集的余圈条件。黏合它们并黏合投影，所得空间对任意 \(Z\) 可在 \(Z\rightarrow X\) 的开原像上检验泛性质，因此是 \(X\times_TY\)。

从仿射环公式出发，先在 \(T,Y\) 仿射时对 \(X\) 使用这一黏合步骤，再交换 \(X,Y\)，处理 \(T\) 仿射的一般情形。最后把 \(T\) 作仿射开覆盖，同时限制 \(X,Y\) 后再次黏合。存在性成立；任意两个实现由泛性质给出互逆同构。

对平方映射，用仿射公式得到闭纤维环
\[
k[t]\otimes_{k[x]}k\cong k[t]/(t^2-a),
\]
其中 \(k[x]\to k\) 取 \(x\mapsto a\)。\(a=0\) 时商环为 \(k[t]/(t^2)\)，唯一素理想为 \((t)\)，长度为二。\(a\ne0\) 时，代数闭性给出 \(b\in k^*\) 使 \(b^2=a\)，而特征不为二使 \(b,-b\) 互异。中国剩余定理给出 \(k[t]/(t^2-a)\cong k\times k\)，两个局部因子各长一，总长度仍为二。原扩张以 \(1,t\) 为自由 \(k[x]\) 基，所以是秩二的有限平坦扩张；这里变化的是纤维的底点数与约化性，纤维代数的 \(k\) 维数始终为二。
\end{proof}

对概形态射 \(u:X\to Y\) 与 \(y\in Y\)，概形纤维的底空间可与集合 \(u^{-1}(y)\) 自然同胚，但它还带有该集合本身没有规定的结构层和剩余域。仿射地，若 \(Y=\operatorname{Spec}A\)、\(y\) 对应 \(\mathfrak p\)，而 \(U=\operatorname{Spec}B\subseteq X\) 映入 \(Y\)，则
\[
U_y=\operatorname{Spec}(B\otimes_A\kappa(\mathfrak p))
=\operatorname{Spec}(B_{A\setminus\mathfrak p}/\mathfrak pB_{A\setminus\mathfrak p}).
\]
局部化与商环的素理想对应，把这些点恰好识别为 \(B\) 中收缩到 \(\mathfrak p\) 的素理想；基本开集的对应还给出子空间拓扑。逐仿射图黏合，便得到上述同胚。平方映射的原点纤维只有一个底点，却保留平方零元素，差别发生在概形结构中。

\ProseParagraphBreak
一般的纤维积则不只由一对底层点决定。例如
\(\operatorname{Spec}\mathbb C\times_{\operatorname{Spec}\mathbb R}\operatorname{Spec}\mathbb C\)
的环同构为
\[
\mathbb C\otimes_{\mathbb R}\mathbb C\longrightarrow\mathbb C\times\mathbb C,
\qquad z\otimes w\longmapsto(zw,z\bar w).
\]
这由 \(\mathbb C[T]/(T^2+1)\cong\mathbb C\times\mathbb C\) 的两个根 \(i,-i\) 给出，故纤维积有两个点，而两个因子的底空间各只有一个点。两份乘法分别采用恒等与复共轭这两种 \(\mathbb R\) 嵌入；相同的底点对仍可产生不同的纤维积点。

\begin{definition}\label{b3:def:etale-scheme-morphism}
设 \(f:X\rightarrow Y\) 为概形态射。若对每个 \(x\in X\)，都存在仿射开邻域 \(U=\operatorname{Spec}B\subseteq X\)、\(V=\operatorname{Spec}A\subseteq Y\)，使 \(x\in U\)、\(f(U)\subseteq V\)，且所对应的环同态 \(A\rightarrow B\) 为有限展示，则称 \(f\) 为\term[jubu-youxian-zhanshi]{局部有限展示}[locally of finite presentation]的态射。若这些邻域还可取到使每个 \(A\rightarrow B\) 平展，即 \(B\) 为有限展示 \(A\) 代数、在 \(A\) 上平坦且 \(\Omega_{B/A}=0\)，则称 \(f\) 为平展态射。
\end{definition}

由\cref{b3:thm:etale-equivalent-characterizations}，定义中的环条件也可写成有限展示且形式平展，或在源的每一点附近具有\cref{b3:def:standard-smooth-model}中的标准平展模型。因此概形的平展性由这些仿射邻域中的有限组方程及其可逆 Jacobian 行列式描述。这里同时选择源的 \(U\) 和目标的 \(V\)，有限展示是在这两个邻域的坐标环之间要求的；茎同态本身的有限性不能替代这份邻域上的有限数据。

\begin{proposition}\label{b3:prop:etale-scheme-stability}
设 \(f:X\rightarrow Y\) 为平展概形态射。对任意概形态射 \(Y'\rightarrow Y\)，基变换 \(X\times_YY'\rightarrow Y'\) 平展。若 \(g:Y\rightarrow Z\) 也是平展概形态射，则复合 \(g\circ f:X\rightarrow Z\) 平展。
\end{proposition}
\begin{proof}
先核对环论中的有限展示、形式平展及局部化稳定性，再把所选的源、目标仿射图缩到相容的位置。后一项保证环论计算确实描述给定点附近的概形态射。

先说明所需的环运算。有限展示在基变换下保持，因为有限组生成元及关系的系数只需映到新底环。它也在复合下保持：若 \(B=A[T_1,\ldots,T_r]/(p_1,\ldots,p_s)\)、\(C=B[S_1,\ldots,S_v]/(q_1,\ldots,q_w)\)，把 \(q_j\) 的有限个 \(B\) 系数提升到 \(A[T_1,\ldots,T_r]\)，则 \(p_i\) 与这些提升后的 \(q_j\) 一起给出 \(C\) 在 \(A\) 上的有限展示。再由 \(B_b=B[T]/(bT-1)\)，目标上的单元素局部化也保持有限展示。\cref{b3:prop:formal-properties-stability}保证形式平展性在复合、基变换及目标局部化下保持；与\cref{b3:thm:etale-equivalent-characterizations}合用，便知平展环同态也保持于这些运算。

还需把仿射图取到相容的位置。两张含同一点的仿射开集 \(V_0,V_1\) 有共同的基本开邻域：先在 \(V_0\) 内取含该点的 \(D(a)\subseteq V_0\cap V_1\)，再在 \(V_1\) 内取含该点的 \(D(c)\subseteq D(a)\)。由\cref{b3:thm:affine-sheaf-sections}，\(c\) 在 \(D(a)\) 上可写为 \(d/a^N\)，故 \(D(c)\) 在 \(V_0\) 中就是 \(D(ad)\)。这样，改变这两张图只需比较单元素局部化。

对基变换，在 \(X\times_YY'\) 的一点取其 \(X\) 分量 \(x\)，并选定义中的平展图 \(U=\operatorname{Spec}B\rightarrow V=\operatorname{Spec}A\)。在该点的 \(Y'\) 分量附近取仿射开集 \(V'=\operatorname{Spec}C\)，使其映入 \(V\)。\cref{b3:prop:scheme-fiber-products-exist}给出包含所选点的仿射开邻域
\[
U\times_VV'=\operatorname{Spec}(B\otimes_AC).
\]
它到 \(V'\) 的环同态为 \(C\rightarrow B\otimes_AC\)，由前面的环计算平展。每一点都可这样处理，故基变换平展。

对复合，在 \(x\in X\) 及 \(f(x)\in Y\) 处分别取平展图 \(U_0\rightarrow V_0\) 和 \(V\rightarrow W\)，其中 \(W\subseteq Z\) 仿射。将 \(V_0,V\) 换成含 \(f(x)\) 的共同基本开邻域 \(V_1\)。由于 \(V_1\) 在 \(V_0\) 中基本开，\(U_1=U_0\cap f^{-1}(V_1)\) 是 \(U_0\) 的基本开集，而 \(U_1\rightarrow V_1\) 是原平展图的基变换。\(V_1\rightarrow V\) 由局部化给出，\(V\rightarrow W\) 平展；前面的环复合结论于是使 \(U_1\rightarrow W\) 平展。这是 \(x\) 处所需的图，故 \(g\circ f\) 平展。
\end{proof}

这些局部化比较也说明定义与仿射覆盖的选择无关。在源上取更小的基本开集、在目标上取基本开集并限制到其原像，仍得到平展环同态；另一组仿射图可在交集上作共同的基本开细化，从而使用同样的局部模型。

\ProseParagraphBreak

例如 \(\mathbb C=\mathbb R[T]/(T^2+1)\)，而 \(2T\) 在 \(\mathbb C\) 中可逆，故 \(\operatorname{Spec}\mathbb C\rightarrow\operatorname{Spec}\mathbb R\) 由标准平展模型给出。若一个概形态射将源同构地识别为目标的开子概形，就称它为\term[kai-jinru]{开浸入}[open immersion]。开浸入也平展：每个源点在目标的某张仿射图 \(\operatorname{Spec}A\) 中都有包含于这个开子概形的基本开邻域 \(D(a)\)，该处的环同态为 \(A\rightarrow A_a\)，它有限展示且形式平展。

拓扑空间的 Hausdorff 性可以用对角在普通乘积中闭来刻画。概形则要求相对纤维积中的对角成为闭嵌入，同时控制底空间与结构层。这里使用的相对纤维积也不同于点集的普通乘积，因此概形分离性不要求其 Zariski 底空间为 Hausdorff 空间。

\begin{definition}\label{b3:def:separated-finite-type-proper}\label{b3:def:closed-immersion}\label{b3:def:separated-morphism}\label{b3:def:finite-type-morphism}\label{b3:def:proper-morphism}
设 \(f:X\rightarrow Y\) 为概形态射。若它把 \(X\) 同胚地识别为 \(Y\) 的闭子空间，且层同态 \(\mathcal O_Y\rightarrow f_*\mathcal O_X\) 满射，则称它为闭嵌入。

若对角态射 \(\Delta_f:X\rightarrow X\times_YX\) 是闭嵌入，则称 \(f\) 具有\term[fenli-xing]{分离性}[separatedness]。

若对每个仿射开集 \(V=\operatorname{Spec}A\subseteq Y\)，其原像可被有限个仿射开集 \(\operatorname{Spec}B_i\) 覆盖，且各 \(B_i\) 为有限型 \(A\) 代数，则称 \(f\) 为有限型态射。

若 \(f\) 分离、有限型，且沿任意 \(Y'\rightarrow Y\) 基变换后的映射都把闭集送到闭集，则称 \(f\) 具有\term[guyou-xing]{固有性}[properness]，或称为固有态射。因此固有性同时要求有限型、相对对角的闭嵌入及普遍闭性。

若结构态射 \(X\to\operatorname{Spec}\mathbb Z\) 分离，则称概形 \(X\) 分离；讨论某个基概形 \(Y\) 上的相对分离性时，则是指所给态射 \(X\to Y\) 的分离性。
\end{definition}
\symidx{\(\Delta_f\)：概形态射 \(f\) 的对角态射}

闭嵌入还有如下局部描述：若 \(i:X\hookrightarrow Y\) 为闭嵌入，则对每个仿射开集 \(V=\operatorname{Spec}A\subseteq Y\)，有 \(i^{-1}(V)\cong\operatorname{Spec}(A/I)\)，且限制态射由商映射 \(A\to A/I\) 给出；反过来，这些局部商映射也确定闭嵌入。这一描述见\cite[Chapter II, Corollary 5.10, p. 116]{Hartshorne1977}。它还说明闭嵌入经任意基变换仍为闭嵌入：在仿射图上，任意 \(A\) 代数 \(B\) 都满足
\[
(A/I)\otimes_AB\cong B/IB,
\]
而这些局部商映射相容，能够黏合。

在仿射情形中，对角对应乘法满射 \(B\otimes_AB\rightarrow B\)，所以两个仿射概形之间的态射分离。对域 \(k\)，\(\operatorname{Spec}k[t]\) 因而分离，但其泛点不闭，底空间不是 Hausdorff。双原点直线却不分离：在来自两张不同仿射图的乘积 \(\operatorname{Spec}k[s,t]\) 中，对角的交是直线 \(s=t\) 去掉原点，缺少的原点对应两个不同原点组成的一对点。这一集合不闭。

\ProseParagraphBreak
有限环映射给出固有态射。它有限型，且源与目标均仿射，故分离；有限性经任意张量基变换保持。对有限环扩张中的闭集 \(V(J)\)，商环仍整于底环除以收缩理想，由\cref{b3:thm:lying-over}其像恰为 \(V(J\cap A)\)。任意概形基变换后，在基底的仿射开覆盖上应用此计算，便知所得映射仍闭。另一方面，对域 \(k\)，\(\mathbb A_k^1\rightarrow\operatorname{Spec}k\) 本身当然把闭集送到闭集，因为目标只有一个点；它也分离且有限型。但沿 \(\mathbb A_k^1\) 基变换后，闭集 \(xy=1\) 的投影为 \(D(x)\)，不再是闭集。因此这个态射不普遍闭，也就不固有，失败只有在改变基底后才显现。

\subsection{射影空间的固有性}
\label{b3:subsec:2404:projective-properness}

射影空间能补足仿射直线投影不闭的缺陷。仍看 \(xy=1\)，把第二个因子写成 \(\mathbb P_k^1\) 上的齐次坐标 \([Y:Z]\)，原仿射图为 \(Z\ne0\)、\(y=Y/Z\)。曲线的闭包为
\[
xY=Z\quad\text{于}\;\mathbb A_k^1\times_k\mathbb P_k^1.
\]
在 \(Y\ne0\) 的图上令 \(z=Z/Y\)，方程成为 \(z=x\)，所以整个闭包就是 \(x\mapsto(x,[1:x])\) 的图像。原来 \(x\ne0\) 的部分在其中稠密；当 \(x=0\) 时，新增的纤维是无穷远点 \([1:0]\)。射影化把原投影遗漏的底点补回，所得像为整条仿射直线。

\ProseParagraphBreak
证明固有性有三项任务：标准图直接给出有限型性，全部乘积图上的对角计算给出分离性，普遍闭性则需要追踪一个射影纤维为何为空。空纤维使每个齐次坐标幂零，有限个坐标将其合成高次齐次分量的共同消失；这些分量有限生成，再用 Nakayama 引理把消失扩大到基底邻域，便证明像的补集开。这里检查任意基点，不限于闭点。Hartshorne 的赋值判据要求态射有限型且源概形 Noether；该书随后在 Noether 概形之间证明射影态射固有\cite[Chapter II, Theorems 4.7 and 4.9, pp. 101--104]{Hartshorne1977}。固有性的定义见同书\cite[Chapter II, p. 100]{Hartshorne1977}。

\begin{theorem}[射影空间闭子概形的固有性]\label{b3:thm:closed-projective-proper}
设 \(Y\) 为概形，\(r\ge0\) 为整数，并令
\(\mathbb P_Y^r=\mathbb P_{\mathbb Z}^r\times_{\operatorname{Spec}\mathbb Z}Y\)。
若 \(X\hookrightarrow\mathbb P_Y^r\) 为闭嵌入，则复合态射 \(X\hookrightarrow\mathbb P_Y^r\rightarrow Y\) 固有。
\end{theorem}
\begin{proof}
先证明投影普遍闭。令 \(A\) 为任意交换环，\(Z\subseteq\mathbb P_A^r\) 为闭集，用齐次理想 \(J\subseteq A[X_0,\ldots,X_r]\) 表示为 \(V_+(J)\)，并记 \(C=A[X_0,\ldots,X_r]/J\)。标准图的局部化与取零次部分都与底环上的张量积相容，所以 \(\operatorname{Proj}(C\otimes_A\kappa(\mathfrak p))\) 是 \(\operatorname{Proj}C\) 在 \(\mathfrak p\) 处的概形纤维。其底空间与 \(Z\) 在该点的集合纤维相同，因此若 \(\mathfrak p\) 不在 \(Z\) 的像中，这个射影纤维便为空。每张 \(D_+(X_i)\) 图都为空，意味着在相应零次局部化环中 \(1=0\)，故对每个 \(i\) 都有整数 \(N_i\ge1\)，使 \(X_i^{N_i}=0\) 于纤维环。取
\[
N>\sum_{i=0}^r(N_i-1).
\]
任一次数 \(N\) 的单项式，至少有一个 \(X_i\) 的指数达到 \(N_i\)，否则总次数至多为右端。因此所有这些单项式在纤维环中为零，即 \(C_N\otimes_A\kappa(\mathfrak p)=0\)；更高次数的部分也为零。

\(C_N\) 由有限个次数 \(N\) 的单项式生成，故是有限生成 \(A\) 模；\(C_N\otimes_A\kappa(\mathfrak p)=0\) 由 Nakayama 引理推出 \((C_N)_{\mathfrak p}=0\)。取有限组生成元的共同零化分母，得 \(a\notin\mathfrak p\) 使 \((C_N)_a=0\)。环 \(C\) 由次数一元素生成，所以 \((C_n)_a=0\) 对所有 \(n\ge N\) 成立。于是 \((C_a)_+^N=0\)，每个齐次素理想都包含这个幂零的正次数理想，故 \(\operatorname{Proj}C_a=\varnothing\)。这正是 \(Z\) 在 \(D(a)\) 上的原像为空。像的补集因此开，证明投影闭；所用有限性来自有限个坐标，而不需要假设 \(A\) 是 Noether 环。

射影空间的标准仿射图及其过渡同态在任意张量基变换后仍给出同样的射影空间，所以上述计算对任何仿射基变换成立。对一般概形基变换，在基底的仿射开集上检验闭集像，便得普遍闭性。

有限型性由标准图立即得到。\(\mathbb P_A^r\) 只有 \(r+1\) 张标准仿射图，每张图为 \(r\) 元多项式环的谱，故到 \(\operatorname{Spec}A\) 的投影有限型。

再逐图证明分离性。记 \(U_i=D_+(X_i)\)。产品由全部 \(U_i\times_AU_j\) 覆盖，而对角在这一产品图上的原像是 \(U_i\cap U_j\)。在 \(U_i\) 上记 \(x_\nu=X_\nu/X_i\)，并约定 \(x_i=1\)；在第二张图 \(U_j\) 上记 \(y_\nu=Y_\nu/Y_j\)，并约定 \(y_j=1\)。因此产品图的坐标环为
\[
R_{ij}=A[x_\nu\ (\nu\ne i),y_\nu\ (\nu\ne j)].
\]
\(i=j\) 时，对角的环同态把 \(y_\nu\) 送到 \(x_\nu\)，其核由 \(y_\nu-x_\nu\) 生成，故是闭嵌入。\(i\ne j\) 时，两组坐标要表示同一方向，就必须有 \(x_jy_\nu=x_\nu\)。特别是 \(\nu=i\) 时 \(x_i=1\)，关系成为 \(x_jy_i=1\)，恰好要求两张图交叠处的 \(x_j\) 可逆。因此 \(U_i\cap U_j\) 的坐标环为 \(A[x_\nu\ (\nu\ne i)]_{x_j}\)，而对角对应
\[
\theta_{ij}:R_{ij}\longrightarrow A[x_\nu\ (\nu\ne i)]_{x_j},
\qquad x_\nu\longmapsto x_\nu,
\quad y_\nu\longmapsto\frac{x_\nu}{x_j}.
\]
其中 \(y_i\) 映为 \(x_j^{-1}\)，故 \(\theta_{ij}\) 满射。其核恰为
\[
L_{ij}=(x_jy_\nu-x_\nu:\nu\ne j).
\]
确实，\(\nu=i\) 的关系为 \(x_jy_i=1\)，使 \(x_j\) 在商环中可逆；其余关系给出 \(y_\nu=x_\nu/x_j\)。因此 \(R_{ij}/L_{ij}\) 与上述局部化环互相同构。这些环同态都来自同一个对角态射，在图的重叠上相容，所以全部局部闭嵌入黏合为对角的闭嵌入。投影因而分离。对任意基概形 \(Y\)，在其仿射开集上应用上述计算，便知 \(\mathbb P_Y^r\to Y\) 同样满足这三项条件。

最后令 \(X\) 为 \(\mathbb P_Y^r\) 的闭子概形。在基底的每张仿射图上，\(X\) 与有限张标准图相交，所得仿射环是相应多项式环的商，仍为有限型底环代数。其对角是射影空间对角沿 \(X\times_YX\to\mathbb P_Y^r\times_Y\mathbb P_Y^r\) 的基变换，故仍为闭嵌入。任意基变换后的 \(X\) 仍是相应射影空间的闭子概形，而闭嵌入与闭投影的复合保持闭集像为闭集。因此 \(X\to Y\) 有限型、分离且普遍闭，即为固有态射。
\end{proof}

% !TeX root = ../../Book3.tex
% 纲要标题：### 24.5 层上同调与 Čech 计算
% 纲要位置：计划纲要.md，第 2762 行。

\section{层上同调与 Čech 计算}
\label{b3:sec:2405}

令 \(k\) 为域，在射影直线 \(\mathbb P_k^1\) 的第一张标准图上取坐标 \(t=x_1/x_0\)。扭曲层 \(\mathcal O(-2)\) 在两张图上都自由，但在交集上，\(t^{-1}\) 不能写成 \(k[t]\) 中的多项式与 \(t^{-2}k[t^{-1}]\) 中元素之差。这种差别来自局部截面之间的相容条件。层上同调用复形记录这些黏合障碍，而具有仿射交集的开覆盖能把它们化为具体的核与余核计算。

% 纲要内容（仅作写作计划，尚非教材正文）：
% * 全局截面的非右正合性、内射层消解及上同调长正合列
% * flasque 层及其消解；Noether 仿射拟凝聚层消失的证明概要与所省略的层论步骤
% * 有限仿射覆盖的所有非空有限交仿射时，Čech 与层上同调的双复形比较；分离性是保证交集条件的充分条件
% * 射影直线扭曲层的零次与一次上同调；同一微分的核与余核对照
% * 双原点直线的交集仍仿射，比较可用；区分黏合、分离、固有及全局截面恢复的条件
% * 局部计算与全局性质的联系，以及进一步代数几何的具体参考位置
% ---

\subsection{全局截面与层上同调}
\label{b3:subsec:2404:02}

层公理解决相容截面的黏合，却不保证局部存在的原像能够选得相容。全局截面函子保持核，但可能把层的满射变成不满射；层上同调记录这种差别。

\begin{definition}\label{b3:def:sheaf-cohomology}
设 \(X\) 为拓扑空间，\(\mathcal F\) 为交换群层。取层范畴中的内射消解 \(0\rightarrow\mathcal F\rightarrow\mathcal I^\bullet\)，定义
\[
H^i(X,\mathcal F)=H^i\bigl(\Gamma(X,\mathcal I^\bullet)\bigr).
\]
这些群称为 \(\mathcal F\) 的\term[ceng-shangtongdiao]{层上同调}[sheaf cohomology]；特别地 \(H^0(X,\mathcal F)=\Gamma(X,\mathcal F)\)。
\end{definition}

上述层范畴有足够多内射对象。具体地，在每点 \(x\) 取茎的单射 \(\mathcal F_x\to E_x\)，其中 \(E_x\) 是内射交换群，记 \(i_x:\{x\}\to X\)。伴随公式
\[
\operatorname{Hom}_X(\mathcal G,i_{x*}E_x)
\cong\operatorname{Hom}_{\mathbb Z}(\mathcal G_x,E_x)
\]
及取茎的正合性使 \(i_{x*}E_x\) 内射。它们的乘积也内射，因为向乘积映射等价于逐项映射，逐项延拓可同时完成。自然映射 \(\mathcal F\to\prod_xi_{x*}E_x\) 为单射：一个非零芽在对应的 \(x\) 分量中仍非零。逐次对余核重复，就得到内射消解；比较映射及其同伦按内射延拓逐次构造，故定义不依赖消解。与环模层的比较见\cite[Chapter III, Proposition 2.2, Corollary 2.3 and Proposition 2.6, pp. 207--208]{Hartshorne1977}。

\ProseParagraphBreak
短正合列 \(0\rightarrow\mathcal F'\rightarrow\mathcal F\rightarrow\mathcal F''\rightarrow0\) 诱导自然的长正合列
\[
0\rightarrow H^0(X,\mathcal F')\rightarrow H^0(X,\mathcal F)
\rightarrow H^0(X,\mathcal F'')\xrightarrow{\delta}
H^1(X,\mathcal F')\rightarrow H^1(X,\mathcal F)\rightarrow\cdots.
\]
构造与\secref{b3:sec:2201}中的导出函子相同：先把两端嵌入内射对象，再取直和为中间项，逐次对余核重复，得到逐项分裂的内射消解短正合列。取全局截面仍逐项分裂，随后用复形的连接同态即得上述序列。尤其 \(\delta\) 把局部可提升、全局未必可提升的截面送到它的障碍类。

\ProseParagraphBreak
若一个交换群层从较大开集到较小开集的限制总是满射，就称它为 flasque 层。内射层是 flasque 层，flasque 层没有正次上同调，因而任意 flasque 消解都可用于计算层上同调；这些层论性质见\cite[Chapter III, Lemma 2.4, Proposition 2.5 and Remark 2.5.1, pp. 207--208]{Hartshorne1977}。

\ProseParagraphBreak
下面的仿射消失定理采用 Hartshorne 的 Noether 仿射版本\cite[Chapter III, Theorem 3.5, p. 215]{Hartshorne1977}。概要说明如何把模消解变成 flasque 层消解；所省略的是一般层论中的截面延拓论证，以及支集上的截面与挠子模之间的识别。

\begin{theorem}[仿射拟凝聚层的上同调消失]\label{b3:thm:affine-quasicoherent-vanishing}
设 \(A\) 为 Noether 交换环，\(M\) 为 \(A\) 模，\(X=\operatorname{Spec}A\)。则对每个整数 \(i>0\)，有
\[
H^i(X,\widetilde M)=0.
\]
\end{theorem}
\begin{proofsketch}
flasque 层的消失可由内射消解说明：内射层是 flasque 层，因为在开集上的截面等价于从该开集上常值整数层的零延拓到它的映射，开集包含给出零延拓层的单射，内射性使映射延拓。若短正合列的核为 flasque 层，商层的局部截面可在全开集提升：将两个局部提升在交集上的差延拓到一个开集，扣去后便可黏合，再对开覆盖归纳。用 Zorn 引理处理无限覆盖，可得到全局提升。因而 flasque 层嵌入内射层后的余核仍 flasque，逐次取余核使全局截面消解正合，证明消失。

对于内射 \(A\) 模 \(E\)，\(\widetilde E\) 是 flasque 层。这一结论见\cite[Chapter III, Proposition 3.4, pp. 214--215]{Hartshorne1977}。对闭集 \(Y=\overline{\operatorname{Supp}\widetilde E}\) 作 Noether 归纳。给开集 \(U\) 上的截面 \(s\)，选 \(D(f)\subseteq U\) 且与 \(Y\) 相交。由\cref{b3:lem:injective-cech-acyclic}，\(E\to E_f\) 满射，所以 \(s|_{D(f)}\) 来自某个全局截面 \(t\)。差 \(s-t|_U\) 支集包含在 \(V(f)\)，正是 \(\widetilde{\Gamma_{(f)}E}\) 在 \(U\) 上的截面；这一识别用有限基本开覆盖把各处分母及零化幂次取成共同上界。该挠子模由同一引理内射，其支集闭包包含在严格较小的 \(Y\cap V(f)\) 中，归纳使差也能全局延拓。于是 \(s\) 能全局延拓，证明限制满射；\(U\cap Y=\varnothing\) 时截面为零，无须选择 \(f\)。

现在取 \(M\) 的内射模消解。伴随层函子由局部化正合，得到 \(\widetilde M\) 的 flasque 消解。取全局截面后回到原模消解，因而正次上同调为零。
\end{proofsketch}

\subsection{仿射交集与 Čech 比较}
\label{b3:subsec:2405:cech-comparison}

开覆盖给出的是一个具体的截面复形；要使它计算层上同调，需要各重叠开集不再留下正次上同调。仿射交集上的拟凝聚层消失正好提供这个条件。

\ProseParagraphBreak
对开覆盖 \(\mathfrak U=(U_i)\)，把各重叠上的截面组成复形
\[
\check C^p(\mathfrak U,\mathcal F)
=\prod_{i_0<\cdots<i_p}\mathcal F(U_{i_0}\cap\cdots\cap U_{i_p}),
\qquad
(\delta s)_{i_0\cdots i_{p+1}}
=\sum_{\nu=0}^{p+1}(-1)^\nu
s_{i_0\cdots\widehat{i_\nu}\cdots i_{p+1}}|.
\]
末尾的限制取到共同交集。删除两指标的两种次序符号相反，故 \(\delta^2=0\)。其零次核就是全局截面，但高次同调不对任意覆盖自动等于刚定义的 \(H^i\)。

\ProseParagraphBreak
Hartshorne 的分离情形见\cite[Chapter III, Theorem 4.5, p. 222]{Hartshorne1977}。下面只要求给定覆盖的有限交仿射；分离性是保证这一条件成立的一种充分条件。

\begin{theorem}[仿射开覆盖的 Čech 比较]\label{b3:thm:affine-cech-comparison}
设 \(X\) 为 Noether 概形，\(\mathcal F\) 为 \(X\) 上的拟凝聚层，\(\mathfrak U=(U_i)_{i=1}^m\) 为有限仿射开覆盖，并假设每个非空有限交集 \(U_{i_0}\cap\cdots\cap U_{i_p}\) 都仿射。则对每个整数 \(q\ge0\)，有自然同构
\[
\check H^q(\mathfrak U,\mathcal F)\cong H^q(X,\mathcal F).
\]
特别地，若 \(X\) 还分离，则任意有限仿射开覆盖都满足上述交集条件。
\end{theorem}
\begin{proof}
先说明 flasque 性如何使 Čech 方向正合。对任意交换群层 \(\mathcal G\)，记 \(U_{i_0\cdots i_p}=U_{i_0}\cap\cdots\cap U_{i_p}\)，并令 \(j_{i_0\cdots i_p}:U_{i_0\cdots i_p}\hookrightarrow X\) 为开嵌入。将 Čech 复形层化，得到
\[
\mathscr C^p(\mathfrak U,\mathcal G)
=\prod_{i_0<\cdots<i_p}
 j_{i_0\cdots i_p*}(\mathcal G|_{U_{i_0\cdots i_p}}).
\]
增广层复形 \(0\to\mathcal G\to\mathscr C^0\to\mathscr C^1\to\cdots\) 正合。在一个点 \(x\) 的茎上，选含 \(x\) 的 \(U_a\)，并把邻域缩小到 \(U_a\) 内。对交替余链在指标前插入 \(a\) 得到次数减一映射 \(h\)，重复指标的分量取零；交替符号直接给出 \(\delta h+h\delta=1\)。这证明茎上的正合性，故也证明层复形正合。

若 \(\mathcal G\) flasque，则它在任意开集上的限制仍 flasque。各个 \(\mathscr C^p\) 也 flasque：开嵌入的直像在开集 \(V\) 上取 \(\mathcal G(V\cap U_{i_0\cdots i_p})\)，限制满射性逐个因子保持，有限乘积亦保持。因此这个层复形是 flasque 消解，取全局截面后计算 \(H^p(X,\mathcal G)\)。而 \(\mathcal G\) 没有正次上同调，所以增广截面复形
\[
0\to\Gamma(X,\mathcal G)\to
\check C^0(\mathfrak U,\mathcal G)\to
\check C^1(\mathfrak U,\mathcal G)\to\cdots
\]
正合。这一层化论证也见\cite[Chapter III, Lemma 4.2 and Proposition 4.3, pp. 220--221]{Hartshorne1977}。

取交换群层范畴中的内射消解 \(\mathcal F\to\mathcal I^\bullet\)，组成第一象限双复形
\[
D^{p,q}=\check C^p(\mathfrak U,\mathcal I^q),
\qquad d_{\mathrm{Tot}}=\delta+(-1)^p d_{\mathcal I}.
\]
各 \(\mathcal I^q\) 是 flasque 层，刚证的增广正合性使每一行的上同调只在 \(p=0\) 剩下 \(\Gamma(X,\mathcal I^q)\)。

对另一方向，限制层消解到任意 \(U_{i_0\cdots i_p}\) 后仍正合，因为可逐茎检验；限制后的各 \(\mathcal I^q\) 仍 flasque。因此
\[
\Gamma(U_{i_0\cdots i_p},\mathcal I^\bullet)
\]
通过这个 flasque 消解计算 \(H^q(U_{i_0\cdots i_p},\mathcal F|_{U_{i_0\cdots i_p}})\)。非空交集仿射且 Noether，\(\mathcal F\) 的限制拟凝聚，故由\cref{b3:thm:affine-quasicoherent-vanishing}，这些群在 \(q>0\) 时为零。空交集的截面复形为零。于是每一列的上同调只在 \(q=0\) 剩下相应交集上的 \(\mathcal F\) 截面。

这两种增广给出自然的复形映射
\[
\begin{tikzcd}[column sep=large]
\Gamma(X,\mathcal I^\bullet)
 \arrow[r,"\text{限制}"]
&\operatorname{Tot}D
&\check C^\bullet(\mathfrak U,\mathcal F)
 \arrow[l,"\text{消解增广}"']
\end{tikzcd}
\]
它们都诱导上同调同构。计算与\cref{b3:thm:cech-local-cohomology}的双复形论证相同：在固定总次数中，从最外侧分量开始，利用相应行或列的正合性把它写成微分的像，扣去这一边界后依次消去其余分量，最后只留下边上的复形；边界的识别使用同一个过程。这里只有 \(0\le p<m\)，每个总次数中均只有有限个分量，所以过程有限。两边分别计算层上同调与 Čech 上同调，由此得到所述同构。消解比较映射及其同伦都传到这个双复形，故同构不依赖消解选择，且关于 \(\mathcal F\) 自然。

最后验证分离性给出的充分条件。若 \(U,V\subseteq X\) 仿射，则 \(U\cap V\) 由分离对角在 \(U\times_{\operatorname{Spec}\mathbb Z}V\) 中给出闭子概形，后者仿射，故交集仿射。对交集个数归纳，所有非空有限交集均仿射。
\end{proof}

这一证明把真正需要的条件留在交集上。例如双原点直线的两张仿射图交于 \(\operatorname{Spec}k[t,t^{-1}]\)，所以这一覆盖也能计算拟凝聚层的层上同调，尽管整个概形不分离。

\begin{proposition}\label{b3:prop:projective-line-cohomology}
设 \(k\) 为域，\(n\in\mathbb Z\)。在 \(X=\mathbb P_k^1\) 上，
\[
\dim_kH^0(X,\mathcal O_X(n))=\max\{n+1,0\},\qquad
\dim_kH^1(X,\mathcal O_X(n))=\max\{-n-1,0\},
\]
且 \(H^i(X,\mathcal O_X(n))=0\) 对 \(i\ge2\) 成立。
\end{proposition}
\begin{proof}
用\cref{b3:tab:projective-line-twist-charts}中的两张标准图及生成元 \(e_1=t^ne_0\)，Čech 复形只有两项，其微分可采用\eqref{b3:eq:projective-line-twist-cech-map}中的 \(\delta_n\)。改变整个微分的符号不影响核与余核。核已在\secref{b3:sec:2403}算出：\(n\ge0\) 时由 \(1,t,\ldots,t^n\) 参数化，\(n<0\) 时为零。余核为
\[
\frac{k[t,t^{-1}]}{k[t]+t^nk[t^{-1}]}.
\]
分母包含指数 \(j\ge0\) 以及 \(j\le n\) 的全部单项式。留下的基恰为 \(t^j\)、\(n<j<0\)，故维数是 \(\max\{-n-1,0\}\)。没有更高的 Čech 项；射影直线是 Noether 概形，两张标准图及其交集均仿射，故\cref{b3:thm:affine-cech-comparison}说明所得即为层上同调。
\end{proof}

例如 \(H^1(\mathbb P_k^1,\mathcal O(-2))\) 由 \(t^{-1}\) 表示。它说明交集上存在不能写成两张图中截面之差的元素，这是一项真正的黏合障碍，而非仅仅引入了新的函子符号。

\ProseParagraphBreak
同一个微分 \(\delta_n\) 的核与余核给出两种互补的信息：前者是可以黏合的局部截面，后者是交集上不能由局部截面之差消去的类。\cref{b3:tab:projective-line-cohomology}把刚才的计算并列写出，表中系数层均为 \(\mathbb P_k^1\) 上的 \(\mathcal O(n)\)。

\begin{table}[htbp]
\centering
\caption{射影直线扭曲层的上同调}
\label{b3:tab:projective-line-cohomology}
\begin{tabularx}{\linewidth}{@{}lXc@{}}
\toprule
上同调 & 复形描述与代表元 & \(k\) 维数 \\
\midrule
\(H^0\) & \(\ker\delta_n\)；\(n\ge0\) 时以 \(1,t,\ldots,t^n\) 表示基，\(n<0\) 时为零 & \(\max\{n+1,0\}\) \\
\(H^1\) & \(\operatorname{coker}\delta_n\)；基由余类 \(t^j\)、\(n<j<0\) 给出 & \(\max\{-n-1,0\}\) \\
\(H^i\)、\(i\ge2\) & 两张开集的 Čech 复形没有更高次项 & \(0\) \\
\bottomrule
\end{tabularx}
\end{table}

\subsection{局部计算与全局性质}
\label{b3:subsec:2404:03}

至此，环、模、分次环与几何对象之间已有三种可实际使用的对应：仿射概形由环及其局部化决定，仿射拟凝聚层由模决定，射影概形则把分次数据放在次数零的局部图中。第三种对应会忘掉某些分次模信息，因此不能把仿射模层等价直接搬到 Proj。

\ProseParagraphBreak
本卷的代数方法也由此获得几何解释。模在素点的局部化变成层的茎，Fitting 理想描述纤维维数发生变化的闭条件；正则局部环描述点的局部结构，光滑性还要比较基域或基概形。局部上同调研究指定闭支集附近的截面障碍，与本节不带支集的层上同调有关，但不能仅因二者使用 Čech 复形便把它们的次数和对象混同。具体联系需要同时写明开集、闭支集和所用增广复形。

\ProseParagraphBreak
Hartshorne 的《Algebraic Geometry》将这些构造作了系统展开：Chapter II, §§1--3 建立层、概形及纤维积；§4 研究分离与固有态射；§5 系统整理拟凝聚层和扭曲层；§7 把线性系与射影态射联系起来\cite[Chapter II, \S\S 1--5 and \S 7]{Hartshorne1977}。本章已经展开基本构造，但射影嵌入的系统判据、凝聚上同调的有限性以及一般基变换定理仍须另行学习。


\begin{chaptersummary}
结构层把每个邻域中的可用分式保留下来，因此能够区别同一素谱上的不同环结构。仿射概形的态射由反向的环同态决定，拟凝聚层由模决定。在仿射概形上，有限生成投射模对应局部有限秩的局部自由层。离开仿射情形后，过渡同构仍是对象的一部分，全局截面一般不能恢复全部局部信息。

\ProseParagraphBreak
Proj 通过零次局部化把齐次数据化为仿射图。扭曲层记录不同图中生成元的次数变化；射影直线上的 Čech 计算展示了相容方程的核与余核各自的意义。纤维积把局部对象的交叠化为新的概形，平展态射在复合与换基下保持，因而容许把可分扩张也作为局部观察。分离性、有限型和固有性进一步限制态射的全局行为，但这些条件与正则、正规或光滑等局部性质不能相互替代。

\ProseParagraphBreak
章首的双原点直线可以由仿射图黏合成概形，却不分离。它的两张图及交集仍然仿射，因此同一组 Čech 截面差仍可计算层上同调。分离性保证仿射开集的有限交仿射；对一组给定覆盖而言，应核对的是这些交集，而不能只凭空间不分离就否定比较。
\end{chaptersummary}

\BookExercises
未另说明时，\(k\) 为域，射影多项式环取标准分次。

\begin{optionalexercise}\label{b3:ex:24-local-gluing-crt}
在 \(X=\operatorname{Spec}\mathbb Z\) 上取 \(M=\mathbb Z/6\mathbb Z\)。说明 \(D(2),D(3)\) 覆盖 \(X\)，并求使局部截面分别为 \(1\in M_2\cong\mathbb Z/3\mathbb Z\)、\(0\in M_3\cong\mathbb Z/2\mathbb Z\) 的全局截面。
\end{optionalexercise}
\begin{solution}
没有素理想同时包含 \(2,3\)，故两开集覆盖全谱。在交集 \(D(6)\) 上 \(M_6=0\)，两截面自动相容。所求整数类满足 \(m\equiv1\pmod3\)、\(m\equiv0\pmod2\)，唯一解为 \(4\pmod6\)。这里的黏合正是模的局部化版本的中国剩余计算。
\end{solution}

\begin{optionalexercise}\label{b3:ex:24-infinitesimal-automorphisms}
设 \(k\) 为域。求 \(k\) 上概形 \(\operatorname{Spec}k[\varepsilon]/(\varepsilon^2)\) 的全部自同构，并说明底拓扑空间为什么不能区分它们。
\end{optionalexercise}
\begin{solution}
态射反向对应 \(k\) 代数同态。幂零元必须送到幂零元，故 \(\varepsilon\mapsto c\varepsilon\)、\(c\in k\)；同态可逆当且仅当 \(c\ne0\)，逆由 \(c^{-1}\) 给出。因此自同构群为 \(k^\times\)。底空间只有一点，所以这些自同构都诱导同一个点映射。
\end{solution}

\begin{optionalexercise}\label{b3:ex:24-arithmetic-two-points}
描述 \(\operatorname{Spec}(\mathbb Z/12\mathbb Z)\) 的点、开集、各点的局部环及结构层。
\end{optionalexercise}
\begin{solution}
中国剩余同构为 \(\mathbb Z/12\mathbb Z\cong\mathbb Z/4\mathbb Z\times\mathbb Z/3\mathbb Z\)。两点分别来自素数 \(2,3\)，各自开闭，底空间为离散两点。点 \(2\) 的局部环是 \(\mathbb Z/4\mathbb Z\)，点 \(3\) 的局部环是 \(\mathbb Z/3\mathbb Z\)。整个空间上的截面为两环的乘积，单点上的截面为相应因子，空集上为零环；限制是投影。点 \(2\) 上的非零幂零元说明它不是剩余域 \(\mathbb F_2\) 所给的约化点。
\end{solution}

\begin{optionalexercise}\label{b3:ex:24-rational-module-fibers}
设 \(p\) 为素数，\(X=\operatorname{Spec}\mathbb Z\)、\(M=\mathbb Z[1/p]/\mathbb Z\)，令 \(\mathcal F=\widetilde M\)。计算各点的茎及纤维，证明所有纤维都为零而 \(\mathcal F\ne0\)。说明为什么这不可能发生于有限生成整数模的伴随层。
\end{optionalexercise}
\begin{solution}
每个 \(M\) 元素都被某个 \(p\) 的幂零化，所以零素理想处的茎为零。若 \(q\ne p\) 为素数，\(p\) 在 \(\mathbb Z_{(q)}\) 中可逆，茎也为零。在 \((p)\) 处，所有不被 \(p\) 整除的整数在每个有限 \(p\) 幂阶循环子群上都可逆，因此局部化不改变 \(M\)，该茎为非零的 \(M\)。乘 \(p\) 在 \(M\) 上满射：\(a/p^n+\mathbb Z\) 是 \(a/p^{n+1}+\mathbb Z\) 的 \(p\) 倍。因此闭点 \((p)\) 的纤维为 \(M/pM=0\)，其余点的纤维随零茎也为零。

有限组元素被同一个 \(p^N\) 零化，但 \(M\) 包含不被该幂零化的元素，所以它不有限生成。若原模有限生成，则各个茎也是有限生成局部模；零纤维由 Nakayama 引理迫使每个茎为零，继而层为零。本例正缺少这项有限性。
\end{solution}

\begin{optionalexercise}\label{b3:ex:24-ideal-sheaf-fiber}
设 \(k\) 为域，\(A=k[x,y]\)、\(X=\operatorname{Spec}A\)。考察局部自由层映射 \(\phi:\mathcal O_X^2\to\mathcal O_X^2\)，矩阵为 \(\operatorname{diag}(x,y)\)。证明它为层单射，求余核，并对所有素点求其纤维映射的秩。
\end{optionalexercise}
\begin{solution}
\(A\) 是整环，故 \((a,b)\mapsto(xa,yb)\) 单射；伴随层函子正合，所以它也给出层单射，余核为 \(\widetilde{A/(x)\oplus A/(y)}\)。在素点 \(\mathfrak p\) 的剩余域中，矩阵仍为对角矩阵，秩等于 \(x,y\) 中不属于 \(\mathfrak p\) 的元素数。因此 \(D(xy)\) 上秩为二，\(V(x)\cup V(y)\) 去掉原点后秩为一，原点 \((x,y)\) 处秩为零。局部化的映射仍单射，纤维映射却会失去单射性；这两个操作分别由正合局部化与通常仅右正合的剩余域张量完成。
\end{solution}

\begin{optionalexercise}\label{b3:ex:24-dual-number-tangent-map}
设 \(k\) 为域，\(r\ge0\) 为整数。求所有将闭点送到原点的 \(k\) 概形态射
\(\operatorname{Spec}k[\varepsilon]/(\varepsilon^2)\rightarrow\mathbb A_k^r\)。
\end{optionalexercise}
\begin{solution}
环同态由每个 \(x_i\) 的像决定。闭点映到原点要求常数项为零，故唯一写成 \(x_i\mapsto a_i\varepsilon\)，其中 \(a_i\in k\) 任意。全部态射与 \(k^r\) 对应。这些向量不是新的普通点，而是以原点为中心的一阶方向。
\end{solution}

\begin{optionalexercise}\label{b3:ex:24-split-vector-bundle}
设 \(k\) 为域，\(a,b\in\mathbb Z\)。在 \(\mathbb P_k^1\) 的标准覆盖上以矩阵 \(\operatorname{diag}(t^a,t^b)\) 黏合秩二局部自由层 \(\mathcal E\)。求其零次、一次上同调维数及最高外幂，并证明 \(a=b=1\) 时它不平凡。
\end{optionalexercise}
\begin{solution}
过渡矩阵分为两块，故 \(\mathcal E=\mathcal O(a)\oplus\mathcal O(b)\)。逐分量的 Čech 复形给
\[
h^0(\mathcal E)=\max\{a+1,0\}+\max\{b+1,0\},\quad
h^1(\mathcal E)=\max\{-a-1,0\}+\max\{-b-1,0\}.
\]
局部最高外积的过渡系数为矩阵行列式 \(t^{a+b}\)，故 \(\bigwedge^2\mathcal E=\mathcal O(a+b)\)。当 \(a=b=1\) 时全局截面维数为四，而平凡秩二层为二，所以不可能同构。
\end{solution}

\begin{optionalexercise}\label{b3:ex:24-double-origin-cech}
设 \(k\) 为域，\(X\) 为将两份 \(\operatorname{Spec}k[t]\) 沿 \(D(t)=\operatorname{Spec}k[t,t^{-1}]\) 按同一坐标黏合得到的双原点直线。计算 \(H^0(X,\mathcal O_X)\) 与 \(H^1(X,\mathcal O_X)\)，核对仿射开覆盖的 Čech 比较条件，并证明 \(X\) 不仿射。
\end{optionalexercise}
\begin{solution}
两张图及其交集均为 Noether 仿射概形，结构层拟凝聚，故\cref{b3:thm:affine-cech-comparison}适用，不需要再假设 \(X\) 分离。两条限制均为通常包含，所以 Čech 微分为
\[
k[t]\oplus k[t]\longrightarrow k[t,t^{-1}],
\qquad(a,b)\longmapsto a-b.
\]
其核为对角的 \(k[t]\)，因此 \(H^0(X,\mathcal O_X)=k[t]\)；其余核为 \(H^1(X,\mathcal O_X)=k[t,t^{-1}]/k[t]\)，以 \(t^{-1},t^{-2},\ldots\) 为基。\(X\) 由两张 Noether 仿射图覆盖，因而 Noether。若它仿射，\cref{b3:thm:affine-quasicoherent-vanishing}将使 \(H^1(X,\mathcal O_X)=0\)，与上述计算矛盾。
\end{solution}

\begin{optionalexercise}\label{b3:ex:24-weighted-twist}
设 \(k\) 为域，\(d\ge2\) 为整数，\(S=k[x]\)、\(\deg x=d\)，\(X=\operatorname{Proj}S\)。对所有整数 \(n,m\) 求 \(\mathcal O_X(n)\)，并判定自然乘法 \(\mathcal O_X(n)\otimes\mathcal O_X(m)\to\mathcal O_X(n+m)\) 何时为同构。
\end{optionalexercise}
\begin{solution}
\(D_+(x)\) 为全空间，\((S_x)_0=k\)，故 \(X\cong\operatorname{Spec}k\)。\(S_x=k[x,x^{-1}]\) 的次数 \(n\) 部分，在 \(d\mid n\) 时为 \(kx^{n/d}\)，否则为零。因此 \(\mathcal O_X(n)\cong\mathcal O_X\) 当且仅当 \(d\mid n\)，其余情形为零层。

若 \(d\mid n\)、\(d\mid m\)，乘法将两个基向量送到 \(x^{(n+m)/d}\)，为同构。若 \(d\nmid(n+m)\)，源与目标都为零，也为同构。余下的情形是 \(d\mid(n+m)\) 而两个加数都不被 \(d\) 整除，此时源为零、目标为 \(\mathcal O_X\)，不可能同构。特别取 \(n=1\)、\(m=d-1\)，就得到扭曲乘法法则失效的实例。只知道底概形是一点，不能删去标准分次的条件。
\end{solution}

\begin{optionalexercise}\label{b3:ex:24-projective-nilpotent-thickening}
设 \(k\) 为域，令 \(S=k[x_0,x_1,z]/(z^2)\)，三个变量次数均为一。求 \(\operatorname{Proj}S\) 的两张标准仿射图，并描述其约化概形。
\end{optionalexercise}
\begin{solution}
每个素理想包含 \(z\)，故 \(D_+(x_0),D_+(x_1)\) 覆盖 Proj。在第一张图中环为 \(k[t,w]/(w^2)\)，其中 \(t=x_1/x_0,w=z/x_0\)；第二张图为 \(k[u,v]/(v^2)\)，其中 \(u=x_0/x_1,v=z/x_1\)。交集变换 \(u=t^{-1},v=w/t\)。各图商去幂零元后按 \(u=t^{-1}\) 黏合为 \(\mathbb P_k^1\)，而原概形的局部环保留平方零方向。
\end{solution}

\begin{optionalexercise}\label{b3:ex:24-conic-is-projective-line}
设 \(k\) 为域。用仿射图证明 \(C=\operatorname{Proj}k[x,y,z]/(xz-y^2)\cong\mathbb P_k^1\)，并求此同构下 \(\mathcal O_C(1)\) 对应的扭曲层。
\end{optionalexercise}
\begin{solution}
若齐次素理想包含 \(x,z\)，也包含 \(y\)，故 \(D_+(x),D_+(z)\) 覆盖。在 \(D_+(x)\) 上令 \(t=y/x\)，有 \(z/x=t^2\)，坐标环为 \(k[t]\)；在 \(D_+(z)\) 上令 \(u=y/z\)，有 \(x/z=u^2\)，环为 \(k[u]\)。交集 \(u=t^{-1}\)，所以两图黏合为射影直线。在两图中 \(\mathcal O_C(1)\) 的基分别为 \(x,z\)，关系 \(z=t^2x\)，因此它对应 \(\mathcal O_{\mathbb P^1}(2)\)。这说明同一概形的射影嵌入还记录了额外的可逆层。
\end{solution}

\begin{optionalexercise}\label{b3:ex:24-fat-projective-point}
设 \(k\) 为域，三个变量取次数一，\(X=\operatorname{Proj}k[x,y,z]/(x^3,y^2)\)。求全局函数、底点数、长度及唯一点处的切空间维数。判断它是否同构于 \(\operatorname{Spec}k[t]/(t^6)\)。
\end{optionalexercise}
\begin{solution}
所有齐次素理想都包含 \(x,y\)，避开无关理想要求不含 \(z\)，所以 \(D_+(z)=X\)。取 \(a=x/z\)、\(b=y/z\)，得到
\[
X\cong\operatorname{Spec}k[a,b]/(a^3,b^2).
\]
因此全局函数环以 \(1,a,a^2,b,ab,a^2b\) 为 \(k\) 基，底空间只有一点，长度为六。极大理想模平方以 \(a,b\) 的类为基，切空间维数为二；\(k[t]/(t^6)\) 的对应维数为一。两者虽有相同点数与长度，局部环的嵌入维数不同，所以不同构。
\end{solution}

\begin{optionalexercise}\label{b3:ex:24-proj-forgets-torsion}
设 \(k\) 为域，在标准分次环 \(S=k[x,y]\) 中取 \(I=x(x,y)^2\)、\(J=(x)\)。计算 \(I^{\mathrm{sat}}\) 及 \(J/I\) 的分次 Hilbert 级数，求零化整个 \(J/I\) 所需的最小无关理想幂，并比较两者给出的射影闭子概形。
\end{optionalexercise}
\begin{solution}
令 \(\mathfrak m=(x,y)\)。因 \(\mathfrak m^2x\subseteq I\)，有 \(J\subseteq I^{\mathrm{sat}}\)。反之，若 \(f\in I^{\mathrm{sat}}\)，则某个 \(y^Nf\) 属于 \(I\subseteq(x)\)；素元 \(x\) 不整除 \(y\)，故 \(x\mid f\)，所以 \(I^{\mathrm{sat}}=J\)。乘以 \(x\) 给出 \(J/I\cong(S/\mathfrak m^2)(-1)\)，其基为 \(x,x^2,xy\) 的类，Hilbert 级数为 \(t+2t^2\)。\(\mathfrak m^2\) 零化它，\(\mathfrak m\) 却不零化 \(x\) 的类，故最小幂为二。这个差模在所有标准图上为零，两份射影闭子概形相同，尽管它们的齐次理想不同。
\end{solution}

\begin{optionalexercise}\label{b3:ex:24-ramified-fibers}
设 \(k\) 为代数闭域，\(n\ge2\) 为整数，考察 \(k[x]\to k[t]\)、\(x\mapsto t^n\)。计算各闭纤维的点数及局部长度，并说明非零闭纤维何时约化。特征 \(p>0\) 时写 \(n=p^e m\)、\(p\nmid m\)；特征零时置 \(e=0\)、\(m=n\)。
\end{optionalexercise}
\begin{solution}
\(x=a\) 的纤维环为 \(k[t]/(t^n-a)\)。\(a=0\) 时只有一点，局部长度为 \(n\)。非零 \(a\) 时，在正特征置 \(q=p^e\)，特征零置 \(q=1\)。唯一取 \(c\in k\) 使 \(c^q=a\)，有 \(t^n-a=(t^m-c)^q\)。\(t^m-c\) 有 \(m\) 个互异根，因为它的根非零且导数 \(mt^{m-1}\) 不在根处消失。由中国剩余定理，纤维环为 \(m\) 份 \(k[\varepsilon]/(\varepsilon^q)\) 的直积，故有 \(m\) 个点，各点局部长度为 \(q\)，总长度为 \(mq=n\)。非零纤维约化恰好要求 \(q=1\)，即正特征不整除 \(n\)。原代数始终以 \(1,t,\ldots,t^{n-1}\) 为自由基，有限平坦性对应固定的总长度，而不要求固定点数或约化性。
\end{solution}

\begin{optionalexercise}\label{b3:ex:24-finite-proper-not-reduced}
设 \(k\) 为域。证明 \(\operatorname{Spec}k[t]/(t^3)\rightarrow\operatorname{Spec}k\) 固有，但源概形不约化。对任意 \(k\) 代数 \(A\) 写出其基变换。
\end{optionalexercise}
\begin{solution}
源环作为 \(k\) 模有限生成，故态射有限，依正文的有限态射论证为固有。\(t\) 的类非零且立方为零，源不约化。基变换为 \(\operatorname{Spec}A[t]/(t^3)\rightarrow\operatorname{Spec}A\)，仍有限自由，基为 \(1,t,t^2\)。固有性约束态射的全局闭性，不会排除幂零结构。
\end{solution}

\begin{optionalexercise}\label{b3:ex:24-localization-not-proper}
判断 \(\operatorname{Spec}\mathbb Z[1/2]\rightarrow\operatorname{Spec}\mathbb Z\) 是否分离、有限型、固有。
\end{optionalexercise}
\begin{solution}
它是仿射态射，故分离；\(\mathbb Z[1/2]\cong\mathbb Z[u]/(2u-1)\)，故有限型。像为 \(D(2)\)，包含一般点却不含闭点 \((2)\)。一般点的闭包是全谱，所以 \(D(2)\) 不闭；映射连本身都不是闭映射，因而不固有。
\end{solution}

\begin{optionalexercise}\label{b3:ex:24-global-sections-not-right-exact}
设 \(k\) 为域，\(d\in\mathbb Z\)。在 \(\mathbb P_k^1\) 上考察
\[
0\longrightarrow\mathcal O(d-2)
\xrightarrow{(-x_1,x_0)}
\mathcal O(d-1)^2
\xrightarrow{(x_0,x_1)}
\mathcal O(d)\longrightarrow0.
\]
证明它正合，判定右端的全局截面映射何时满射，并求其余核。解释相同的局部满射为何会随扭曲产生不同的全局行为。
\end{optionalexercise}
\begin{solution}
\Cref{b3:prop:projective-line-section-obstruction}给出的正合列张量可逆层 \(\mathcal O(d)\) 后仍正合，得到所给序列。若 \(d<0\)，右端没有非零全局截面，映射自动满射。若 \(d\ge1\)，任意次数 \(d\) 的齐次单项式都可写成 \(x_0\) 或 \(x_1\) 乘次数 \(d-1\) 的单项式，故全局截面映射也满射。只有 \(d=0\) 时，源的全局截面为零、目标为 \(k\)，余核为 \(k\)；其余 \(d\) 的余核均为零。对应的连接同态把 \(1\) 送到 \(H^1(\mathcal O(-2))\) 中的非零 Laurent 类，记录这个唯一未能消去的全局提升障碍。局部正合性没有改变，扭曲后的全局截面却受到各层次数的限制。
\end{solution}

\begin{optionalexercise}\label{b3:ex:24-punctured-plane}
设 \(k\) 为域，令 \(R=k[x,y,z]\)、\(U=\operatorname{Spec}R\setminus\{(x,y,z)\}\)。计算全部 \(H^q(U,\mathcal O_U)\)，并证明 \(H^1(U,\mathcal O_U)=0\) 仍不足以推出仿射性。
\end{optionalexercise}
\begin{solution}
用 \(D(x),D(y),D(z)\) 覆盖，所有非空有限交都为 Noether 仿射概形，所以 Čech 比较适用。三个局部化环在 \(R_{xyz}\) 中的交为 \(R\)，故 \(H^0(U,\mathcal O_U)=R\)。支集 Čech 复形是三个一元复形
\[
0\longrightarrow k[x]\longrightarrow k[x,x^{-1}]\longrightarrow0
\]
及相应 \(y,z\) 复形在 \(k\) 上的张量积。每个一元复形按 Laurent 单项式分裂为恒等映射的可缩复形，以及只在次数一出现的负次幂空间。张量展开后，含可缩因子的项可逐项消去，惟有三个负次幂空间的张量积留在次数三。因此支集上同调只在三次非零；开覆盖复形去掉最左边的 \(R\)，在正次数恰少一层，得到
\[
H^1(U,\mathcal O_U)=0,\qquad
H^2(U,\mathcal O_U)=
\bigoplus_{a,b,c\ge1}k x^{-a}y^{-b}z^{-c},
\]
且 \(H^q=0\) 对 \(q\ge3\) 成立。\(U\) 由三张 Noether 仿射图覆盖，故为 Noether 概形；若它仿射，拟凝聚层的高次上同调消失将与非零的 \(H^2\) 矛盾。仿射性要求全部正次消失，单独一次消失无法保证这一点。
\end{solution}

\begin{optionalexercise}\label{b3:ex:24-maps-projective-to-affine}
设 \(k\) 为域，\(r\ge0\) 为整数。证明每个 \(k\) 概形态射 \(\mathbb P_k^1\rightarrow\mathbb A_k^r\) 都是常值态射。
\end{optionalexercise}
\begin{solution}
\Cref{b3:prop:scheme-hom-affine-target}给出任意源到仿射目标的态射对应。这里全局函数环为 \(\Gamma(\mathbb P_k^1,\mathcal O)=k\)，保持 \(k\) 结构的态射由 \(k[x_1,\ldots,x_r]\rightarrow k\) 给出，即选定 \(r\) 个标量。其像因而是这个 \(k\) 值点，态射经 \(\operatorname{Spec}k\) 分解。
\end{solution}

\begin{optionalexercise}\label{b3:ex:24-field-valued-projective-points}
设 \(K/k\) 为域扩张，\(r\ge0\) 为整数。证明 \(\mathbb P_k^r\) 的 \(K\) 值点与 \(K^{r+1}\setminus\{0\}\) 模去 \(K^\times\) 的同比例关系对应。
\end{optionalexercise}
\begin{solution}
一个 \(K\) 值点是 \(k\) 态射 \(\operatorname{Spec}K\rightarrow\mathbb P_k^r\)。源只有一点，它落在某张 \(D_+(x_i)\) 中，在这张图上给出 \(r\) 个坐标比 \(a_j/a_i\in K\)，可取 \(a_i=1\)。换到另一张非零坐标图时，全部坐标同时除以该坐标，恰好按同比例关系变化。反过来，非零向量可选一项为分母，由坐标比给出图内的态射；不同选择在交集一致，因而定义同一个 \(K\) 值点。这里对域成立的向量描述不能未经修改就用于一般基环；一般情形还要记录可逆模及其生成截面。
\end{solution}

